% 日本語 LuaLaTeX。2026-11-22 / v3.7
% エンジン、出力先、SyncTeX は repo ルートの .latexmkrc に従う。
\RequirePackage{fix-cm}
\documentclass[a4paper,11pt]{ltjsarticle}
\usepackage{amsmath,amssymb,amsthm,mathtools,mathrsfs}
\DeclareFontFamily{U}{rsfs}{\skewchar\font127}
\DeclareFontShape{U}{rsfs}{m}{n}{<-6> rsfs5 <6-8> rsfs7 <8-> rsfs10}{}
\usepackage[margin=25mm]{geometry}
\usepackage{booktabs,tabularx,array,enumitem}
\usepackage[unicode,pdfusetitle,colorlinks=true,
  linkcolor=blue,citecolor=blue,urlcolor=blue]{hyperref}
\numberwithin{equation}{section}
\theoremstyle{plain}
\newtheorem{theorem}{定理}[section]
\newtheorem{proposition}[theorem]{命題}
\newtheorem{lemma}[theorem]{補題}
\newtheorem{corollary}[theorem]{系}
\theoremstyle{definition}
\newtheorem{definition}[theorem]{定義}
\newtheorem{assumption}[theorem]{仮定}
\newtheorem{example}[theorem]{例}
\newtheorem{problem}[theorem]{問題}
\theoremstyle{remark}
\newtheorem{remark}[theorem]{注意}
\renewcommand{\proofname}{証明}
\newcommand{\kk}{k}
\newcommand{\OO}{\mathcal O}
\newcommand{\PP}{\mathbb P}
\newcommand{\Gm}{\mathbb G_m}
\newcommand{\HH}{\mathcal H_2}
\newcommand{\HHstack}{\mathfrak H_2}
\DeclareMathOperator{\Hom}{Hom}
\DeclareMathOperator{\RHom}{RHom}
\DeclareMathOperator{\End}{End}
\DeclareMathOperator{\Ext}{Ext}
\DeclareMathOperator{\Aut}{Aut}
\DeclareMathOperator{\GL}{GL}
\DeclareMathOperator{\PGL}{PGL}
\DeclareMathOperator{\SO}{SO}
\DeclareMathOperator{\SL}{SL}
\DeclareMathOperator{\Spin}{Spin}
\DeclareMathOperator{\Gr}{Gr}
\DeclareMathOperator{\Spec}{Spec}
\DeclareMathOperator{\rank}{rank}
\DeclareMathOperator{\Pic}{Pic}
\DeclareMathOperator{\Jac}{Jac}
\DeclareMathOperator{\coh}{coh}
\DeclareMathOperator{\qgr}{qgr}
\DeclareMathOperator{\Qgr}{Qgr}
\DeclareMathOperator{\Perf}{Perf}
\DeclareMathOperator{\Sym}{Sym}
\DeclareMathOperator{\coker}{coker}
\DeclareMathOperator{\im}{im}
\DeclareMathOperator{\Supp}{Supp}
\title{非可換二次曲面のavatarとAS正則$\mathbb Z$-代数：\\
  八・十・十二本矢の分類と内在的忘却}
\author{}
\date{2026年11月22日\quad v3.7（全noncubic pointed認識と相対有限End）}
\begin{document}
\maketitle
\begin{abstract}
本稿の第一の主結果は、指定された最小単純射影分解型を持つ
square/rotated square型AS正則$\mathbb Z$-代数と
適切な反標準avatarの\emph{同型類の全単射}である。
AS条件からの階数強制と点Hilbert曲面・Picard線束の抽出、
任意のavatarからのacyclic spherical helixによるAS--Gorenstein逆構成、
全整数Serre completionによる互逆性、および左右Noether性を証明する。
さらに独立に最小分解を課した十本矢型を定義し、
Ginzburg次数零集中、Keller--Yang mutation、
全周期の評価・余評価により、
square型全体とcubic型・最小十本矢型の非交和の全単射を証明する。
十本矢型のavatar条件は$G\simeq FP$である。
十二本矢はsquareの有限foundationの八矢と
四つのrelation-dual戻り矢からなるQP表示であり、
任意の十二本矢potentialのAS分類を意味しない。
AS分解のEuler漸化式を四周期有理母関数として解き、
全$m$で$\dim A_{i+4m,i}=(2m+1)^2$を証明する。
非被約反標準因子でもpencil trickにより
最小十本矢壁が$P\simeq GF^{-1}$に正確に一致する。
指定mutationについてKeller--Yangの局所条件を確認し、
mutation先の左右Noether性と
$\Qgr,\qgr$のAbel圏同値を共通ample cornerから導く。
square/rotated側ではcoherentな有界導来Morita同値を示すが、
これをO-marked構造対象を保つ同値と同一視しない。
二つの最小十本矢壁の交点は非零二捩れ$P^2=\OO_D,\ G=FP$であり、
その反標準曲線はsmooth、既約nodal、$I_2$の三型だけとなる。
さらに曲線と第一次数二pencilを固定すると、
この二重壁は二つのCartier壁の
\emph{scheme-theoreticに横断的な}交点であり、
各幾何fiber上の長さはsmooth/nodal/$I_2$で$3/1/1$となる。
滑らかな楕円族では二重壁は相対的に有限étale三次である。
これはavatarの標識付きパラメータ空間についての結論で、
内在的O-markedモジュライへの相対同型ではない。
さらにこの交点が、任意のcubic AS代数と
\emph{導来同値にならない}正確なlocusであることを証明する。
滑らかな楕円曲線では任意のcubic候補foundationを
数値Serre作用・determinant・pencil乗法のkernelで排除し、
退化二型へは二捩れを保つ平滑化と
exceptional perfect複体の形式持上げで拡張する。
この導来非同値は、内在的noncubic O-marked heartとの
pointed同値の構成とは異なる問題である。
underlying Z-代数のgroupoidと適切な相対perfect AS族にも
構成を拡張するが、未修正のstable QP groupoidとは区別する。

非可換二次曲面の四頂点表示について、矢数の変化をモジュライの族として記述する。
有限 foundation の生成矢、関係の双対を加えた箙、非最小な potential 表示を区別する。
固定した Hom 次元を持つ有限有向代数の結合律スキームを構成し、
二つの隣接乗法の階数から、八本、二種類の十本、十二本への層別を得る。
少ない矢数の開部分は、正確に十本または十二本を要求する層ではなく、
非最小表示を含めた共通モジュライの開部分である。
重み四 potential の二つの二次係数による消去は任意の基底上で明示できるが、
通常の right-equivalence 商には消去される対の自己同型が残る。
従って Jacobian 代数の一致を、そのまま表示スタックの同値とすることはできない。

本稿の中心目標は、構造対象を指定したO-marked非可換二次曲面、
線形代数的potential、反標準avatarの三者に、体上同型類の明示全単射
またはそれを誘導する相対モジュライ射を与えることである。
滑らかな楕円cubic型では、VdB四つ組とsquare avatarの双方向公式と
O-markingを保つruling反射による同型類商を証明する。
potentialからavatarへの一般表示次数は4、O-markedまでの合成は8となる。
さらにnoncubicの\textbf{全幾何点}について、二つの被約二根対の積
$q=f_1f_2$による有限三重quartic-resolventと
算術種数一の二重被覆$C_q:y^2=q$を構成し、
四根相異・一根共有・二根対一致に対応する
smooth/nodal/I2の三型を分類する。
各型で非零二捩れ$P\in\Pic^0(C_q)[2]$を使った
$(2,2)$反標準square avatarを明示し、
Bのnoncubic O-marked全幾何同型類と全単射を得る。
さらに平方型を相対幾何変換により、$\Sigma_2$ の負断面 $S$ を避ける
$D\in|-K_{\Sigma_2}|$ と $0\ne P\in\Pic^0(D)[2]$ へ移すことで、
全noncubic O-marked同型類を $\Sigma_2$ の反標準 avatar でも分類する。
$\chi=1$ の二つの特殊同型類は $D$ の既約 nodal 型と $I_2$ 型で区別される。
同時に中央二重平方環の可換化の点曲線$E_f$から
$E_f\to C_q$とrank-one torsion-free反不変module$M$への
相対射を構成する。
今回さらに、四根相異なる場合の$E_f\to C_q$は
\emph{次数二のétale isogeny}であり、
その二曲線が一般に同型でないことを証明する。
$E_f$上にも二つのconic商に由来する次数二pencil
$F_d,F_c$が存在し、
$0\ne F_cF_d^{-1}\in\Pic^0(E_f)[2]$となる。
$E_f$側と$C_q$側の二捩れ付き楕円曲線は
互いに\emph{双対二重isogeny}で交換される。
noncubicの二つの非自由MCM型の線形行列の
kernel射は$E_f$の二つのpencilを直接回収し、
その二つの隣接Hom積の像も
conic Veroneseの対称積像と一致する。
さらに、点曲線を二conic二重被覆のファイバー積
$E_f=Q_c\times_{\PP^1}Q_d$として全退化型で同定する。
smooth、共有根のnodal、根対一致の$I_2$の
全noncubic幾何点で四頂点End代数の全Homと積を
点曲線への制限から同定する。
mutation後の四項は一添字ずらし、
$(H_1,H_2,H_3,H_4)$に合わせることで
構造対象を保つpointed exact同値を証明する。
さらに、根を順序付きで分解した普遍係数空間では、
有限End代数の同型が任意の非被約基底変更と両立することを証明する。
これは相対$\qgr$圏の同値やroot-label忘却の降下とは異なる。
$C_q$・$\Sigma_2$の別avatarとのpointed比較も未解決である。
実例$f_1=x(x-1),f_2=(x-2)(x-3)$では
$j(E_f)=28756228/3$、$j(C_q)=35152/9$で、
点曲線とspectral曲線の単純な同一視は誤りである。
smoothでは$M=P$だが、
共有根で$M$は表面$A_1$型の非自由reflexive moduleとなるため、
通常の可逆な$\Pic^0$線束として固定$C_q$族に延長できない。
さらに共有根の全空間を最小解消しても、反不変moduleの
可逆化は例外$\PP^1$上に奇数次数（自然な可逆化では次数1）を持つ。
元の基底での垂直因子twistではその奇偶は変わらず、
通常の多重次数零の相対Picardには最小解消だけで移れない。
一方、反不変moduleの積$\mu:M\otimes M\to\OO_C$の像は
$(f_1,f_2,y)$で、余核は二根対の共通零点の構造層である。
この式は非被約基底を含む任意の基底変更と可換し、
smooth/nodal/$I_2$で欠損長は$0,1,2$となる。
各幾何fiberでは$M\simeq M^\vee$だが、
nodal極限は$\nu_*\OO_{\PP^1}(-1)$、
$I_2$極限は二つの$\OO_{\PP^1}(-1)$の正規化からの直像となる。
$I_2$では$M$自体の自己同型群が$\Gm^2$となり、
通常のsimpleなfine compactified Picardには入らない。
二次積を保持した固定曲線上の慣性は$\mu_2^2$に縮小する。
さらにこれらを任意の基底で自然化し、
二形式表示の商スタック$[P/G]$から
二重被覆・torsion-free sheaf・自己双対的二次積の組
$(C,M,\mu)$のfppf相対groupoidへの実際の射を構成する。
ただし標的のArtin性や、この射の全忠実性・
O-marked heartとの同値はここでは示していない。
従って、体上全単射、torsion-free spectral相対射、
元の可逆Picard avatarへの相対射を明確に区別する。
古典Picard零切断の重複、wallの次元差、
元AS avatarの$\qgr$心の相対認識の不足も区別する。
さらに未公刊のsquare／rotated square分類草稿のうち本研究に不可欠な
反標準曲線のコホモロジー消滅、二つのpencilによる再埋め込み、
有限directed foundationの行列・scheme-theoreticな復元、
正負全周期にわたる四つのcanonical exact complexを
仮定と証明を含めて本稿内に組み直す。
さらに最小AS分解の必要条件、spherical helixからの全代数の
逆構成、有限foundationの相互逆性を統合して、
体上のAS分類の証明を完成させる。
曲線上の完全列だけから任意基底上の相対AS分類や
内在的O-marked heartとの同値が従うわけではない。
幾何側では、$\Sigma_0$ または $\Sigma_2$ 上の反標準因子と
非自明な $\Pic^0$ パラメータを用い、直積型 square avatar と
二種類の十本矢 chart の関係を記述する。
統一偏極と rotated square の Hilbert 偏極は異なり、
$\Sigma_2$ の出現だけでは最小十二本矢型を意味しない。
最後に、構造対象のみを指定した強い内在的モジュライ $\HHstack$ への忘却を扱う。
相対 heart、反 Serre の豊富な exact 平方根、同じ根の全軌道環の AS4 被覆を
確認した部分では忘却射が定義される。
この適合部分の開性や、統一 avatar と全相対 AS 代数とのスタック同値は
別の認識問題として残し、有限表示の代数性とは区別する。
Okawa--Ueda の quintuple と conifold potential の既存 GIT 商を比較対象に加え、
quintuple 表示の四次元スカラー慣性が pointed な忘却で消えること、
ただし一般点の有限慣性$\mu_2^2$は剛体化後も残り、
点曲線のJacobian二捩れ翻訳群に一致することを証明する。
有効potential/quintuple表示スタックの一般次数四射は
慣性・接写像が同型で、適切な稠密開では表現可能有限étaleとなる。
さらに幾何的点曲線の完全忠実な点対象から、
一般の有限慣性が内在pointed忘却でも忠実に残ることを証明する。
論文Bの普遍cubic smooth atlasを合わせると、剛体化quintupleから
$\HHstack$への射はsmooth相対次元零で、一般開では
表現可能étale、Kodaira--Spencer同型、$HH^1=0,\dim HH^2=3$。
さらに二つの内在的ruling foundationの比較から、
一般の幾何fiberは二点となり有限étale次数二を得る。
有効potential表示から$\HHstack$への合成射は一般次数八である。
一般有限étale二重射のcanonical deck対合を取ると、
内在$\HHstack$の稠密開は剛体化quintupleの一般開の
$S_2$商スタックに同値となり、対合は二つのnormalized
ruling foundationの交換と幾何点上で一致する。
一般openには$\delta,\eta-\delta\notin E[2]$を課す。
一方、$j(E)\ne0,1728$の滑らかな許容四つ組で
交換自己同値が生じるのは$\delta\in E[2]$または
$\eta-\delta\in E[2]$のとき、かつそのときに限る。
正則半単純特別点の全慣性は位数8となる。
同じquintuple同型類が一つでも相対2-fiberには
pointed比較同型の剰余類から二対象が残り、
忘却はstackとして表現可能étale・非分岐である。
さらに片側の二捩れ条件のみ成立する正則半単純楕円点では、
translation型の全慣性群は$C_4\times C_2$、
反転型では$C_2^3$である。
さらに$e_1,e_3$を二乗固有値の基本対称式、
$c=\prod_i x_i$とすると、二捩れ特別座は
$e_3\pm e_1c=0$で与えられる。
OU不変量では$D_\pm=f_2^3-3f_2f_4+2f_6\pm6f_2g_4$となり、
正則半単純開の粗忘却の分岐座は$D_+D_-=0$である。
二つの片側特別座では慣性の三次元接表現が
$(1,1,-1)$となり、粗射は形式的に$Z=z^2$となる。
両条件が同時に成立する$e_1=e_3=0$の交点では
源慣性が$C_2^3$、標的慣性が$D_8\times C_2$へ増大し、
粗射の局所型は$k[[s,a,c]]\to k[[s,a,b,c]]/(b^2-ac)$となる。
これらの粗空間分岐にもかかわらず、忘却のスタック射自体はétaleである。
さらに両方のrulingがcubicである開では$j(E)=0,1728$でも
追加の複素乗法自己同型が三束保存・交換に関与できず、
level-twoのLegendre変形座標によって同じ法表現・粗分岐が成り立つ。
他方、$\eta=2\delta$、$\delta\notin E[2]$の一cubic境界では
もう一方のfoundationがnoncubicになるため、忘却はétaleでもfiberは一点。
特に$j=0$の非零3捩れCM固定点では全慣性$A_4$が現れ、
接の$C_3$重み$(2,1,1)$による孤立巡回商型の粗形式環となる。
他方、論文BのArtin foundation持上げと相対mutationにより、
任意のpointed Artin変形では二つのnormal foundationがそれぞれ
同型類一意に持ち上がり、相対kernel/cokernel mutationが交換する。
さらにBの有限型smooth atlasの相対foundationを用い、
二つのrulingを保つか交換するかを相対perfect Homの階数で開閉判定し、
全内在$\HHstack$上に表現可能な有限étale次数二の
normal ruling選択被覆$\mathfrak H_2^{\mathrm{rul}}$を構成する。
その選択foundationの相対End代数は有限foundation表示スタック
$\mathfrak F_0$への開埋込を与え、実現可能な$S_2$-不変開
$\mathfrak F_{\HH}$について
$\HHstack\simeq[\mathfrak F_{\HH}/S_2]$が全相対群oidの同型となる。
従って一cubic境界でも内在的な第二のnoncubic foundationは
有限型・非被約基底にわたって二択の一方として存在する。
この実現可能開の逆像は$\mathbb A^{64}$の不変開なので、
Noether性により有限個の多項式の非同時消滅条件で記述できる。
これは全実現可能開に対する非構成的存在証明であり、その具体式・次数上界は未計算。
一方Fが可逆なcubic幾何部分では、四つの隣接flatteningの
行列式（次数4）とadjugate列の二次行列式（次数6）で
計20式の具体的開判定を得る。
また可換$\PP^1\times\PP^1$のfactor交換から、
このruling被覆に自己同値と両立する大域sectionは存在しない。
忘却射の相対ファイバーには表示の軌道だけでなく自己同型群の商が現れる。
既存の次数四 GIT 射を Serre 平方根の忘却と同一視しない。
OU射の座標置換群はS4、有効符号群はF2²、Galois閉包の群も別のS4であり、滑らかな楕円曲線上の
Picard二等分はE[2]-torsorの次数4となる。
節点・尖点へ退化するとPicard根の幾何fiberは2点・1点へ変わる。
OU対角四重射自体は、二乗座標を分岐点とするspectral楕円曲線上で、
明示したsectionの二等分Kummer torsorに幾何的に同型である。
対角potentialのgeometric quintuple点曲線は二quadric完全交叉であり、
その鉛筆の判別式が同じspectral曲線となるため、Jacobianとの同型と
GIT不変量による$j$の公式も得る。Hodge対合でrulingの符号を決定し、
滑らかな対角稠密開ではOU四重射とquintupleのPicard根torsorの
同型を偶数符号・置換・スカラーに同変に構成する。
正則半単純な非対角potential族には相対Hodge-Jacobian同型と
特性係数による相対半点を任意底変更と可換に構成する。
potentialを選ばない一般quintupleから内在$\HHstack$への
有限foundation表示内の実現可能開を具体的な有限minor/相対AS条件として
明示すること（有限多項式の存在までは証明）、avatar/QP表示との全相対認識、
および局所粗モデルの大域化は未解決である。
さらに非線型geometric quintupleの全相対族からは
根を標識として選ぶことなくpointedな内在的H2への射が構成でき、
選んだ四つのPicard根は同じ相対tails対象に忘却される。
剛体化後の忘却の表現可能性と実際のファイバーは引き続き未決定である。
\end{abstract}
\tableofcontents

\section{Introduction}
\label{sec:qa-introduction}

本論文の第一の主定理は、
\emph{avatarと、最小射影分解型を指定したAS正則$\mathbb Z$-代数}
の同型類の全単射である。
その証明は第~\ref{sec:qa-as-avatar-classification}節に、
十本矢の独立ASクラスとQP表示の正確な関係、
rank-three壁、mutation後のNoether性・商圏同値は
第~\ref{sec:qa-as10-mutation-classification}節に置く。
同節では二重壁のscheme-theoretic横断性と
非零二捩れの相対有限étale三次も記述する。
二つのmutation壁の交点と、cubic型への到達不能性の完全判定は
第~\ref{sec:qa-as-cubic-obstruction}節で証明する。
最小分解からの有理Hilbert級数の導出と
AS--Gorenstein性の証明を別の論理段階として扱う。
また noncubic 二形式に自然な点曲線$E_f$と
spectral 曲線$C_q$が2-isogenousだが一般に非同型であることを
第~\ref{sec:qa-point-spectral-isogeny}節で説明し、
全noncubic幾何点では$E_f$の二conicファイバー積と
二pencilを用いて、mutation後のAS heartとの
pointed認識を証明する。
第~\ref{sec:qa-point-spectral-isogeny}節では
MCM行列kernelから二pencilを回収し、
四頂点Endの全Homと積を点曲線への制限で同定する。
標準頂点$H_0$でなくshift後の$H_1$に構造対象を対応させる。
順序付き分解の普遍族では、有限graded End代数の比較を
任意の非被約基底変更へ延長する。
ただし相対$\qgr$圏・Serre根の一括降下と
無標識stackの群oid認識は別問題である。
ここでは有限directed foundation、relation-dual QP、
underlying AS代数、O-marked内在的heartの四水準を混同しない。

本論文の第二の問題は、\textbf{O-marked quadric、potential、avatar}を
単に隣接する記述として並べるのではなく、
具体的な同型類の全単射と、それを誘導するモジュライ間の射で結ぶことである。
特に、potentialの頂点と四項foundation、avatarの二ruling、
内在的な構造対象$\OO$をそれぞれ区別して忘却を行う。
滑らかな楕円cubicの三線束辞書と$r_1$商、
共有根・根対一致を含む全noncubic境界の二捩れavatarと二根対の
幾何的全単射、およびrank-one torsion-free spectral中間射と
可逆Picard延長障害を本稿の具体的な比較定理とする。
幾何的基礎は第~\ref{sec:qa-self-contained-foundations}節と
第~\ref{sec:qa-self-contained-helix}節で自己完結に証明し、
同型類・相対spectral射の証明は
第~\ref{sec:qa-avatar-charts}節に置く。
共有根の最小解消と奇偶障害、
さらに二次積のscheme-theoretic欠損と三型の
正規化極限・固定曲線上の慣性・
二次積付きtorsion-free相対groupoidへの射は
第~\ref{sec:qa-nodal-resolution-parity}節で証明する。
$\Sigma_2$の最深境界・非被約因子・古典零切断・全相対AS認識は
結論に含めず、必要な追加構成を明記する。


三次元 AS 正則 $\mathbb Z$-代数の四周期表示では、同じ幾何的構成に対して
八本、十本、十二本の矢を持つ箙が現れる。
これらを別々の分類集合として並べるだけでは、
生成元と関係が同時に消去される変形や、異なる foundation の比較を捉えられない。
本稿は、まず共通の有限代数のモジュライを構成し、その上の二つの開条件と
二つの階数低下条件として矢数を整理する。
さらにgeometric cubic部分の隣接四重tensorの非退化性は
四次4式・六次16式の明示的多項式判定へ書き換える。
次に反標準 avatar の幾何的 chart を比較し、
最後に構造対象付き Abel 圏への忘却に必要な条件を分離し、
Okawa--Ueda の quintuple/potential GIT と比較する。
特に滑らかな楕円型の稠密開では二つのnormalized rulingが
有限étale二重忘却の二sheetになり、そのstack商として
内在的モジュライを記述できる。
さらに二捩れ特別座では、内在的慣性の法表現を求め、
OU不変量による粗分岐式と交点の横断$A_1$局所型まで計算する。
これらを既知のpotential/quintupleの次数四射や
Picard根の四選択とは区別して扱う。
複素乗法で$j=0,1728$となっても二cubic開では追加の交換はなく、
一cubic境界ではfiberのsheetを一つ失うことも区別する。

\subsection{三つの表示と内在的モジュライの比較}

本稿で扱う問題には、次の三つの水準がある。
とくに既存GIT商を内在的モジュライと結び付ける相対比較を本稿Cの
中心研究課題とする（第~\ref{sec:qa-git-comparison}節）。
\begin{enumerate}
\item \emph{有限表示のモジュライ。}
  四頂点間の Hom 次元と結合律を固定した代数スタックは有限型 Artin である。
  隣接する二つの乗法が両方とも同型である部分、一方が同型である二つの部分、
  双方に一次元の欠損を許す部分を比較する。
  これが八本・十本・十二本の開部分の正確な有限代数版である。
\item \emph{AS 代数と avatar。}
  四周期の全 AS 条件と幾何的 avatar の対応には、有限 foundation の分類に加えて
  全単純分解、noetherian 性、相対平坦性、降下が必要である。
  体上の既存原稿の分類結果と、今回の有限代数・幾何的 chart の証明を区別する。
\item \emph{内在的な対象。}
  構造対象 $\OO$ だけを指定した $\HHstack$ は、選んだ箙、potential、
  foundation、平方根、ambient を忘れる。
  しかしそれらの存在条件を失うわけではない。
  特に AS3 正則性や $\Sigma_2$ という有限表現モジュライだけから、
  $\HHstack$ の全相対条件は導けない。
\end{enumerate}

有限代数版では、二種類の十本矢 chart の交わりが八本矢 chart である。
二つの十本矢 chart の和集合の補集合が最小十二本矢の層となる。
これに対し、各 chart の内部で正確に十本が最小生成矢として残る部分は
八本矢部分の境界である。
このため、本稿では上限を許す chart と最小本数が一定の層を同じ記号で表さない。

直積型 collection
\[
 (\OO,\OO(1,0),\OO(0,1),\OO(1,1))
\]
そのものの有限箙は八本、四本の最小関係の双対を加えた relation QP は十二本である。
固定した一頂点 mutation により、八本と一方の十本を含む chart が得られる。
従って直積型と十本矢 chart の対応は mutation を介した対応であり、
元の箙の矢数を十本と数えるという意味ではない。

\subsection{先行研究と依存関係}

有限箙の関係式を moduli stack として扱う一般的枠組みは
Abdelgadir--Okawa--Ueda \cite[Definitions 3.5, 4.1; Propositions 3.6, 4.2]{AOU2014}
による。本稿は四頂点の Hom 次元を固定して結合律を具体的な行列方程式で書き、
二つの乗法の階数が最小箙に与える影響を計算する。
これは一般論の新しい証明を主張するものではない。
Derksen--Weyman--Zelevinsky \cite[Theorem 4.6; Section 5]{DWZ2008}
は完備道代数の potential の splitting と mutation を構成した。
ここで必要な消去は内部次数四のため有限回の三角変換で済み、
通常の道代数にも任意基底にも同じ式を適用できる。
ただし pair の消去と quotient stack の慣性の消去は別の操作である。

AS 正則 quadric の cubic 型については Van den Bergh \cite{VdB2011}、
その幾何的データと potential のモジュライについては
Okawa--Ueda \cite[Theorem 1.1]{OU2014} が基礎となる。
同定理の四次有限射は potential の GIT 商から quintuple の GIT 商への射であり、
内在的 $\HHstack$ への根・構造対象忘却の次数を述べるものではない。
表示側の四次元スカラー慣性の核を命題~\ref{prop:qa-ou-scalar-gerbe}、
剛体化後にも残る一般有限慣性を命題
~\ref{prop:qa-ou-generic-quintuple-finite-inertia}で計算する。
OUの一般安定化群という表現は単位連結成分と全慣性を区別して読む。
また、potential$\to$quintupleの既知GIT次数四射を
定理~\ref{thm:qa-ou-rigid-potential-quintuple-generic-etale}で
一般の表現可能有限étaleスタック射へ強める。
さらに命題~\ref{prop:qa-ou-generic-inertia-faithful-to-h2}では
点曲線の点対象からpointed忘却の一般慣性単射を証明する。
論文Bの普遍smoothnessと合わせた一般étale性・KS同型は
定理~\ref{thm:qa-quintuple-h2-generic-etale}、
一般次数二とpotential合成の次数八は
系~\ref{cor:qa-quintuple-h2-generic-degree-two}と
系~\ref{cor:qa-potential-h2-generic-degree-eight}に示す。
さらに稠密開でのスタック的な商同型は
定理~\ref{thm:qa-h2-generic-ruling-quotient}で与える。
Artin形式変形上の二normal foundationの一意持上げは
命題~\ref{prop:qa-ou-formal-two-ruling-lifts}に示す。
さらに全H2のnormal ruling選択の有限étale二重被覆は
定理~\ref{thm:qa-global-two-ruling-etale-cover}、
finite foundationの実現可能開への開埋込と
全相対$S_2$商同型は
定理~\ref{thm:qa-intrinsic-h2-open-in-foundation-stack}に与える。
有限多項式判定の存在は系~\ref{cor:qa-realizable-foundation-finite-open-criterion}、
被覆の可換gerbe上での非自明性は系
~\ref{cor:qa-ruling-cover-nontrivial-commutative-gerbe}で確認する。
二捩れ特別座のspectral方程式と粗局所分岐は
命題~\ref{prop:qa-ou-spectral-torsion-divisors}、
定理~\ref{thm:qa-ou-special-reflection-coarse-branch}、
系~\ref{cor:qa-ou-coarse-branch-invariant-equations}に、
交点の慣性・横断$A_1$モデルは
定理~\ref{thm:qa-ou-double-torsion-intersection-local-model}に示す。
片側二捩れ特別座の全慣性群型は
命題~\ref{prop:qa-elliptic-special-inertia-group-types}に示す。
特殊ruling交換の正確な二捩れ条件は
命題~\ref{prop:qa-elliptic-ruling-swap-two-torsion}に証明する。
特別点の2-fiberは命題
~\ref{prop:qa-special-ruling-two-fiber-no-stack-ramification}に示す。
ここで用いる square・rotated square の反標準 avatar と十本矢 mutation は
大川--植田の日本語草稿 v7 \cite{AvatarVSeven2026} に、
十二本矢型の追加と統一 avatar は研究ノート \cite{Q12Avatar2026} に記載されている。
幾何的な相対 twist は \cite{Q12Weyl2026}、
これらの関係の整理は \cite{Synthesis2026} を参照する。
これらは未公刊の研究資料であり、そこに記載された全 AS 分類の証明を
今回独立に検証済みの結果として引用しない。
本稿では必要な有限線形代数・消去・chart の証明を本文に置き、
追加の再構成入力を要する主張を明記する。

内在的な $\HH$ とその相対化 $\HHstack$ は、同じ研究の主論文
\cite{QuadricIntrinsic2026} の対象である。
分類主論文は構造対象を保つ exact 同値による分類を扱い、
新しい相対モジュライ原稿 \texttt{nc-quadric-moduli.tex} は
相対認識と境界の幾何の証明正本である。
本稿CにOkawa--Uedaの既存GITとの具体的な相対比較を集約する。
本稿は表示の層別と avatar の幾何を独立に説明する。
主論文中の全分類や相対再構成の証明は本稿に重複させない。
本稿の結果も、全 $\HHstack$ の代数性、全 avatar の相対 AS 認識、
または公刊に向けた包括的な新規性確認の完了を意味しない。

\subsection{規約}

$\kk$ は標数零の代数閉体とする。
族は $\kk$ 上局所有限型スキームの fppf site 上で考え、非被約基底を含む。
矢の積は合成順で、$ab$ は先に $b$、次に $a$ を通る道である。
有限 foundation は四つの順序付き頂点を保持し、
右加群規約に対応する導来同値と標準 heart の exact 同値を区別する。
基底上の係数を消去に使う条件は可逆性であり、単に零でないことではない。

% !TEX root = ../quiver-avatar-moduli.tex
% 2026-11-21 / v1.0
\section{矢の本数、安定化表示、および開部分スタック}
\label{sec:qa-presentations}

本節では、有限 foundation の最小矢、relation-dual の戻り矢を加えた
箙、重み付き箙の非最小表示を区別する。
この区別のもとで、八本矢のモジュライが十本矢のモジュライの開部分となる
という主張を、任意の底上で正確に定式化する。
基礎体を $\kk$ とし、積 $ba$ は $a$ の後に $b$ を合成する規約とする。
本節の代数的な消去には標数の制限を要しない。

\subsection{直積型 collection の最小表示}

\begin{proposition}[直積型 collection の矢と関係]
\label{prop:qa-product-count}
$X=\PP^1\times\PP^1$ とし、
\[
 (E_0,E_1,E_2,E_3)
   =(\OO,\OO(1,0),\OO(0,1),\OO(1,1))
\]
と置く。有限代数 $B=\End_X(\bigoplus_{i=0}^3E_i)$ の最小箙は
\[
 a_i:0\longrightarrow1,\quad b_j:0\longrightarrow2,\quad
 u_j:1\longrightarrow3,\quad v_i:2\longrightarrow3
 \qquad(i,j=0,1)
\]
の八本の矢を持ち、最小関係は
\begin{equation}
\label{eq:qa-square-relations}
 r_{ij}=u_ja_i-v_ib_j\qquad(i,j=0,1)
\end{equation}
の四本である。従って relation-dual の戻り矢
$c_{ij}:3\to0$ を加えた relation QP は十二本矢であり、potential は
\begin{equation}
\label{eq:qa-square-potential}
 W_{\mathrm{sq}}=\sum_{i,j=0}^1c_{ij}(u_ja_i-v_ib_j)
\end{equation}
である。これは道長に関して reduced である。
\end{proposition}

\begin{proof}
$V=H^0(X,\OO(1,0))$、$U=H^0(X,\OO(0,1))$ はともに二次元である。
非恒等 Hom のうち、$\Hom(E_0,E_3)$ 以外の非零成分は
$V,U,U,V$ の四つである。
$\Hom(E_1,E_2)=H^0(X,\OO(-1,1))=0$ であり、
$\Hom(E_0,E_3)=V\otimes U$ への二つの道の合成はともに同型である。
従って二段階の道空間 $(U\otimes V)\oplus(V\otimes U)$ から
$V\otimes U$ への核は四次元で、\eqref{eq:qa-square-relations} が基底となる。
それより長い道は存在せず、これらが関係の全てである。
関係に一次部分はないので八本の矢は最小であり、
関係を他の矢で左右から乗じて新たな関係を得る余地もない。
四本の双対矢を加えると十二本となる。
\eqref{eq:qa-square-potential} は道長三なので、その二次部分は零である。
\end{proof}

\begin{remark}[十二本という数だけでは箙を指定できない]
\label{rem:qa-two-twelve-quivers}
命題~\ref{prop:qa-product-count} の箙を $Q_{\mathrm{sq}}$ と書く。
後述する $Q_{12}$ は四辺に各二本の矢を持つ有向四角形と四本の対角矢からなる。
$Q_{\mathrm{sq}}$ は $3\to0$ に四本の矢を持ち、$1\to2$ の矢を持たない。
従って両者は同じ十二本矢の箙ではない。
また、直積型 collection 自身の relation QP を十本矢と数えることはできない。
十本矢との対応を主張する際には、mutation と collection の変更を明示する必要がある。
この数え方は \cite[Remark 38.2]{AvatarVSeven2026} と一致する。
\end{remark}

\subsection{二つの可縮 pair を持つ共通表示}

四頂点を $\mathbb Z/4\mathbb Z$ で標識し、$Q_8$ を
各 $i\to i+1$ に二本の矢を持つ有向四角形とする。
その矢の重みを一とする。
$Q_{10}^{\mathrm{odd}}$ はこれに
\[
 x:1\longrightarrow3,\qquad y:3\longrightarrow1
\]
を加えた箙で、追加矢の重みは二とする。
$Q_{10}^{\mathrm{even}}$ は代わりに
\[
 u:0\longrightarrow2,\qquad v:2\longrightarrow0
\]
を加えたもの、$Q_{12}$ は両方の pair を加えたものとする。
以下の potential は重み四の巡回同値類であり、potential の定数倍は
同一視しない。

\begin{lemma}[共通 potential の形]
\label{lem:qa-universal-potential}
$Q_{12}$ の重み四の potential は一意に
\begin{equation}
\label{eq:qa-universal-potential}
 W=H+yq+px+vq'+p'u+\lambda yx+\mu vu
\end{equation}
と書ける。ここで $H$ は $Q_8$ の道長四の potential、
$q:1\to3$、$p:3\to1$、$q':0\to2$、$p':2\to0$ は
$Q_8$ の道長二の線形結合であり、$\lambda,\mu$ は係数である。
従って基底を固定した potential 空間は
\[
 P_8\simeq\mathbb A^{16},\qquad
 P_{10}^{\mathrm{odd}}\simeq P_{10}^{\mathrm{even}}\simeq\mathbb A^{25},
 \qquad P_{12}\simeq\mathbb A^{34}
\]
である。これらの記述は任意の $\kk$-scheme 上で成立する。
\end{lemma}

\begin{proof}
重み四の cycle が重み二の矢を持たなければ、四辺を一周するので
$2^4=16$ 個の道が巡回同値類の基底を与える。
重み二の矢を一つ持つ場合は、残りの二本が $Q_8$ の道長二の道となり、
四本の対角矢それぞれについて四個の可能性がある。
重み二の矢を二つ持つ cycle は $yx$ と $vu$ の二つだけである。
この三種類は巡回回転で混ざらないので、直和分解と次元の式を得る。
同じ道の基底を任意の可換係数環上で用いれば相対的な主張も従う。
\end{proof}

\begin{proposition}[任意の底上の有限な消去]
\label{prop:qa-relative-elimination}
$T$ 上で \eqref{eq:qa-universal-potential} が与えられているとする。
$\lambda\in\Gamma(T,\OO_T)^\times$ なら、
\[
 x'=x+\lambda^{-1}q,\qquad y'=y+\lambda^{-1}p
\]
によって
\begin{equation}
\label{eq:qa-one-pair-elimination}
 W=\lambda y'x'+H-\lambda^{-1}pq+vq'+p'u+\mu vu
\end{equation}
となる。従って通常の重み付き道代数で定義した Jacobian 代数について
\[
 \Jac(Q_{12},W)
 \simeq
 \Jac(Q_{10}^{\mathrm{even}},H-\lambda^{-1}pq+vq'+p'u+\mu vu)
\]
が成り立つ。$\mu$ が可逆な場合も同様で、両方が可逆なら
\begin{equation}
\label{eq:qa-two-pair-elimination}
 \Jac(Q_{12},W)
 \simeq\Jac(Q_8,H-\lambda^{-1}pq-\mu^{-1}p'q').
\end{equation}
これらの構成は任意の基底変換と両立する。
\end{proposition}

\begin{proof}
$p,q$ は $x,y,u,v$ を含まないので、表示した変換は逆も多項式である
重み付き道代数の自己同型である。式は積を展開して確認できる。
変換後の cyclic derivative は $\lambda x'$ と $\lambda y'$ を含むため、
Jacobian 商では $x'=y'=0$ である。
残りの cyclic derivative は残余 potential のものに一致する。
二つの pair の変換は互いの矢を変更せず、順序を問わず実行できる。
用いた演算は係数の加減乗除と、表示した可逆元の逆元だけである。
従って同型は任意の基底変換後にも同じ式で成立する。
\end{proof}

\begin{remark}[一般の splitting theorem との関係]
一般の QP の splitting theorem は完備道代数上の
\cite[Theorem 4.6]{DWZ2008} による。
命題~\ref{prop:qa-relative-elimination} では混合項が追加矢を含まないため、
有限回の三角変換だけで分離できる。従って完備化も極限操作も不要である。
この相対計算は AS 正則性や各次数成分の平坦性を主張していない。
それらは別の仮定または定理を必要とする。
\end{remark}

幾何点上の最小矢数は次の四つの場合に分かれる。
表の $\lambda=0$、$\mu=0$ は、その幾何点の剰余体での等式である。
\begin{center}
\begin{tabular}{c|c|c}
係数 & 残る最小対角矢 & 最小矢数 \\
\hline
$\lambda\ne0,\ \mu\ne0$ & なし & $8$ \\
$\lambda=0,\ \mu\ne0$ & $x,y$ & $10$（odd）\\
$\lambda\ne0,\ \mu=0$ & $u,v$ & $10$（even）\\
$\lambda=0,\ \mu=0$ & $x,y,u,v$ & $12$
\end{tabular}
\end{center}
消去後の potential は道長三以上なので、Jacobian ideal に一次関係はない。
これが表の最小性の根拠である。
非被約な底上では、全幾何点で係数が零でも係数自身が零とは限らない。
例えば $\lambda=\epsilon$ on $\Spec\kk[\epsilon]/(\epsilon^2)$ は可逆でなく、
最後の pair を消去できない。

\subsection{通常の QP スタックに残る慣性}

ここでは頂点を固定する全ての重み付き道代数自己同型を射とする。
$G_8=(\GL_2)^4$ と置くと、重み一の矢の像は基底変更で尽くされる。
十本矢では、さらに各追加矢の非零スカラー倍と、
同じ始点・終点を持つ道長二の道の付加が許される。従って
\[
 G_{10}=(G_8\times\Gm^2)\ltimes\mathbb G_a^8,
 \qquad
 G_{12}=(G_8\times\Gm^4)\ltimes\mathbb G_a^{16}.
\]
群の作用は potential の変数変換であり、
$\mathcal P_r=[P_r/G_r]$ を通常の重み付き QP 商スタックとする。
この節では頂点スカラーの rigidification は行わない。

\begin{proposition}[可縮 pair の残余慣性]
\label{prop:qa-stabilized-inertia}
上の群と頂点標識を用いると、
\begin{align}
 [\{\lambda\ne0\}\subset P_{10}^{\mathrm{odd}}/G_{10}]
   &\simeq \mathcal P_8\times B\Gm,\label{eq:qa-qp-gerbe-one}\\
 [\{\lambda\mu\ne0\}\subset P_{12}/G_{12}]
   &\simeq \mathcal P_8\times B(\Gm^2).
   \label{eq:qa-qp-gerbe-two}
\end{align}
特に、消去によってこれらの通常の QP スタックを
$\mathcal P_8$ と同定することはできない。
\end{proposition}

\begin{proof}
十本矢の場合、命題~\ref{prop:qa-relative-elimination} により
$p=q=0$ とでき、さらに一方の矢のスケールを変えて $\lambda=1$ とできる。
これは係数環上の可逆元の逆を用いる代数的な正規化である。
正規化した切片は $H+yx$ の形の potential の空間 $P_8$ である。

この切片の二点 $H+yx$ と $\widetilde H+yx$ を結ぶ自己同型を調べる。
重み一部分の作用を $g\in G_8$、対角矢の作用を
$x\mapsto\alpha x+q_0$、$y\mapsto\beta y+p_0$ と書く。
変換後に $yq_0$ と $p_0x$ の混合項が消えるには、
$\alpha,\beta$ が可逆なので $p_0=q_0=0$ が必要である。
$yx$ の係数から $\alpha\beta=1$ を得る。
従って切片間の全ての射は
\[
 (g,t):\quad H\longmapsto gH,\quad
 x\longmapsto tx,\quad y\longmapsto t^{-1}y
\]
であり、その群は $G_8\times\Gm$ である。
後半の $\Gm$ は $P_8$ に自明に作用する。
各対象が fppf 局所にこの切片へ移り、切片間の射がちょうど上の群で
記述されるので、商スタックの同値~\eqref{eq:qa-qp-gerbe-one} を得る。
二つの pair について同じ正規化を独立に行えば
\eqref{eq:qa-qp-gerbe-two} を得る。
ここでは切片とその安定化群を明示したので、単に gerbe と述べるだけでなく
直積との同値まで得られている。
\end{proof}

\begin{remark}[表示を忘れる操作]
消去される pair の変換 $x\mapsto tx$、$y\mapsto t^{-1}y$ は
Jacobian 商では恒等である。
従って通常の QP groupoid から代数 groupoid への忘却は忠実でない。
同じ障害は \cite[Proposition 52.3]{AvatarVSeven2026} で指摘されている。
一般に関係式の表示の自己同型と商代数の自己同型を区別すべきことは
\cite[Proposition 3.3]{AOU2014} にも現れる。
有限有向箙の関係式スタックの有限型性
\cite[Definition 3.5, Proposition 3.6]{AOU2014} は有用であるが、
その定理だけから、ここで扱う無限次元 graded 代数の全次数平坦な族や
内在的モジュライの代数性は従わない。
\end{remark}

\subsection{代数の射を用いた正確な開埋込み}

\begin{definition}[代数の fppf スタック]
\label{def:qa-algebra-stack}
$\mathfrak M_{12}^{\mathrm{alg}}$ の $T$ 上の対象は、
$A_0=\OO_T^4$ と四つの標識付き冪等元を持つ非負次数付き代数
$A=\bigoplus_{d\ge0}A_d$ で、各 $A_d$ が有限階数局所自由、かつ
fppf 局所に $\Jac(Q_{12},W)$ の形の表示を持つものとする。
射は頂点と次数を保つ全ての代数同型とし、potential、矢の基底、
表示からの全射はデータに含めない。
$\mathfrak M_{10}^{\mathrm{odd,alg}}$、
$\mathfrak M_{10}^{\mathrm{even,alg}}$、$\mathfrak M_8^{\mathrm{alg}}$
もそれぞれの箙を用いて同様に定義する。
これらは fppf 降下を持つ代数のスタックの充満部分スタックである。
この定義に AS 正則性は含めない。
\end{definition}

各対象について、添字を $\mathbb Z/4\mathbb Z$ で読み、有限表示層
\begin{equation}
\label{eq:qa-indecomposables}
 C_i(A)=\operatorname{coker}\left(
 e_{i+2}A_1e_{i+1}\otimes_{\OO_T}e_{i+1}A_1e_i
 \longrightarrow e_{i+2}A_2e_i\right)
\end{equation}
を定める。これは重み二における新しい生成元の層であり、
定義は表示に依存しない。

\begin{theorem}[矢の少ない代数スタックの開性]
\label{thm:qa-open-algebra-stacks}
可縮 pair を追加することによって、充満な開埋込み
\begin{equation}
\label{eq:qa-open-diagram}
 \begin{aligned}
 \mathfrak M_8^{\mathrm{alg}}&\lhook\joinrel\longrightarrow
 \mathfrak M_{10}^{\mathrm{odd,alg}}
 \lhook\joinrel\longrightarrow\mathfrak M_{12}^{\mathrm{alg}},\\
 \mathfrak M_8^{\mathrm{alg}}&\lhook\joinrel\longrightarrow
 \mathfrak M_{10}^{\mathrm{even,alg}}
 \lhook\joinrel\longrightarrow\mathfrak M_{12}^{\mathrm{alg}}.
 \end{aligned}
\end{equation}
を得る。
両十本矢開部分の交わりは $\mathfrak M_8^{\mathrm{alg}}$ である。
正確には、それぞれの像は
\begin{align*}
 \mathfrak M_{10}^{\mathrm{odd,alg}}&=\{C_0=C_2=0\},\\
 \mathfrak M_{10}^{\mathrm{even,alg}}&=\{C_1=C_3=0\},\\
 \mathfrak M_8^{\mathrm{alg}}&=\{C_0=C_1=C_2=C_3=0\}
\end{align*}
である。開埋込みという主張は任意の scheme からの基底変換について成立し、
各スタックの代数性を仮定しない。
\end{theorem}

\begin{proof}
十本矢の表示に $vu$ を追加しても Jacobian 代数は変わらず、
八本矢からは $yx$ と $vu$ を追加できる。
従って同じ代数を対象として送る関手が存在する。
射は両辺とも全ての次数付き頂点固定代数同型なので、関手は充満忠実である。

局所表示~\eqref{eq:qa-universal-potential} で
\eqref{eq:qa-indecomposables} を計算する。
重み二の関係は四本の重み二の矢による cyclic derivative だけである。
道長二の積を割ると、これらの関係は
$\lambda x=\lambda y=\mu u=\mu v=0$ となる。従って局所的に
\[
 C_1\simeq C_3\simeq\OO_T/(\lambda),\qquad
 C_0\simeq C_2\simeq\OO_T/(\mu).
\]
特に $C_0=C_2=0$ と $\mu$ の可逆性は同値であり、
命題~\ref{prop:qa-relative-elimination} により
$Q_{10}^{\mathrm{odd}}$ の表示を得る。他の二つの場合も同様である。
逆に小さい箙の表示では、対応する生成元層は明らかに零である。
これにより像の表示を得る。

有限表示層 $C_i$ の消滅 locus は、その support の補集合という開集合である。
また cokernel は任意の基底変換と両立する。
従って任意の $T\to\mathfrak M_{12}^{\mathrm{alg}}$ による各関手の
基底変換は、対応する $C_i$ の消滅で定まる $T$ の開部分そのものである。
これは開埋込みの主張と交わりの式を同時に証明する。
\end{proof}

\begin{remark}[最小型、相対型、および追加条件]
\label{rem:qa-open-limits}
十本矢の開部分に八本矢の対象が入るのは、
定義~\ref{def:qa-algebra-stack} が非最小の十本矢表示も許すためである。
幾何点で最小矢がちょうど十本の対象だけを取れば、八本矢の対象は含まれない。
相対的に固定した最小型を扱うときは、残る $C_i$ が階数一の局所自由層で
あることを課す。この条件は局所表示では対応する mass 係数が底上で
零となることを要求し、全幾何点で零という条件より強い。

本節の開性は、同型と基底変換に不変な追加条件を課した共通の
部分スタックへ制限しても保たれる。
ただし AS 条件を満たす部分がどの部分であるか、その部分の代数性、
avatar からの全域被覆、さらに構造対象を保つ内在的圏への射は
この定理の結論に含まれない。
また $\mathfrak M_r^{\mathrm{alg}}$ では四頂点と周期的な graded 表示を
保持している。これを忘れた $\mathcal O$-marked の exact 同値を、
次数付き代数同型と同一視してはならない。
\end{remark}

% !TEX root = ../quiver-avatar-moduli.tex
% 2026-11-21
\section{有限 foundation の二つの十本矢開部分}
\label{sec:qa-finite-foundations}

本節では、矢の本数による開部分を実際の代数スタックとして構成する。
対象は四頂点の有限 directed algebra であり、AS 正則性や標準 heart は
この構成の条件に含めない。最後に数える矢は、有限代数そのものの矢と
最小関係に双対な戻り矢を合わせた relation completion の矢である。
従って、ここで得る開部分の定理と、AS 代数・avatar の相対分類定理とは
異なる主張である。

\subsection{全ての積を保持する有限表示}

\(U,V,W\) を二次元、\(E,F\) を四次元、\(N\) を六次元の
\(\kk\)-ベクトル空間とし、頂点成分を
\[
 A_{01}=U,\quad A_{12}=V,\quad A_{23}=W,\qquad
 A_{02}=E,\quad A_{13}=F,\quad A_{03}=N
\]
と固定する。対角成分は \(A_{ii}=\kk\)、逆向き成分は零とする。
非自明な積は四写像
\[
 \begin{aligned}
 a&:V\otimes U\longrightarrow E,&
 b&:W\otimes V\longrightarrow F,\\
 c&:W\otimes E\longrightarrow N,&
 d&:F\otimes U\longrightarrow N
 \end{aligned}
\]
であり、結合律はただ一つの写像等式
\begin{equation}
 c(1_W\otimes a)=d(b\otimes1_U)
 \quad\text{in }\Hom(W\otimes V\otimes U,N)
 \label{eq:qa-associativity}
\end{equation}
である。この等式の定める閉アフィンスキームを \(Z\) と書く。
\(Z\subset\mathbb A^{128}\) は有限型である。
全ての積を保持するため、後に \(a\) または \(b\) の階数が落ちても
有限代数のベクトル空間としての階数二十四は変わらない。

\begin{definition}[foundation の許容開部分]
\label{def:qa-foundation-open}
\(Z\) の開部分 \(Z^{\mathrm{nd}}\) を次の条件で定める。
\begin{enumerate}
\item \((c,d): (W\otimes E)\oplus(F\otimes U)\to N\) は全射である。
\item 全ての幾何ファイバーで、\(\ker a\) と \(\ker b\) は
非零の分解可能テンソルを含まない。
\end{enumerate}
ここで \(Z^{\mathrm{nd}}\) の記号は上記の有限条件だけを意味し、
AS 正則性を意味しない。
\end{definition}

\begin{lemma}
\label{lem:qa-foundation-open-ranks}
定義~\ref{def:qa-foundation-open} は基底変更と両立する開条件である。
各幾何点で \(\rank a,\rank b\in\{3,4\}\) である。
\end{lemma}
\begin{proof}
全射性は最大 minor の非消滅条件である。
\(a\) について、\(Z\times\PP(V)\times\PP(U)\) 上で
\(a(v\otimes u)=0\) となる閉 incidence を取る。
射影空間の積からの射影は proper なので、その \(Z\) への像は閉である。
その補集合が第二条件を表す。\(b\) についても同じである。

二次元空間のテンソル積の射影化 \(\PP^3\) において、
分解可能テンソルは Segre 二次曲面をなす。
代数閉体上の射影直線はこの二次曲面と交わる。
従って二次元以上の核は非零の分解可能テンソルを含むので、
許される核の次元は零または一である。
\end{proof}

\[
 G=\GL(U)\times\GL(V)\times\GL(W)
       \times\GL(E)\times\GL(F)\times\GL(N)
\]
を各空間の基底変更で作用させる。
頂点の対角単元 \((\lambda_0,\lambda_1,\lambda_2,\lambda_3)\)
が \(A_{ij}\) 上で \(\lambda_j\lambda_i^{-1}\) 倍する作用の像を
\(K\subset G\) と書く。
共通スカラーを除けば \(K\simeq\Gm^3\) であり、
\(K\) は \(Z\) 上自明に作用する中心部分群である。

\begin{proposition}[有限代数の代数スタック]
\label{prop:qa-foundation-stack}
\[
 \mathfrak F^{\mathrm{lin}}_{\le12}=[Z^{\mathrm{nd}}/G],
 \qquad
 \mathfrak F_{\le12}=[Z^{\mathrm{nd}}/(G/K)]
\]
は有限型 Artin スタックである。
前者は上記の階数の頂点成分と積を持つ有限代数のスタックであり、
後者は頂点ごとの比較同型のスカラーによる中心 \(K\) を剛体化したものである。
この命題ではスタックの滑らかさも AS locus の開性も主張しない。
\end{proposition}
\begin{proof}
各頂点成分を局所自由層とする同じデータは、fppf 局所に基底を選ぶと
\(Z^{\mathrm{nd}}\) の点で表される。
代数の同型はこれらの空間の基底変更と正確に一致する。
従って前者の商は指定した moduli を表す。
\(G\) と \(G/K\) は滑らかな有限型代数群なので、
両商は有限型 Artin スタックである。
\(K\) は各対象に中心自己同型として含まれるため、
\([Z^{\mathrm{nd}}/G]\to[Z^{\mathrm{nd}}/(G/K)]\) はその剛体化である。
\end{proof}

\subsection{最小生成元と最小関係の数}

\begin{proposition}[八・十・十二本の有限代数的判定]
\label{prop:qa-foundation-arrow-count}
\(Z^{\mathrm{nd}}\) の幾何点に対応する代数を \(A\) とし、
\[
 \delta_a=4-\rank a,\qquad \delta_b=4-\rank b
 \quad\in\{0,1\}
\]
と置く。\(A\) の最小 directed quiver は
\(6+\delta_a+\delta_b\) 本の矢を持ち、最小関係は
\(2+\delta_a+\delta_b\) 本である。
従って、最小関係ごとに一本の戻り矢を加えた relation completion は
\begin{equation}
 8+2(\delta_a+\delta_b)
 \label{eq:qa-foundation-arrow-count}
\end{equation}
本の矢を持つ。
\end{proposition}
\begin{proof}
隣接頂点間には各二本、合計六本の矢がある。
\(E/\operatorname{im}a\) と \(F/\operatorname{im}b\) の基底を取り、
それぞれ \(0\to2\)、\(1\to3\) の追加生成矢とする。
\((c,d)\) の全射性により、\(0\to3\) の全ての元は
\(0\to2\to3\) または \(0\to1\to3\) を経由する積で生成されるので、
\(0\to3\) の追加生成矢はない。
従って生成矢の数は \(6+\delta_a+\delta_b\) である。

\(0\to2\)、\(1\to3\) の最小関係は、それぞれ
\(K_a=\ker a\)、\(K_b=\ker b\) であり、その次元は
\(\delta_a,\delta_b\) である。
これらの関係を課す前の \(0\to3\) の道の空間は
\[
 (W\otimes V\otimes U)
 \oplus\bigl(W\otimes(E/\operatorname{im}a)\bigr)
 \oplus\bigl((F/\operatorname{im}b)\otimes U\bigr)
\]
という、次元 \(8+2\delta_a+2\delta_b\) の空間である。
短い関係の生成する部分空間は最初の直和成分内の
\[
 (W\otimes K_a)+(K_b\otimes U)
\]
であり、この和は直和である。
実際、両核が一次元で共通の非零テンソルがあれば、
それは \(w\otimes\kappa_a=\kappa_b\otimes u\) と書ける。
\(w\) に対して非零となる線形形式で縮約すると、
\(\kappa_a\) が分解可能になり、定義に反する。
一方の核が零の場合は明らかである。

従って短い関係を課した後の \(0\to3\) 成分は次元八である。
結合律~\eqref{eq:qa-associativity} により \(N\) への写像が定まり、
これは \((c,d)\) の全射性によって全射である。
残る核は二次元であり、ちょうど二本の新たな最小関係を与える。
添字次数は三であるが、追加矢があるときは道長二と三の項を含み得る。
既に短い関係の生成する部分で割っているため、これらは最小関係であり、
四頂点の有向性から、それ以外の始点・終点に新たな関係はない。
これで関係数と \eqref{eq:qa-foundation-arrow-count} を得る。
\end{proof}

添字の差による重みを有限代数の矢 \(i\to j\) に与え、
同じ始点・終点を持つ関係の戻り矢には重み \(4-(j-i)\) を与える。
関係 potential の各項は重み四となる。
隣接矢六本と \(0\to3\) の二つの関係の戻り矢は、
四周期の各隣接頂点を結ぶ二本ずつの重み一の矢を与える。
\(\delta_a=1\) なら \(0\to2\)、\(2\to0\) の重み二の矢を一組、
\(\delta_b=1\) なら \(1\to3\)、\(3\to1\) の重み二の矢を一組加える。
従って重み付きの箙自体は、八本矢、二種類の十本矢、十二本矢の
それぞれの配置に正確に一致する。
この配置の一致は、potential の復元や AS 正則性、
QP の同型と有限代数の同型の比較まで含む主張ではない。

\subsection{二つの開部分とその共通部分}

\begin{theorem}[矢数を上限とする開部分]
\label{thm:qa-foundation-open-diagram}
\(Z^{\mathrm{nd}}\) 上で \(a\)、\(b\) が可逆となる開部分を
それぞれ \(Z_a,Z_b\) と書く。
\[
 \begin{aligned}
 \mathfrak F_{\le10}^{(a)}&=[Z_b/(G/K)],&
 \mathfrak F_{\le10}^{(b)}&=[Z_a/(G/K)],\\
 \mathfrak F_8&=[(Z_a\cap Z_b)/(G/K)]
 \end{aligned}
\]
と置く。上付き \((a)\)、\((b)\) は、階数低下を許す写像を示す。
これらは \(\mathfrak F_{\le12}\) の開部分スタックであり、
\begin{equation}
 \mathfrak F_8
 =\mathfrak F_{\le10}^{(a)}
       \cap\mathfrak F_{\le10}^{(b)}.
 \label{eq:qa-foundation-chart-intersection}
\end{equation}
両十本矢開部分の合併は、最小 relation completion の矢数が
高々十である locus である。
その補集合は、幾何点上では \(\rank a=\rank b=3\) の
最小十二本矢 locus である。
\end{theorem}
\begin{proof}
可逆性は行列式の非消滅であり、\(G/K\) の作用で不変である。
不変開部分からの商の射は representable な open immersion である。
交わりと合併は同じ基底上で計算できる。
矢数との対応は命題~\ref{prop:qa-foundation-arrow-count} による。
\end{proof}

\begin{remark}[最小本数の層と開部分の区別]
\label{rem:qa-exact-arrow-strata}
例えば \(\mathfrak F_{\le10}^{(a)}\) は八本矢の点も含む。
最小十本矢の点だけを取った層は
\(\mathfrak F_{\le10}^{(a)}\setminus\mathfrak F_8\) であり、
八本矢の moduli を開部分として含むのはこの層ではない。
非被約基底について「階数三の層」と呼ぶときは、単なる幾何点条件でなく、
対応する大きな minor の消滅と三次 minor の可逆性を課した
scheme-theoretic な階数 locus を用いる。
本節の開部分は可逆性で定義されており、この区別による曖昧さを持たない。
また、選んだ最小関係とその戻り矢から作る QP の自己同型群を、
有限代数の自己同型群と同一視したわけではない。
QP から有限代数への回復と射の比較には別の証明が必要である。
\end{remark}

\subsection{cubic表示の幾何条件に対する次数六以下の明示的検査}
\label{subsec:qa-explicit-geometric-fourfold-test}

既存のnormal foundation chartで最初の積が同型となる場合、
論文B~\cite[\texttt{prop:cubic-foundation-open}]{QuadricModuli2026}
は二次元空間$V_0,V_1,V_2,V_3$の
四重テンソル$w\in V_0\otimes\cdots\otimes V_3$の
geometric性を使ってcubic逆構成を行う。
以下はこのgeometric開を
\emph{完全に明示的な次数四・六の多項式条件}に書き直す。

\begin{lemma}[四重テンソルのgeometric性のadjugate判定]
\label{lem:qa-explicit-cubic-geometric-adjugate}
$k$を代数閉体とし、各$V_i$を二次元として、
$w\in V_0\otimes V_1\otimes V_2\otimes V_3$
を取る。$i\in\mathbb Z/4\mathbb Z$について、
隣接二因子の縮約行列を
\[
 M_i:(V_i\otimes V_{i+1})^\vee
       \longrightarrow V_{i+2}\otimes V_{i+3}
\]
とする。基底を選んで$M_i$を$4\times4$行列とみなす。
$A_i=\operatorname{adj}(M_i)$をその古典的adjugateとし、
その第$j$列を二因子の$2\times2$テンソルに
reshapeした行列を$C_{ij}$（$1\le j\le4$）とする。
\[
 d_i=\det M_i,\qquad
 h_{ij}=\det C_{ij}
 \tag{Geo-poly}\label{eq:qa-cubic-geometric-minors}
\]
と置く。$d_i$は$w$の成分の次数四、
$h_{ij}$は次数六の多項式である。

Van den Berghのgeometric条件
\[
 (\phi_i\otimes\phi_{i+1})\mathbin{\lrcorner}w\ne0
 \quad\text{for all }i,\quad
 0\ne\phi_i\in V_i^\vee,\ \
 0\ne\phi_{i+1}\in V_{i+1}^\vee
\]
と次の有限条件は同値である。
\[
 \boxed{\quad
  \text{各$i=0,1,2,3$について、}
  \quad d_i\ne0
  \ \text{または}\ 
  h_{i1},h_{i2},h_{i3},h_{i4}
  \text{の少なくとも一つが非零}.\quad}
 \tag{Geo-test}\label{eq:qa-cubic-geometric-finite-test}
\]
従ってgeometric四重テンソル開は、
\emph{四つの次数四式と十六の次数六式だけ}
を使って定義される。
各$M_i$の行・列の基底変更に対し
この非消滅条件は不変なので、
有限foundationの相対chartへ降下する。

これは前段の一般存在定理
~\ref{cor:qa-realizable-foundation-finite-open-criterion}
で未計算だった\emph{全実現可能開}の方程式ではなく、
$F$が可逆であるcubic部分における
geometric条件の具体的な判定式である。
\end{lemma}
\begin{proof}
固定した$i$について、$K_i=\ker M_i$を
四次元空間$(V_i\otimes V_{i+1})^\vee$の部分空間とする。
geometric条件は、$K_i$に非零の
\emph{分解可能}二因子テンソルがないことと同値。
これは
$\PP(K_i)\cap\operatorname{Seg}(\PP^1\times\PP^1)=\varnothing$
という条件である。
Segre像は$\PP^3$の滑らかな二次曲面で、
分解可能な$2\times2$行列は行列式零によって特徴付けられる。

$\operatorname{rank}M_i=4$なら$K_i=0$なので条件は成立し、
$d_i\ne0$と同値である。
$\operatorname{rank}M_i\le2$なら
$\dim K_i\ge2$なので$\PP(K_i)$は少なくとも直線を含む。
代数閉体上で任意の射影直線は二次曲面と交わる
（線上で行列式を制限した二次形式に零点がある、
恒等零の場合も交わる）。
従ってgeometric条件は失敗する。
このとき全ての$3\times3$小行列式が零なので
$A_i=0$となり、全ての$h_{ij}$は零である。

残る$\operatorname{rank}M_i=3$の場合、
$A_i$は階数一の非零行列で、像が
一次元の$K_i$と一致する。
その非零列は$K_i$を張るので、
$K_i$が分解可能テンソルを含まないことは、
いずれかの列をreshapeした$2\times2$行列の
行列式が非零となることと同値である。
これがいずれかの$h_{ij}\ne0$という条件になる。

基底変更で$K_i$は
$(\operatorname{GL}(V_i)\otimes\operatorname{GL}(V_{i+1}))^\vee$
により移り、Segreの分解可能座も不変だから、
判定条件は基底によらない。
四つの$i$について適用すれば結論を得る。
\end{proof}

\subsection{既存の normal foundation chart の位置}

\begin{proposition}[一方の積を可逆にした表示]
\label{prop:qa-existing-foundation-chart}
\(b\) が可逆な開部分では \(F\simeq W\otimes V\) と同定でき、
結合律は \(d\) を \(a,c\) から一意に決める。
従って、この開部分の線形データは既存の normal foundation 表示
\[
 F_{\mathrm{old}}:P\otimes W_1\to W_2,
 \qquad G_{\mathrm{old}}:Q\otimes W_2\to W_3
\]
において
\((P,W_1,W_2,Q,W_3)=(V,U,E,W,N)\)、
\((F_{\mathrm{old}},G_{\mathrm{old}})=(a,c)\) としたものに一致する。
この同定の下で \((c,d)\) の全射性は \(c\) の全射性に等しい。
既存の表示の全開部分と比較する際には、
本節の分解可能な核を持たない条件を追加する。
\end{proposition}
\begin{proof}
\eqref{eq:qa-associativity} の右辺で \(b\otimes1_U\) が同型なので、
\(d=c(1_W\otimes a)(b^{-1}\otimes1_U)\) である。
残る空間と二写像は表示したものに一致する。
\(d\) は \(c\) を経由するので、両者の像の和は \(c\) の像である。
\end{proof}

この比較により、\cite[\texttt{prop:foundation-stack}]{QuadricModuli2026} に現れる一方の積を恒等化した chart と、
本節の二つの積がともに階数三になる境界とを区別できる。
古い二写像表示で \(\rank F_{\mathrm{old}}=2\) として得られる十二本矢表示は、
本節の許容開部分には入らない。
その核は二次元であり分解可能テンソルを含むからである。
また、その十二本矢は一方の対角位置に二組の追加矢を持つのに対し、
本節の最小十二本矢は両対角位置に一組ずつの追加矢を持つ。
矢の総数だけから、これらを同じ型と扱ってはならない。

% !TeX root = ../quiver-avatar-moduli.tex
% 2026-11-22: square / rotated-square の幾何的・有限代数的基礎を本文内で証明
\section{Square と rotated square の自己完結な幾何的基礎}
\label{sec:qa-self-contained-foundations}

この節は，未公刊草稿 \cite{AvatarVSeven2026} の第I--II部から
本稿Cの比較に必要な部分を取り出して，仮定と証明を与え直す．
とくに \emph{反標準曲線上の幾何的 avatar，同型類の相互変換，
有限 directed foundation の復元} は以下の証明だけで読める．
これと，無限周期の AS--Gorenstein 分解，Noether 性，
内在的 Hilbert 関手の認識とは別の定理である．
後者の全証明を有限次元の計算から自動的に推論しない．

\subsection{定義と非被約曲線を含む消滅補題}

$k$ は標数零の代数閉体とする．
$Y$ は，$h^2=2,\ hf=1,\ f^2=0,\ -K_Y=2h$ を満たす
標識付き $\Sigma_0=\PP^1\times\PP^1$
（$h=f+g$）または
$\Sigma_2=\PP_{\PP^1}(\OO\oplus\OO(2))$
（負断面 $S=h-2f$）とする．
$D\in|2h|$ を有効 Cartier 因子とし，
$Y=\Sigma_2$ では $D\cap S=\varnothing$ を課す．
以下，$F=\OO_D(f)$，$H=\OO_D(h)$，
$C=\OO_D(h-2f)$ とする．
$\Sigma_2$ の場合には $C\simeq\OO_D$ である．
曲線は非被約・可約でもよい．

\begin{lemma}[部分曲線と Picard 消滅]
\label{lem:qa-core-curve-vanishing}
上の $D$ について
$\omega_D\simeq\OO_D$，
$h^0(\OO_D)=h^1(\OO_D)=1$ が成り立つ．
真の非零な有効部分曲線 $0<Z<D$ には
$Z(D-Z)\ge2$，したがって $p_a(Z)\le0$ が成り立つ．
さらに $R\in\Pic^0(D)$ とすると，
\begin{enumerate}
\item $R\not\simeq\OO_D$ なら $H^0(D,R)=H^1(D,R)=0$;
\item $FR$ と $(H F^{-1})R$ は大域生成な次数二の
      pencil で，$h^0=2$，$H^1=0$;
\item $HR$ は very ample で，$h^0=4$，$H^1=0$;
\item $CR\not\simeq\OO_D$ なら $H^0(D,CR)=H^1(D,CR)=0$.
\end{enumerate}
この主張は $D$ の滑らかさも被約性も仮定しない．
\end{lemma}
\begin{proof}
随伴公式と
$0\to\omega_Y\to\OO_Y\to\OO_D\to0$
から最初の主張が従う．
$\Sigma_0$ で $Z=af+bg$（$0\le a,b\le2$）なら
$Z(D-Z)=2a+2b-2ab$ であり，$0,D$ 以外は二以上である．
$\Sigma_2$ では $D\cap S=\varnothing$ より
$Z\cdot S=0$ であるから，$Z=aS+bf$ と書けば $b=2a$．
$0<Z<D$ のとき $Z\equiv h$ であり，
$Z(D-Z)=h^2=2$ となる．
随伴公式
$2p_a(Z)-2=-Z(D-Z)$ も得られる．

ここで非被約曲線にも使える二つの観察を証明しておく．
線束 $N$ の切断が generic に非零となる最大の部分曲線を
$Z$ とすると，その切断は
$N|_Z(-(D-Z))$ の generic に単射な切断を誘導する．
ゆえに
\[
 \deg_Z N\ge Z(D-Z).                                      \tag{*}
\]
また長さ $\ell$ の零次元部分 scheme $\xi\subset D$ と
線束 $M$ に対して $H^1(I_\xi M)\ne0$ なら，
Serre 双対性から $I_\xi M\to\omega_D$ の非零射を得る．
像の generic annihilator が定める $Z$ 上に射を因数分解し，
$Z$ の双対化層と余核の長さを比較すると
\[
 \deg_Z M\le 2p_a(Z)-2+\ell .                             \tag{**}
\]
ここで $\omega_Z=\mathcal{H}om_D(\OO_Z,\omega_D)$
という随伴を用いる．両議論は局所的な annihilator と
Cohen--Macaulay 曲線の双対性だけに依存するので，
nilpotent layer を持つ部分曲線にも適用できる．

$R\in\Pic^0(D)$ の各成分次数は零である．
$R$ に切断があれば，真部分 $Z$ への
generic support は $(*)$ に反し，全体への切断なら
零点を持たず $R\simeq\OO_D$ となる．
$\omega_D\simeq\OO_D$ と双対性により (1) が従う．
$F$ と $HF^{-1}$ は各部分曲線上非負次数を持つ．
$\Sigma_2$ では $HF^{-1}\simeq F$ なので同じ主張である．
$(**)$ を $\ell=0,1$ に適用すると
$H^1$ の消滅と長さ一の点の生成が従い，
Riemann--Roch から切断数は二である．
$H$ は真部分曲線上正次数を持つ
（$\Sigma_2$ の唯一の例外 $S$ は $D$ と交わらない）．
$(**)$ を $\ell=0,1,2$ に適用すると $HR$ は
全ての長さ二の cluster を分離するので very ample，
かつ $h^0(HR)=4$ となる．
最後に $\Sigma_2$ では $C=\OO_D$ なので (1) に帰着する．
$\Sigma_0$ では $C=g-f$ であり，
真部分曲線 $Z=af+bg$ に対する
$C\cdot Z=a-b\ge 2a+2b-2ab$ は成立しない．
$(*)$ は $CR$ の切断を排除し，全体への切断は
$CR\simeq\OO_D$ を強制する．
$C^{-1}R^{-1}$ に同じ議論を適用し，双対性で
$H^1(CR)=0$ を得る．
\end{proof}

\subsection{二つの pencil と $\Sigma_2$ への再埋め込み}

\begin{lemma}[独立な pencil と一致する pencil]
\label{lem:qa-core-two-pencils}
$A,B\in\Pic(D)$ を大域生成の次数二線束とし，
$h^0(A)=h^0(B)=2$，
$AB$ を very ample で $h^0(AB)=4$ とする．
\begin{enumerate}
\item $H^*(BA^{-1})=0$ なら，完全 pencil の積は
$D\hookrightarrow\PP H^0(A)^\vee\times\PP H^0(B)^\vee$
という閉埋め込みで，像は $(2,2)$ Cartier 因子である．
\item $B\simeq A$ かつ $p:D\to\PP^1$ が有限なら，
$A^2$ の完全線形系から
$D\hookrightarrow\Sigma_2$ を得る．
像は負断面を避ける $|2h|$ の Cartier 因子であり，
$\OO_D(f)=A,\ \OO_D(h)=A^2$ である．
この埋め込みは標識付き曲面の同型を除き一意である．
\end{enumerate}
\end{lemma}
\begin{proof}
(1) $A$ の評価列を $B$ で tensor した
\[
 0\to BA^{-1}\to H^0(A)\otimes B\to AB\to0
\]
に大域切断を取る．消滅仮定によって
$H^0(A)\otimes H^0(B)\xrightarrow{\sim}H^0(AB)$
である．従って Segre 埋め込みとの合成は
very ample な $|AB|$ であり，もとの積も閉埋め込みである．
Cohen--Macaulay 曲線の滑らかな曲面内への余次元一閉埋め込みは
Cartier であり，各 ruling との交点数二から類は $(2,2)$ となる．

(2) $H^0(A^2)$ は四次元で，
$\Sym^2H^0(A)$ は三次元部分空間となる．
$p$ が有限ゆえ非定数であり，
$\Sym^2H^0(A)\to H^0(A^2)$ は単射である．
補元を $z$ とすると，$A^2$ への商
$\OO_D\oplus p^*\OO_{\PP^1}(2)\to A^2$ は
$(z,1)$ で与えられる．
これは quotient 規約の
$\Sigma_2=\PP_{\PP^1}(\OO\oplus\OO(2))$
への射を定める．第二成分が可逆なので負断面を避ける．
$|h|$ との合成は完全な $|A^2|$ であり，閉埋め込みである．
$fD=2,\ hD=4$ から像は $2h$ である．
$z$ の変更は三角形の束自己同型で吸収される．
\end{proof}

\begin{theorem}[square と rotated square の avatar の相互変換]
\label{thm:qa-core-avatar-square-rotated}
次の二種類の標識付き曲線型 avatar の同型類は自然に全単射する：
\[
\begin{array}{ll}
\text{\rm square:}&
(D\subset\PP^1\times\PP^1;\ F,G,P),\
P\in\Pic^0(D)\setminus\{\OO_D\},\\
\text{\rm rotated:}&
(D\subset Y;\ F,H,P),\
Y=\Sigma_0\text{ または }\Sigma_2,\
P\not\simeq\OO_D,\ C,\
D\cap S=\varnothing\text{ if }Y=\Sigma_2.
\end{array}
\]
ここで square の $F,G$ は順序付き ruling，
rotated の $F,H$ は $f,h$ の制限，$C=HF^{-2}$ である．
同型類の対応は
\begin{equation}
\begin{aligned}
\mathrm{rotated}\longmapsto\mathrm{square}:&
  \quad F_s=F,\quad G_s=HF^{-1}P^{-1},\\
\mathrm{square}\longmapsto\mathrm{rotated}:&
  \quad F_r=F_s,\quad H_r=F_sG_sP .
\end{aligned}
\label{eq:qa-core-rsq-sq}
\end{equation}
さらに rotated の Hilbert 曲面が $\Sigma_2$ となるのは
$G_sP\simeq F_s$ の場合に限る．
\emph{square の許容条件に}
$P\ne F_sG_s^{-1}$ \emph{や}
$P\ne G_sF_s^{-1}$ \emph{を追加してはならない．}
\end{theorem}
\begin{proof}
rotated からは補題~\ref{lem:qa-core-curve-vanishing}より
$F$ と $HF^{-1}P^{-1}$ は二つの大域生成 pencil で，
積 $HP^{-1}$ は very ample である．
その比 $HF^{-2}P^{-1}=CP^{-1}$ は
仮定により非自明で，補題の (4) から acyclic である．
補題~\ref{lem:qa-core-two-pencils}(1) により
$(2,2)$ Cartier 因子へ埋め込める．

逆に square から $F_s$ と $G_sP$ を取る．
これらは大域生成次数二 pencil で，
積 $F_sG_sP$ は very ample である．
$G_sP\not\simeq F_s$ なら，比
$G_sF_s^{-1}P$ は非自明なので
補題~\ref{lem:qa-core-curve-vanishing}(4) から acyclic であり，
$\Sigma_0$ 上に再埋め込みができる．
$G_sP\simeq F_s$ のとき
$F_s^2\simeq F_sG_sP$ は very ample である．
$F_s$ は $G_s$ とともに各成分上正次数を持ち
（$P$ の各成分次数は零），
対応する $D\to\PP^1$ は有限である．
補題~\ref{lem:qa-core-two-pencils}(2) で $\Sigma_2$ を得る．
この場合 $H_r=F_s^2$ なので $C_r=\OO_D$ である．
一般の場合 $C_r=G_sF_s^{-1}P$ であり，
$P=C_r$ なら $F_s=G_s$ となる．
元の $(2,2)$ 埋め込みでは二つの ruling は異なり
$H^*(G_sF_s^{-1})=0$ なので，これは起こらない．
最後に式~\eqref{eq:qa-core-rsq-sq}の合成は
$F_s,G_s$，または $F_r,H_r$ を厳密に回収する．
完全 pencil と補元の一意性により曲面上の同型類も回収する．
\end{proof}

\subsection{有限 directed foundation と反標準行列式}

以下 square の場合には $H_s=F_sG_s$ と書く．
四頂点の directed algebra は，$i<j$ に
$\Hom_D(L_i,L_j)$，対角に $k$，逆向きに零を置き，
積を切断の乗法で定めたものとする．
この定義は $\Perf(D)$ の full dg endomorphism algebra
を意味しない．

\begin{proposition}[二種類の有限 foundation と点曲線の復元]
\label{prop:qa-core-directed-foundations}
曲線型 avatar に対して
\[
\begin{aligned}
B_{\rm sq}&=\End_{\rm dir}(P,F_s,G_s,F_sG_s),\\
B_{\rm rsq}&=\End_{\rm dir}
 (P,C_r,H_rF_r^{-1},H_r^2F_r^{-1})
\end{aligned}
\]
を考える．square の Hom 次元上三角行列は
$\left(\begin{smallmatrix}
1&2&2&4\\0&1&0&2\\0&0&1&2\\0&0&0&1
\end{smallmatrix}\right)$，
rotated のそれは
$\left(\begin{smallmatrix}
1&0&2&6\\0&1&2&6\\0&0&1&4\\0&0&0&1
\end{smallmatrix}\right)$
である．いずれも対応する四頂点箙とその最小関係を与え，
四頂点の点表現を三頂点へ忘れる射は，
関係の普遍評価行列の余核の相対射影化として
\emph{もとの} $D\hookrightarrow Y$ を回収する．
特に square の行列は rank 四から三への Cartier 退化で，
rotated は rank 二から一への退化である．
\end{proposition}
\begin{proof}
線束の差の $H^0$ を補題~\ref{lem:qa-core-curve-vanishing}
で計算すれば表示の Hom 次元を得る．
square で非自明な矢は四つの二次元空間
$H^0(F_sP^{-1}),H^0(G_sP^{-1}),H^0(F_s),H^0(G_s)$
であり，関係は $0\to3$ の四次元空間である．
rotated では $0\to2$ と $1\to2$ の二次元空間，
$2\to3$ の四次元空間が生成矢で，
$0\to3$ と $1\to3$ に各二次元の関係がある．
後者の二つの関係は pencil の評価列から直接得られる．

square での四つの関係を同時に扱うには，
$\PP^1$ の対角の Koszul 分解を二つ外積して
$\PP^1\times\PP^1$ の対角分解を作る．
それを $\iota_*(H_sP)$ に作用させて大域切断を取ると，
補題の消滅と $h^0=2,2,4$ から
\[
 0\longrightarrow\OO_Y(-f)^{\oplus2}
       \oplus\OO_Y(-g)^{\oplus2}
 \longrightarrow\OO_Y^{\oplus4}
 \longrightarrow\iota_*(H_sP)\longrightarrow0
\]
となる．曲面上で双対化し，Cartier 双対性を使うと
\begin{equation}
0\to\OO_Y^{\oplus4}
 \xrightarrow{\kappa_{\rm sq}}
 \OO_Y(g)^{\oplus2}\oplus\OO_Y(f)^{\oplus2}
 \to\iota_*(H_sP^{-1})\to0 .
\label{eq:qa-core-square-matrix}
\end{equation}
大域切断と $(-f),(-g)$ への tensor は
元の四矢空間と四関係を与えるので，
$\kappa_{\rm sq}$ は普遍評価行列である．
余核は局所的に $\OO_Y/(q)$ の自由 rank-one
加群であり，局所行基本変形により
$\kappa_{\rm sq}\sim\operatorname{diag}(1,1,1,q)$．
ゆえに $\PP_Y(\coker\kappa_{\rm sq})=V(q)=D$ である．

rotated では $N=H_rF_r^{-1}P^{-1}$ とおく．
$\Sigma_0$ では同じ積の対角分解と
$H^*(P^{-1})=H^*(C_rP^{-1})=0$ を用い，
$\Sigma_2$ では有限平坦な二重射
$p:D\to\PP^1$ の $p_*N\simeq\OO_{\PP^1}^{\oplus2}$
を用いる．いずれも
\[
0\to\OO_Y(-h)^{\oplus2}\to\OO_Y^{\oplus2}
  \to\iota_*N\to0
\]
となる．$\OO_Y(h)$ で tensor すると
\begin{equation}
0\to\OO_Y^{\oplus2}
 \xrightarrow{\kappa_{\rm rsq}}
 \OO_Y(h)^{\oplus2}
 \to\iota_*(H_r^2F_r^{-1}P^{-1})\to0 .
\label{eq:qa-core-rotated-matrix}
\end{equation}
元の二関係と八道から大域切断で
$\kappa_{\rm rsq}$ を回収する．
局所標準形は $\operatorname{diag}(1,q)$ であり，
その相対射影化は再び $D$ である．
最後の普遍線束の商を
$H_s$ または $H_r^2F_r^{-1}$ で正規化すると
$P$ が回収される．
これは単なる幾何点の行列式の一致ではなく，
余核による scheme-theoretic な復元である．
\end{proof}

\begin{remark}[この節から AS 分類へ進むための残りの入力]
\label{rem:qa-core-as-scope}
上の構成は \emph{任意の}曲線型 avatar に対して成立する．
有限 foundation から全整数添字の AS 正則
$\mathbb Z$-代数を得るには，さらに四つの canonical
kernel--cokernel 列を全周期で構成し，
mutation と Serre 関手を比較する必要がある．
その際に非可換圏の実際の最小射影分解と
AS--Gorenstein 双対性を点検しなければならない．
本稿の他の節が有限 foundation と AS 代数を
区別するのはこのためである．
\end{remark}

% !TeX root = ../nc-quadric-avatar.tex
% 2026-11-22: 幾何的avatarの四つのcanonical列
\section{反標準 avatar の canonical complexes と周期}
\label{sec:qa-self-contained-helix}

この節は前節の幾何的補題から四項の canonical complex を
構成する．これは有限 foundation と全周期の関係を明示する
ためのものであり，任意の相対基底上の AS 正則性と
内在的な心の同定を無条件に主張するものではない．

\begin{proposition}[rotated square の四つの完全列]
\label{prop:qa-core-rotated-complexes}
定理~\ref{thm:qa-core-avatar-square-rotated} の
rotated square 曲線型 avatar を固定し，
\[
\begin{gathered}
L_0=P,\quad L_1=HF^{-2},\quad
L_2=HF^{-1},\quad L_3=H^2F^{-1},\\
L_4=H^2P,\quad L_5=H^3F^{-2},\quad
L_6=H^3F^{-1}P^2,\quad L_7=H^4F^{-1}P^6 .
\end{gathered}
\]
これらの線束の間には directed foundation の評価写像と
最小関係から作られる次の四つの局所分裂完全列がある：
\begin{align}
0&\to L_0\to L_2^{\oplus2}\to L_3^{\oplus2}
     \to L_4\to0,\label{eq:qa-core-rsq-complex0}\\
0&\to L_1\to L_2^{\oplus2}\to L_3^{\oplus2}
     \to L_5\to0,\label{eq:qa-core-rsq-complex1}\\
0&\to L_2\to L_3^{\oplus4}
     \to L_4^{\oplus2}\oplus L_5^{\oplus2}
     \to L_6\to0,\label{eq:qa-core-rsq-complex2}\\
0&\to L_3\to L_4^{\oplus2}\oplus L_5^{\oplus2}
     \to L_6^{\oplus4}\to L_7\to0.
     \label{eq:qa-core-rsq-complex3}
\end{align}
任意の $q\in\mathbb Z$ に対して，$i=0,1,2,3$ とし
\[
\begin{array}{c|cccc}
i&0&1&2&3\\ \hline
G_i&0&h-2f&h-f&2h-f\\
d_i&0&0&2&6\\
\varepsilon_i&1&0&0&0
\end{array}
\]
と書くと，
\begin{equation}
L_{4q+i}=\OO_D(G_i+2qh)
           P^{\,4q(q-1)+d_iq+\varepsilon_i}
\label{eq:qa-core-rsq-periodic-L}
\end{equation}
は同じ四列の全整数平行移動を持つ．
\end{proposition}
\begin{proof}
$A=L_2L_0^{-1}=HF^{-1}P^{-1}$ は
補題~\ref{lem:qa-core-curve-vanishing} により二切断の
大域生成 pencil である．評価列
$0\to A^{-1}\to\OO_D^{\oplus2}\to A\to0$
を $L_2$ で tensor すると
\[
0\to L_0\to L_2^{\oplus2}\to L_2^2L_0^{-1}\to0 .
\]
他方 $L_3(L_2^2L_0^{-1})^{-1}=FP$ も pencil なので，
その評価列を $L_3$ で tensor した
\[
0\to L_2^2L_0^{-1}\to L_3^{\oplus2}\to L_4\to0
\]
と splice すれば第一列となる．
第二列は $L_2L_1^{-1}=F$ と
$L_3H^{-1}=HF^{-1}$ という二 pencil の評価列
\[
0\to L_1\to L_2^{\oplus2}\to H\to0,\qquad
0\to H\to L_3^{\oplus2}\to L_5\to0
\]
を splice して得られる．
ここで $HF^{-1}$ は二切断の大域生成 pencil である．

第三列には，二つの pencil $a_+=FP$，$b_+=HF^{-1}$
を用いる．比 $b_+a_+^{-1}=HF^{-2}P^{-1}$ は
補題~\ref{lem:qa-core-curve-vanishing} より acyclic，
積 $a_+b_+=HP$ は very ample である．
よって補題~\ref{lem:qa-core-two-pencils} で得る
別の $(2,2)$ 埋め込み上に，
命題~\ref{prop:qa-core-directed-foundations} の
四行四列の対角分解を適用する．
その曲線への導来制限では，
Cartier 因子の標準 Tor 項
$\operatorname{Tor}_1(\iota_*Q,\OO_D)=Q(-D)$
が左端に現れる．
$Q=HP^2$ を代入し，全体を $L_3$ で tensor すると
第三列の四項と係数 $4,2+2,1$ が得られる．
これは局所的に長さ一の Cartier 解消を制限して
得た完全列なので，単に Euler 特性の等号ではない．

第四列も，二つの pencil
$a_0=HF^{-1}P$，$b_0=FP^2$ を使う．
その比 $a_0b_0^{-1}=HF^{-2}P^{-1}$ は acyclic，
積は $HP^3$ で very ample である．
同じ四行四列分解を $Q=HP^2$ に適用し，
導来制限と双対化を行った後 $L_6$ で tensor すると
第四列を得る．
第三，第四列の中間写像は対角分解の評価から得られ，
大域切断に戻すと四本の矢と二本ずつの最小関係に一致する．
全列は線束の局所分裂完全列である．

周期については，各連続四項を共通線束で
tensor して最初の項を $P$ に正規化すると，
二つの ruling の役割を担う線束は
$F_q=FP^{2q}$ と $H_q=HP^{4q}$ となる．
$H_qF_q^{-2}=HF^{-2}$ は不変なので，これらは
同じ非自明性と消滅条件を満たす．
前節の二 pencil の再埋め込みを各 $q\in\mathbb Z$
に適用でき，上の四列を繰り返して構成できる．
線束の差分を帰納的に追うと
式~\eqref{eq:qa-core-rsq-periodic-L} を得る．
\end{proof}

\begin{corollary}[square の四つの完全列]
\label{cor:qa-core-square-complexes}
square avatar の $L_0=P,L_1=F,L_2=G,L_3=FG$ に対し
\[
(L_4,L_5,L_6,L_7)
=(F^2G^2P,\ F^3G^2P^2,\ F^2G^3P^2,\ F^3G^3P^4)
\]
とおくと，
\begin{align*}
0&\to L_0\to L_1^{\oplus2}\oplus L_2^{\oplus2}
       \to L_3^{\oplus4}\to L_4\to0,\\
0&\to L_1\to L_3^{\oplus2}\to L_4^{\oplus2}
       \to L_5\to0,\\
0&\to L_2\to L_3^{\oplus2}\to L_4^{\oplus2}
       \to L_6\to0,\\
0&\to L_3\to L_4^{\oplus4}
       \to L_5^{\oplus2}\oplus L_6^{\oplus2}
       \to L_7\to0
\end{align*}
が canonical な局所分裂完全列となる．
\end{corollary}
\begin{proof}
定理~\ref{thm:qa-core-avatar-square-rotated} により，
square の一つ後方または前方の foundation を
rotated の四項と比較する．
正確には rotated の $(L_{-1},L_0,L_1,L_2)$ を
$L_{-1}=F^{-1}P^2$ から出発して
共通線束 $FP^{-1}$ で tensor すると
$(P,F,G,FG)$ となる．
前命題の列を一つ添字移動し，
全項にこの共通線束を tensor すれば四列を得る．
この操作は曲線上の線束の標準化であって，
非可換圏に $(-)\otimes FP^{-1}$ という自己同値を
仮定したものではない．
\end{proof}

\begin{proposition}[全整数 helix の前向き Ext 消滅]
\label{prop:qa-core-forward-ext-vanishing}
上記の square および rotated square の曲線上の線束列について
\[
 \Ext_D^p(L_i,L_j)=0\qquad(p>0,\ i<j)
\]
が成り立つ．曲線が可約・非被約であっても同様である．
これは曲線上の forward acyclicity の主張であり，
この段階で例外 collection の fullness や
非可換 AS--Gorenstein 性を主張するものではない．
\end{proposition}
\begin{proof}
Gorenstein 曲線 $D$ 上で
$\Ext^p_D(L_i,L_j)=H^p(D,L_jL_i^{-1})$ であり，
$p\ge2$ は自動的に零である．
rotated の連続する四商は，$q$ 番目の区間で
\[
 C P^{-1},\quad FP^{2q},\quad HP^{4q},
       \quad FP^{2q+1},
 \qquad C=HF^{-2}.
\]
唯一の総次数零の商 $CP^{-1}$ は
avatar の除外条件 $P\ne C$ と
補題~\ref{lem:qa-core-curve-vanishing}(4) で acyclic である．
その他の商は各成分上非負次数の線束であり，
少なくとも一つ含めば全体上の次数は正である．
複数の連続商の tensor 積は，単独の
$CP^{-1}$ 以外は各部分曲線上非負次数を持ち，
$D$ 全体上正次数を持つ．
ここで $C$ 自体の成分次数は零ではない場合があるが，
$C$ の両側には $F$ または $H$ の商が隣接し，
それらとの積は各部分曲線に非負次数を持つ．
補題~\ref{lem:qa-core-curve-vanishing}の証明で得た
cluster 消滅の不等式 $(**)$ を $\ell=0$ に適用すれば
$H^1=0$ である．

square では四つの連続商は
\[
 FP^{2q-1},\quad GF^{-1},\quad FP^{2q},
       \quad FGP^{4q+1}
\]
である．$GF^{-1}$ の単独の acyclicity は
$\Sigma_0$ 上の Cartier 制限列，
または補題~\ref{lem:qa-core-curve-vanishing}(4)
による．それ以外は上と同様に
nef な正次数の商を少なくとも一つ含み，
$GF^{-1}$ を隣接商と組にすると
各部分曲線上の次数は非負になる．
同じ cluster 判定から $H^1=0$ が従う．
\end{proof}

\begin{remark}[AS 正則性への論理的な境界]
\label{rem:qa-core-helix-not-as}
前命題と系は，両端と中間項を持つ
\emph{実際に完全な}曲線上の複体を与える．
非可換な全 $\mathbb Z$-代数の AS 正則性には，
これらを canonical exceptional helix の
最小射影分解として実現し，
反対向きの $\Ext^3$ による Gorenstein 性を示す
別の導来圏比較が必要である．
この比較のためには acyclic spherical helix と
有限 directed algebra の Serre 関手を用いる．
その\emph{体上の全AS正則$\mathbb Z$-代数への橋}を
第~\ref{sec:qa-as-avatar-classification}節で
最小分解の必要条件・逆構成・相互逆性まで証明する．
ただし相対基底変更と元の O-marked heart の直接同値は
この四列だけからは結論しない．
\end{remark}

% !TeX root = ../nc-quadric-surface-avatar.tex
% 2026-11-22 / AS/avatar classification, source v7 Parts I-II
\section{AS正則\texorpdfstring{$\mathbb Z$}{Z}-代数とavatarの相互逆な分類}
\label{sec:qa-as-avatar-classification}

本節の目的は有限directed foundationの記述を
\emph{全整数頂点のAS正則$\mathbb Z$-代数の同型類}
との全単射へ昇格させることである。
元の証明は\cite[第2--8節，第11--12節]{AvatarVSeven2026}にある。
前二節で証明した曲線上の消滅・評価行列・canonical複体を入力とし、
残っていたAS--Gorenstein性、全helixの一意性と
必要条件の逆向きもここで示す。
一般的なspherical helix実現定理だけは
\cite[Theorem 5]{OkawaUedaSpherical2020}を外部入力とし、
本件でどのようにAS分解を得るかを明記する。

\subsection{最小分解で指定する二型}
\label{subsec:qa-as-minimal-types}

$A=\bigoplus_{j\ge i}A_{ji}$は局所有限、
$A_{ii}=k$、$A_{ji}=0$ ($j<i$)の
positively directed $\mathbb Z$-代数とする。
$P_i=e_iA$、$S_i$を右射影・右単純加群とする。
頂点をすべて保つ代数同型で分類し、
矢の基底や四周期同型の選択は追加データとしない。

\begin{definition}[二つのAS型]
\label{def:qa-as-square-rotated-minimal}
$I=\mathrm{rsq},\mathrm{sq}$に対し
$A\in\mathrm{AS}_{I}$とは、
すべての$i$で最小右射影分解
\begin{equation}
0\to P_{i-4}\to Q_{2,i}\to Q_{1,i}
 \to P_i\to S_i\to0
\label{eq:qa-as-two-minimal-resolutions}
\end{equation}
を持ち、さらに
\begin{equation}
\Ext_A^q(S_i,A)=0\ (q\ne3),\qquad
\Ext_A^3(S_i,A)\simeq {}_\ell S_{i-4}
\label{eq:qa-as-two-gorenstein}
\end{equation}
を満たすものとする。
$Q_{2,i},Q_{1,i}$は次の四周期表で指定する。
\[
\begin{array}{c|c|l|l}
I&i\bmod4&Q_{2,i}&Q_{1,i}\\ \hline
\mathrm{rsq}&0&P_{i-2}^{\oplus2}&P_{i-1}^{\oplus2}\\
&1&P_{i-3}^{\oplus2}&P_{i-2}^{\oplus2}\\
&2&P_{i-3}^{\oplus4}&P_{i-2}^{\oplus2}\oplus P_{i-1}^{\oplus2}\\
&3&P_{i-3}^{\oplus2}\oplus P_{i-2}^{\oplus2}&P_{i-1}^{\oplus4}\\ \hline
\mathrm{sq}&0&P_{i-3}^{\oplus2}\oplus P_{i-2}^{\oplus2}&P_{i-1}^{\oplus4}\\
&1&P_{i-2}^{\oplus2}&P_{i-1}^{\oplus2}\\
&2&P_{i-3}^{\oplus2}&P_{i-2}^{\oplus2}\\
&3&P_{i-3}^{\oplus4}&P_{i-2}^{\oplus2}\oplus P_{i-1}^{\oplus2}
\end{array}
\]
左右Noether性・幾何性・周期のtrivializationは仮定しない。
\end{definition}

\begin{lemma}[局所双対性と全$\mathbb Z$-代数の回収]
\label{lem:qa-as-serre-unique-recovery}
$A\in\mathrm{AS}_I$、$E_i=\pi P_i\in\Qgr A$とする。
$\mathcal D_{\rm lin}=\operatorname{thick}\{E_i\mid i\in\mathbb Z\}$
において連続する四対象はstrong exceptional tilting generatorである。
そのEnd代数$B$は大域次元二で、
Serre関手$S_B$により
\begin{equation}
E_{i+4}\simeq S_B^{-1}(E_i)[2],\qquad
A_{ji}=\Hom(E_i,E_j)\quad(j\ge i)
\label{eq:qa-as-serre-completion}
\end{equation}
を得る。全代数は初期四頂点のEnd代数から
頂点を保つ同型を除いて一意に決まる。
その数値Euler行列は
\[
G_{\rm rsq}=\begin{pmatrix}
1&0&2&6\\0&1&2&6\\0&0&1&4\\0&0&0&1
\end{pmatrix},\quad
G_{\rm sq}=\begin{pmatrix}
1&2&2&4\\0&1&0&2\\0&0&1&2\\0&0&0&1
\end{pmatrix},
\]
かつ$\dim A_{i+4m,i}=(2m+1)^2$である。
\end{lemma}
\begin{proof}
最小分解を各頂点で評価したEuler漸化式は、
負の成分が零で対角が$k$という初期条件から
全Hom次元を一意に定める。
可換基準collection
$(\OO,\OO(1,0),\OO(0,1),\OO(1,1))$と
$(\OO,\OO(-1,1),\OO(0,1),\OO(1,2))$
にも同じ漸化式が成立し、上の行列を得る。
四周期のHilbert関数は
$h^0(\PP^1\times\PP^1,\OO(2m,2m))=(2m+1)^2$。
これを可換基準モデルとの類比だけに依存させず、
最小分解表から有理母関数として直接検算する。
$r=i\bmod4$、$h_{r,n}=\dim A_{i,i-n}$、
$H_r(z)=\sum_{n\ge0}h_{r,n}z^n$とおき、
$Q_{1,i},Q_{2,i}$中の$P_{i-k}$の重複度を
それぞれ$a_{r,k},b_{r,k}$とする。
表のEuler漸化式は
\[
 (1-z^4)H_r(z)
 =1+\sum_{k=1}^3
 (a_{r,k}-b_{r,k})z^kH_{r-k\bmod4}(z).
 \tag{Euler4}\label{eq:qa-as-four-phase-euler}
\]
これら四本の一次方程式の有理関数解について、
$\zeta^4=1$を満たす原始四乗根$\zeta$を取ると、
各$r$で
\begin{equation}
 \frac14\sum_{j=0}^3H_r(\zeta^jz)
 =\frac{1+6z^4+z^8}{(1-z^4)^3}
 =\sum_{m\ge0}(2m+1)^2z^{4m}.
 \label{eq:qa-as-four-phase-rational-proof}
\end{equation}
左辺は四の倍数次数への射影である。
式の有理関数恒等性は

\noindent\texttt{computations/verify\_as\_avatar\_euler.py}で
四元連立の厳密な逆行列計算として再現できる。
ここで確認したのは\emph{最小分解が完全である場合の
Euler数値帰結}であって、
分解の完全性やAS--Gorenstein性を
母関数だけで証明したのではない。

まず有限区間で局所双対の消滅を確認する。
$P_j$ を頂点 $b$ 未満の部分加群で割った有限長商を
$U_{j,b}$ とする。頂点によるフィルトレーションの
因子は $S_\ell^{\oplus\dim A_{j\ell}}$（$b\le\ell\le j$）である。
\eqref{eq:qa-as-two-gorenstein} を射影 $P_i$ の成分で読むと、
$\Ext_A^q(S_\ell,P_i)$ は $(q,\ell)=(3,i+4)$ の場合だけ
$k$ となる。従ってフィルトレーションの長完全列から
$\Ext_A^q(U_{j,b},P_i)$ は $q\ne3$ では零であり、
$q=3$ の次元は $b\le i+4\le j$ なら $\dim A_{j,i+4}$、
それ以外は零である。
さらに切断を下げるとき、追加される因子の
$\Ext_A^3(-,P_i)$ は零なので、$b\le i+4$ 以後の
遷移写像は同型となる。
従ってこの計算はフィルトレーションの Euler 特性だけでなく、
導来 torsion を計算する直極限の各頂点での消滅を与える。

加群としての同定には、有限長加群 $N$ を
頂点 $i+4$ 未満の部分と $i+4$ 以下の部分で切る。
その二つの外側の部分の $\Ext_A^\bullet(-,P_i)$ は上の計算で零である。
従って包含と商の誘導する同型により
$\Ext_A^3(N,P_i)\simeq\Hom_k(N_{i+4},k)$ を得る。
この同型は $\Ext_A^3(S_{i+4},P_i)\simeq k$ の同定を固定すれば
$N$ に自然である。
特に $U_{j,b}$ に適用した同型は $j$ を変える道の作用と可換であり、
直極限で $R\tau(P_i)\simeq D(Ae_{i+4})[-3]$ となる。
この計算を tails に移すため、$J_{i,b}$ を
$P_i$ の頂点 $b$ 未満の部分加群とする。
$P_i/J_{i,b}$ は有限長で、各単純の長さ三の解消から
有限生成射影による長さ高々三の解消を持つ。
$J_{i,b}$ も有限生成射影による長さ高々二の解消を持つ。
さらに上の添字規約では、$J_{i,b}$ の射影解消を
頂点 $b$ 未満の射影の直和で選べる。
従って下に有界な torsion 加群 $N$ に対して、
十分小さい $b$ で $\Ext_A^p(J_{i,b},N)=0$ が全 $p$ で成り立つ。
一般の torsion 加群はそのような部分加群の filtered colimit
であり、有限射影解消により同じ極限消滅が成り立つ。
従って Gabriel 商の section 関手には
論文Bの \cite[式 \texttt{eq:cubic-tails-section-colimit} と
\texttt{eq:cubic-tails-derived-localization}]{QuadricModuli2026}
と同じ局所化三角を適用できる。
ここで用いたのは有限区間の射影解消と torsion 上の消滅である。

この三角を $P_j$ に適用して頂点 $i$ で評価すると、
\[
 \Hom(E_i,E_j)=A_{ji},\qquad
 \Ext^2(E_i,E_j)=D(A_{i,j+4}),\qquad
 \Ext^q(E_i,E_j)=0\quad(q>0,\ q\ne2)
\]
を得る。特に長さ四の窓の全対で高次 Ext が消える。
最小分解を$\pi$で移した四項完全列を
正負双方へ繰り返すと連続四対象は全$E_j$を生成する。
有限quiverの最小関係はsourceからsinkへ向かい、
関係同士のoverlapを持たない。
よってその有限End代数の単純分解の長さは高々二。
$E_j$ を末項とする四項窓の直和を $T$ とする。
$B=\End(T)$ は有限次元かつ有限大域次元なので、
導来 Morita 同値で Serre 関手は $DB$ との導来テンソルに対応する。
この窓の各 $E_i$ について $\Hom(E_i,E_{j-4})=0$ であり、
先の局所双対から
\[
 \RHom(E_i,E_{j-4}[2])\simeq D\Hom(E_j,E_i)[0]
\]
となる。この式は $i$ の窓内の射に自然である。
実際、$\Ext^2$ の同定は有限切断の $i$ 番射影への
道の作用を双対化したものである。
従って $\RHom(T,E_{j-4}[2])$ は、ただ一つのコホモロジーとして
右 $B$ 加群 $D\Hom(E_j,T)$ を持つ。
これは $S_B(E_j)$ に対応する加群であり、
Morita 同値により $S_B(E_j)\simeq E_{j-4}[2]$ を得る。
窓を変えても生成する圏は同じなので、その Serre 関手は同一である。
よって全 $j$ で $S_B^{-1}[2](E_j)\simeq E_{j+4}$ となる。
有限窓の導来圏でこの自己同値を反復し、合成を通常の射の合成に
取れば、全整数の代数の復元 \eqref{eq:qa-as-serre-completion} を得る。
これは周期同型の選択を分類データへ加える主張ではない。
\end{proof}

\subsection{AS条件からavatarを復元する}
\label{subsec:qa-as-to-avatar}

\begin{lemma}[AS条件が強制する階数]
\label{lem:qa-as-required-ranks}
rotated型の有限quiverを
$0\overset{2}{\longrightarrow}2
 \overset{4}{\longrightarrow}3$、
$1\overset{2}{\longrightarrow}2$と書き、
$R_0,R_1$を二つの二次元関係空間とする。
$V=H^0(2\to3)$とし、
二つの関係評価行列を$M_0(v),M_1(v)$、
随伴関係の$4\times4$行列を
$T(c,d),\widetilde T(a,b)$とすると
\begin{enumerate}
\item $R_0,R_1\subset V\otimes k^2$に
 非零の分解可能テンソルはない。
\item $\det M_0(v)=\det M_1(v)=0$、
 $v\ne0$なら両行列の階数は一。
\item $c,d$がともに非零なら
 $\operatorname{rk}T(c,d)\ge3$、
 $a,b$がともに非零なら
 $\operatorname{rk}\widetilde T(a,b)\ge3$。
\end{enumerate}
\end{lemma}
\begin{proof}
まず無限の関係を有限窓から回収する。
有限窓の代数 $B$ は大域次元二であり、$R_0,R_1$ は
有向道代数の関係 ideal の最小生成である。
未変形の Keller の relation completion
\cite[Theorem 6.10]{KellerDeformedCY2011} を用いると、
$T_B\RHom_{B^e}(B,B^e)[2]$ の零次コホモロジーは
有限窓の関係に双対戻り矢を加えた通常の Jacobian 代数である。
ここでは元の矢のテンソル次数を零、戻り矢を一とする。
この次数は一窓をまたぐ回数であり、関係の巡回微分は
両辺で同じ回数を保つ。
逆 Serre 双加群 $\RHom_{B^e}(B,B^e)[2]$ は
$S_B^{-1}[2]$ を表すので、その $m$ 次テンソル成分の
頂点対 $(i,j)$ は前補題により $\Hom(E_i,E_{j+4m})$ になる。
積は双加群のテンソル積、すなわち添字を移した射の合成である。
従ってこの Jacobian 代数を四頂点ずつ展開すると元の $A$ を得る。
以下の周期商で使う関係はこの通常の道 ideal の全生成である。

四周期の有限quiverのsourceとsinkを連結し、
各ブロック内の元の関係とその随伴微分関係を
すべて課した商を考える。
一ブロックのboundary cornerを$W$とすると、
$m$ブロックのcornerは$W^{\otimes m}$である。
境界間の道は一意に分割される。
以下では各場合ごとに境界を選び、
境界をまたぐ関係の像が零であることを確認する。
そのとき残る関係は各ブロックの一因子だけに作用するので、
ブロックごとの商のテンソル積が全境界 corner となる。
$\dim W\ge2$なら成分次元は少なくとも$2^m$となり、
前補題の$(2m+1)^2$に反する。

使う関係を先に書く。
$v_0,\dots,v_3$を$V$の基底とし、
$r_{0,\alpha}=\sum_{\ell,i}A_{\ell i\alpha}v_\ell a_i$、
$r_{1,\beta}=\sum_{\ell,j}B_{\ell j\beta}v_\ell b_j$
（$\alpha,\beta=0,1$）とする。
次周期の戻り矢$c_\alpha,d_\beta$を加えると、
各$r_{0,\alpha},r_{1,\beta}$に加えて
\[
\begin{aligned}
 \sum_{\ell,\alpha}A_{\ell i\alpha}c_\alpha v_\ell&=0
 &&(i=0,1),\\
 \sum_{\ell,\beta}B_{\ell j\beta}d_\beta v_\ell&=0
 &&(j=0,1),\\
 \sum_{i,\alpha}A_{\ell i\alpha}a_i'c_\alpha+
 \sum_{j,\beta}B_{\ell j\beta}b_j'd_\beta&=0
 &&(\ell=0,1,2,3)
\end{aligned}
\tag{Adj-rel}\label{eq:qa-as-forcing-adjoint-relations}
\]
および周期移動を持つ。以下は消した矢で切った
\emph{全随伴関係の像}を含む商である。

(1) 分解可能な関係を$r_{0,0}=v_0a_0$へ正規化する。
他方の枝の頂点 $1$ とその周期移動を零にし、
$v_0=a_0=c_1=0$を課す。境界は各周期の頂点 $3$ とする。
式~\eqref{eq:qa-as-forcing-adjoint-relations}の
$c_0$係数は$v_0$か$a_0$を含み、
$c_1$係数は零となるため、境界をまたぐ $c v$ 型の関係は消える。
残る境界から次の境界への道は三本で、
$r_{0,1}'c_0$による関係は高々一本。
従って$\dim W\ge3-1=2$である。
$R_1$の場合は二枝を交換する。

(2) $M_0(v)=0$かつ$\det M_1(v)=0$とする。
$v=v_0$とし、$M_1(v_0)$の左右kernelに
沿って$b_0,d_0$を選ぶ。
$v_1=v_2=v_3=b_1=d_1=0$を課し、境界を各周期の頂点 $2$ とする。
境界をまたぐ $v a,v b$ 型の $r_0,r_1$ はこの切断で消える。
$c v,d v$ 型の関係も、$M_0(v_0)=0$ と
$d_0$ の kernel 条件により消える。
$a'c,b'd$ 型のうち $\ell=0$ の式は $b_0$ の kernel 条件から零となる。
四本の$a_i'c_\alpha$と一本の$b_0'd_0$、
合計五本の境界道に、
$\ell=1,2,3$の高々三関係だけが残る。
よって$\dim W\ge5-3=2$である。
$M_1(v)=0$の場合も同様。

(3) $\operatorname{rk}T(c,d)\le2$なら
$c=d=(1,0)$に正規化し、
$A_{\ell i0},B_{\ell j0}$の四係数列が
$\langle v_0,v_1\rangle$に入るようにできる。
境界を各周期の頂点 $3$ とし、$v_0=v_1=c_1=d_1=0$を課す。
境界をまたぐ $c v,d v$ 型の関係は係数列が
$\langle v_0,v_1\rangle$ に入るため消える。
残った$a_0'c_0,a_1'c_0,b_0'd_0,b_1'd_0$を
高々二関係で割る空間$U'$は
$\dim U'\ge2$である。
次の境界までの道空間は
\[
 W=(\langle v_2',v_3'\rangle\otimes U')/
       \langle r_{0,1}'c_0,r_{1,1}'d_0\rangle
\]
であるから$\dim W\ge2\dim U'-2\ge2$。
$c_0,d_0$を可逆元や恒等射と扱っていないので、
この次元計算は四周期の正当な次数付き商に対するもの。
各場合で$\dim W\ge2$は指数増大に矛盾する。
左最小分解に同じ議論を適用すると
$\widetilde T$の階数条件を得る。
\end{proof}

\begin{proposition}[一点表示の幾何]
\label{prop:qa-as-extraction}
任意の$A\in\mathrm{AS}_{\rm rsq}$から
\[
Y=\PP_{\PP^1}(\mathcal E)\simeq\Sigma_0\ \text{または}\ \Sigma_2,
\quad
D\hookrightarrow Y,\quad
P\in\Pic^0(D)
\]
を回収できる。曲線型では
$D\in|2h|$、$\Sigma_2$なら$D\cap S=\varnothing$、
かつ$P\ne\OO_D,\OO_D(h-2f)$となる。
初期四項は
\[
(P,\OO_D(h-2f),\OO_D(h-f),\OO_D(2h-f)).
\]
$D=Y$の場合には$\Sigma_0$の古典型一点を得る。
\end{proposition}
\begin{proof}
保持三頂点の安定表現は
$\{([b],[v])\in\PP^1\times\PP^3\mid M_1(v)b=0\}$
である。$[b]$ を固定すると、関係は $V$ の二つの一次式となる。
この二式が従属なら、$R_1$ の非零元の $b$ による縮約が零となり、
その元は $V\otimes\ker b$ に属する。
$\ker b$ は一次元なので、これは分解可能テンソルを与え、
前補題に反する。
よって普遍関係写像
$\OO(-1)^2\to V\otimes\OO_{\PP^1}$
は各点で階数二であり、単元 minor により局所分裂単射である。
その余核 $\mathcal E$ は rank 二、次数二の大域生成束であり、
$M_1(v)b=0$ の incidence は $\PP(\mathcal E)$ と
試験スキーム上の商線束の関手として同型である。
$\mathcal E$ の分裂次数は大域生成性から非負、和は二なので
$\OO(1)^2$ または $\OO\oplus\OO(2)$ となり、
$Y$ の二型を得る。
忘れた頂点を加える評価は
$\kappa:\OO_Y^2\to\OO_Y(h)^2$である。
零点では $\det M_0(v)=\det M_1(v)=0$ であり、
前補題により $M_0(v)$ の階数は一なので、
局所的に$\kappa\sim\operatorname{diag}(1,q)$。
従って$\PP_Y(\coker\kappa)=V(q)=D$は
scheme-theoretic閉埋め込みであり、
$\det\kappa$の類は$2h$。
$\Sigma_2$ の負断面 $S$ は
$\mathcal E=\OO\oplus\OO(2)\twoheadrightarrow\OO$ の商に対応する。
この商では $\OO_Y(h)|_S\simeq\OO_S$ であり、
普遍評価 $V\otimes\OO_S\to\OO_S$ は一つの固定商 $[v_*]$ となる。
従って incidence 式 $M_1(v_*)b=0$ はすべての $[b]\in\PP^1$ で成り立つ。
二つの独立な $b$ を代入すると $M_1(v_*)=0$ を得る。
もし $\det M_0(v_*)=0$ なら、両行列式が零なのに
$M_1$ が階数零となり、前補題の第二条件に反する。
従って $\det M_0$ は $S$ に沿って単元であり、
スキームとして $D\cap S=\varnothing$ である。
この議論は $\Sigma_2$ 上で $\det\kappa$ が恒等的に零となる場合も排除する。
普遍右表現の双対と余核$Q$を用いて
$P=Q^{-1}\OO_D(2h-f)$と正規化する。

点表示の移動を随伴関係
\eqref{eq:qa-as-forcing-adjoint-relations} で記述する。
既知の評価を $v,a,b$ とすると、$c,d$ は
$M_0(v),M_1(v)$ の関係空間側の核線から得る。
いずれも階数一なので核線は局所直和因子であり、
単元 minor による記述は基底変更と可換である。
$c,d$ を代入した $T(c,d)$ では、既知の $v$ が左核に入るため
行列式は零となり、前補題の第三条件によって階数は三である。
従って右核から次の組 $(a',b')$ が一意に定まる。
その一方が零なら、他方と $c$ または $d$ によって
非零の関係テンソルの縮約が零となる。
この関係は二次元矢空間の一次元部分に支持されるので
分解可能となり、前補題の第一条件に反する。
従って $a',b'$ はともに非零である。

次に $\widetilde T(a',b')$ には既知の $(c,d)$ が核を与える。
同じ階数三の議論により、新しい $v'$ が一意に定まる。
随伴関係は $M_0(v')a'=M_1(v')b'=0$ を含むので、
$(v',a',b')$ は再び元の四頂点の点表示を与える。
逆向きにも同じ核線を用いると、階数一と階数三での
一意性により元の表示に戻る。
すべての写像は単元 minor 上の核線で構成されるので
逆性はスキーム上で成立する。
この移動の自己同型を $\sigma$ とし、
矢を $L_i\to L_j$ の評価と書くために普遍右表現を双対化する。
初期の線束を $L_0=P,L_1=HF^{-2},L_2=HF^{-1},L_3=H^2F^{-1}$
と正規化すると、二つの行列の余核は $Q$ および $HF$ である。
後者は $b:B\otimes\OO\twoheadrightarrow F$ の核に $H$ を捩った
部分束が $M_1(v)$ の像になることから従う。
したがって二つの階数一行列の核は、
行列式を取るとそれぞれ $QH^{-2}$ と $FH^{-1}$ である。
その双対が $L_4L_3^{-1}$ と $L_5L_3^{-1}$ を与えるので、
$L_4=H^2P$、$L_5=H^3F^{-2}$ を得る。

続く二つの階数三の写像では、
核と余核の線束を含む完全列の型は
\[
\begin{aligned}
0&\to L_6^{-1}\to
 (L_4^{-1})^{\oplus2}\oplus (L_5^{-1})^{\oplus2}
 \xrightarrow{T(c,d)} (L_3^{-1})^{\oplus4}
 \to L_2^{-1}\to0,\\
0&\to L_7^{-1}\to (L_6^{-1})^{\oplus4}
 \to (L_4^{-1})^{\oplus2}\oplus (L_5^{-1})^{\oplus2}
 \to L_3^{-1}\to0 .
\end{aligned}
\]
最初の余核は既知の $v$、次の余核は既知の $(c,d)$ の評価である。
二次元・四次元の定数係数空間の行列式を除けば、
$L_6=L_4^2L_5^2L_2L_3^{-4}=H^3F^{-1}P^2$ と
$L_7=L_6^4L_3L_4^{-2}L_5^{-2}=H^4F^{-1}P^6$ となる。
よって移動後の偏極は
$\sigma^*F=L_6L_5^{-1}=FP^2$ と
$\sigma^*H=L_7L_6^{-1}=HP^4$ である。
曲線型では、二つの式から $\sigma^*(HF^{-2})=HF^{-2}$ となる。
$\Sigma_0$ 上で成分の類を $af+bg$ と書くと、
$HF^{-2}$ のその成分上の次数は $a-b$ である。
この値が同じで類が異なる二成分が共存すると仮定すると、
一方の二係数は他方よりともに少なくとも一大きい。
小さい方も非零の有効類なので、二成分の和のいずれかの
係数が三以上となり、$D\in|2f+2g|$ に反する。
従って $\sigma$ は各成分の数値類を保ち、
特に $F$ の多重次数を保つ。
$\Sigma_2$ では各成分が負断面を避けるため、その類は
正整数倍の $h$ であり、類の異なる二成分は $2h$ 内に共存しない。
この場合も $F$ の多重次数は保たれる。
従って $\sigma^*F=FP^2$ を各成分で比較して $\deg P=0$ を得る。
ここで次数は被約既約成分への制限で測る。
多重次数零が $\Pic^0(D)$ を与える理由は、
まず $D_{\mathrm{red}}$ の正規化に移すと、
各正規化成分の Picard 群の単位成分が次数零部分だからである。
正規化による Picard 群の核は、導手の有限スキーム上の
単元群の商として記述され、連結である。
実際、代数閉体上の有限局所代数の単元群は
剰余体の乗法群と冪零イデアルに由来する連結な群からなる。
さらに $D$ の冪零イデアル $\mathcal J$ に対し
$1+\mathcal J$ を冪で濾過すると、逐次商は加法的な準連接層
$\mathcal J^r/\mathcal J^{r+1}$ である。
曲線上でこれらの $H^2$ は零で、$H^1$ はベクトル群なので、
被約化への Picard 制限は全射で核は連結である。
従って $P$ は非被約の場合も $\Pic^0(D)$ に属する。
曲線型では $\kappa$ の行列式は非零で、
曲面上の二項解消の余核は $\iota_*Q$ である。
これを $-2h+f$ と $-h-f$ で捩ると、余核はそれぞれ
$\iota_*P^{-1}$ と $\iota_*(HF^{-2}P^{-1})$ となる。
$\Sigma_0$ では $h=f+g$ なので、一つ目の捩れの
二つの曲面項はそれぞれ $-f-2g,-g$、
二つ目は $-2f-g,-f$ であり、Künneth によってすべて acyclic である。
$\Sigma_2$ では、相対次数 $-1$ の項の順像は全次数で零であり、
$-f$ の順像は $\OO_{\PP^1}(-1)$ である。
残る $-2h+f$ については相対双対性と
$\omega_\pi=\OO_Y(-2h+2f)$ により
$R^1\pi_*\OO_Y(-2h+f)=\OO_{\PP^1}(-1)$、$\pi_*=0$ となる。
従って両曲面で
$H^\bullet(P^{-1})=H^\bullet(HF^{-2}P^{-1})=0$ を得る。
自明線束には単位切断があるので、二つの除外条件が従う。
$\det\kappa\equiv0$ なら、負断面での可逆性により
$\Sigma_2$ は起こらない。
$Y=\Sigma_0$ 上で $\kappa$ は至るところ階数一なので、
余核 $Q$ は $Y$ 上の線束であり、点表示のスキームは $Y$ 全体である。
この場合の四周期の自己同型 $\sigma$ は
$\omega_Y^{-1}\simeq H^2$ を保つ。
$\Pic(Y)\simeq\mathbb Z^2$ には二捩れがないので、
$H$ は反標準束の一意な平方根として $\sigma$ に保たれる。
従って四周期の式 $\sigma^*H=HP^4$ から $P^4=\OO_Y$ を得て、
同じ Picard 群の計算から $P=\OO_Y$ となる。
この場合の線形代数には曲線型の単射解消を使わない。
\end{proof}


\subsection{Avatarからの全AS代数の逆構成}
\label{subsec:qa-as-avatar-inverse}

\begin{lemma}[canonical truncation]
\label{lem:qa-as-canonical-truncation}
rotated square avatarから前節の$L_i$を全整数に作る。
任意の$i$でcanonical四項列を
初期四頂点の共変Yonedaへ制限すると、
頂点単純加群の最小射影分解になる。
特に三回の右mutationは
\begin{equation}
T^{-1}_{L_{i+3}}T^{-1}_{L_{i+2}}T^{-1}_{L_{i+1}}(L_i)
 \simeq L_{i+4}[-2]
\label{eq:qa-as-canonical-mutation}
\end{equation}
を満たす。
\end{lemma}
\begin{proof}
$pencil$の評価列で作った
\eqref{eq:qa-core-rsq-complex0}と
\eqref{eq:qa-core-rsq-complex1}は
Yonedaで最初の二種の最小関係と一次syzygyを与える。
後二列は四行四列の対角分解の
Cartier曲線への導来制限と双対化から来るので
Tor項を含めて完全である。
各位相で最初に現れる非零な矢空間への評価は、
pencil または四次線束の全切断空間を使うので全射である。
第一・第二位相の残る頂点では、第一 pencil の評価列を
その頂点の線束で双対化して用いる。
評価の余核は pencil trick により次数二の線束の $H^1$ に
含まれるので零となる。第三位相の後二頂点では、四次線束 $H$
に対して $a_+$ または $b_+$ の評価列を用いる。
乗法の余核は $H a_+^{-1}$ または $H b_+^{-1}$ の $H^1$ に
含まれ、これらも次数二の pencil なので零である。
第四位相の最後の頂点では、第四列を作る前の曲面上の
$0\to\OO^{\oplus4}\to\OO(a)^{\oplus2}\oplus\OO(b)^{\oplus2}
 \to\iota_*Q\to0$
の大域切断が、$H^1(\OO)=0$ により所要の全射を与える。
中間の完全性もEuler特性だけには依存しない。
canonical複体を一時的に
$0\to L_i\to M_1\to M_2\to L_{i+4}\to0$と書く。
共変Yonedaの頂点$b$における
$\Hom(M_1,L_b)\to\Hom(L_i,L_b)$のkernelの次元を
計算すると
\[
\begin{array}{c|c|ccc}
i\bmod4&\text{相対次数}&b=1&b=2&b=3\\ \hline
0&(0,0,2,6)&0&0&2\\
1&(0,2,6,8)&0&2&4\\
2&(0,4,6,6)&0&2&2\\
3&(0,2,2,4)&0&0&4
\end{array}
\]
となる。ここで $L_b$ は窓の第 $b$ 頂点の略記である。
窓内の逆向きの通常の Hom は、
命題~\ref{prop:qa-core-forward-ext-vanishing}と
$\omega_D\simeq\OO_D$ による Serre 双対で零である。
従って通常の Hom を使ったこの計算は directed Yoneda の評価と一致する。
同じ双対性により、窓の外の端についても
$\Hom_D(L_{i+4},L_b)=\Ext_D^1(L_b,L_{i+4})^\vee=0$ となる。
特に総次数が等しい対でも、前向き消滅の次数零の場合がこれを保証する。

短完全列 $0\to\operatorname{coker}(L_i\to M_1)\to
M_2\to L_{i+4}\to0$ に $\Hom_D(-,L_b)$ を適用すると、
$\Hom_D(M_2,L_b)$ から $\Hom_D(M_1,L_b)$ への写像は単射である。
像は表の kernel に含まれ、その次元は表の数と一致する。
従って窓の各頂点で中間は完全である。
頂点 $i$ では対角の $k$ が augmentation として残り、
すべての微分は異なる頂点間の成分からなるので解消は最小である。
三回のmutationは頂点1,2,3を順次消す
Yoneda truncationなので、
$D$上に実現すると最後の$L_{i+4}[-2]$のみが残る。
\end{proof}

\begin{proposition}[spherical helixによるAS正則性]
\label{prop:qa-as-helix-reconstruction}
任意の曲線型rotated square avatarから
$B=\End_{\rm dir}(L_0,L_1,L_2,L_3)$を作る。
その標準射影対象を$E_0,\ldots,E_3$とし、
$E_{i+4}=S_B^{-1}(E_i)[2]$で延長すると
\[
 A(D,P)_{ji}=\Hom(E_i,E_j)\quad(j\ge i)
\]
は$\mathrm{AS}_{\mathrm{rsq}}$に属する。
\end{proposition}
\begin{proof}
$\omega_D=\OO_D$かつ
$H^0(\OO_D)=H^1(\OO_D)=k$だから
各$L_i$は一次元spherical objectである。
全整数のcanonical完全列と
命題~\ref{prop:qa-core-forward-ext-vanishing}、
前補題はacyclic spherical helixを与える。
このacyclicityは同文献Definition~4の全$i<j$に対する
次数零以外のHom消滅である。
これを入力として、対応するexceptional helixの
前向きHomが次数零に集中することに
\cite[Theorem 5]{OkawaUedaSpherical2020}を用いる。
指定した最小分解の頂点・重複度と双対の同定には、
以下のcanonical複体とSerre traceの計算を使う。

$C$を四線束のfull dg Endとする。
$B$は前向き次数零Homと対角の恒等射を含み、
$B\to C$ は dg の実現写像である。
$H^0(C)$ の後ろ向き成分は前向き Ext 消滅の双対により零で、
$H^0(B)\to H^0(C)$ は同型となる。
その cone の唯一のコホモロジーは $H^1(C)$ である。
Serre trace による pairing は合成に関して自然であるから、
$H^1(C)=DB$ は $B$ 双加群としての同型である。
単一のコホモロジーをもつ cone は双加群の導来圏で
$DB[-1]$ に同型なので、三角形
\begin{equation}
 DB[-2]\longrightarrow B\longrightarrow C
 \longrightarrow DB[-1]
 \label{eq:qa-as-dg-triangle}
\end{equation}
がある。
有限代数$B$のSerre関手を用いると
realizationとその右随伴の間に
$S_B[-2]\to\mathrm{id}\to\iota^*\iota_*$を得る。
前補題とこの三角形から
$E_i$のrealizationは全整数で$L_i$となり、
\begin{equation}
\RHom(E_i,E_{j-4})\longrightarrow
\RHom(E_i,E_j)\longrightarrow
\RHom_D(L_i,L_j)
\label{eq:qa-as-rhom-comparison}
\end{equation}
が成立する。
$0<j-i<4$ では $j-4<i$ であり、
$i$ と $j-4$ は同じ窓に入るので最初の項は零である。
従って実現は各窓の前向き Hom 上で合成を保つ同型を与え、
前補題の窓内の計算を例外対象の窓に移せる。
$j-i=4$ の最初の項は対角の $k$ である。
$j-i>4$ では最初の項に距離についての帰納法を使うと、
最後の項も前向き Ext 消滅により次数零に集中するため、
中央の前向き Hom も同じ性質を持つ。
この比較の向きは \cite[式 (9), (30)--(32)]{OkawaUedaSpherical2020}
の $E_{j-4}=S_B(E_j)[-2]$ と一致する。

四項複体を得るため、まず窓 $[i-4,i-1]$ を取る。
前補題の共変 Yoneda 解消の三つの射影項をこの窓に戻すと、
$E_{i-4}$ から始まる逐次右 mutation の複体になる。
実際、三項の 畳み込み は後三対象への導来 Hom が零であり、
$E_{i-4}$ への射の cone は後三対象で生成される。
従って窓の半直交分解によって一意に定まる右 mutation である。
窓の Serre 関手による同定から、その最後の項を $E_i$ とした
四項複体 $K_i$ の全複体は零である。
$K_i$ は次数 $-3,-2,-1,0$ に置き、両端は
$E_{i-4},E_i$ とする。中間項の重複度は前補題の解消から
定義~
ef{def:qa-as-square-rotated-minimal} の rotated 型の表となる。

$RHom(E_j,K_i)$ の項別スペクトル系列で次数零の行を取ると、
右射影複体の頂点 $j$ における評価になる。
$j<i$ では前向き消滅と幅四の窓の半直交性により
他の行は零なので、この行は完全である。
$j=i$ では唯一の例外が左端の
$Ext^2(E_i,E_{i-4})=k$ である。実際、
$S_B(E_i)=E_{i-4}[2]$ なのでこれは Serre 双対から従う。
全複体の消滅により、この $k$ から次数零の行の右端の
コホモロジーへの $d_3$ が同型である。
$j>i$ での評価は directed の定義により零である。
従って右射影複体の唯一のコホモロジーは次数零の $S_i$ である。

双対複体を左頂点 $j$ で評価するには $RHom(K_i,E_j)$ を用いる。
$j>i-4$ では同じ前向き消滅と半直交性により次数零の行が完全であり、
$j<i-4$ では directed の定義により評価は零である。
$j=i-4$ では、次数 $0,1,2,3$ のうち最初の項から来る
$Ext^2(E_i,E_{i-4})=k$ だけが次数零の行の外に残る。
全複体が零なので $d_3$ はこれを次数零の行の
次数三のコホモロジーに同定する。
従って $Ext^3_A(S_i,A)={}_ell S_{i-4}$ であり、
他の次数の Ext は零である。
異なる頂点間の成分しか現れないので右解消は最小である。
\end{proof}


\begin{proposition}[有限Endと全代数の相互逆性]
\label{prop:qa-as-inverse-foundations}
$A\in\mathrm{AS}_{\rm rsq}$から取り出した
$(Y,D,P)$の初期四頂点Endは
\[
A|_{\{0,1,2,3\}}\simeq
\End_{\rm dir}
(P,\OO_D(h-2f),\OO_D(h-f),\OO_D(2h-f))
\]
であり、古典型は$Y$上の可換基準collectionである。
逆にavatarの有限Endから構成した代数の
点表現と余核線束は元の$(Y,D,P)$を回収する。
\end{proposition}
\begin{proof}
後者は命題~\ref{prop:qa-core-directed-foundations}の
二行二列評価行列の局所標準形
$\operatorname{diag}(1,q)$から
$\PP_Y(\coker\kappa)=D$を回収し、
忘れられた線束の正規化で$P$を得ることによる。
前者は、AS代数から作った
\[
0\to\OO_Y^{\oplus2}\xrightarrow{\kappa}
\OO_Y(h)^{\oplus2}\to\iota_*Q\to0
\]
を $\OO_Y(-h)$ で捩ると、$\OO_Y(-h)$ の全コホモロジー消滅により
$H^0(QH^{-1})$ は $0\to2$ の二矢空間と同型になる。
元の列の大域切断は
\[
 0\longrightarrow R_0\longrightarrow
 H^0(Y,\OO_Y(h))\otimes H^0(D,QH^{-1})
 \longrightarrow H^0(D,Q)\longrightarrow0
\]
である。最初の射は普遍評価行列の係数なので、
元の二関係の包含そのものである。
最後の射は可逆層の切断の乗法であり、
$0\to3$ の合成を同定する。

保持した $1,2,3$ の三頂点では、$f,h,h+f$ から $2h$ を引いた
$f-2h,-h,f-h$ のすべてが acyclic なので、
それぞれの切断空間は $Y$ から $D$ への制限で変わらない。
$Y$ 上の $f$ と $h$ の乗法は、
$\PP^1$ 上で $\mathcal E$ の評価列を $\OO(1)$ で捩って得る
全射 $V\otimes H^0(\OO(1))\to H^0(\mathcal E(1))$ であり、
核は $R_1$ である。
従って $1\to3$ の合成も元の二関係で割った積と一致する。
残る
$0\to1$のHomの消滅は
$H^\bullet(\OO_D(h-2f)P^{-1})=0$である。
古典型では $P=\OO_Y$ なので $Q=\OO_Y(2h-f)$ である。
一定階数一の $\kappa$ の像を $I$ とすると、
$0\to I\to\OO_Y(h)^2\to Q\to0$ の行列式から
$I=\OO_Y(f)$、さらに $\ker\kappa=\OO_Y(-f)$ となる。
$\OO_Y^2\twoheadrightarrow I$ の二切断は、
$\OO_Y(f)$ の全生成 pencil の基底をなす。
また $I\hookrightarrow\OO_Y(h)^2$ は、
$\OO_Y(h-f)$ の全生成 pencil の基底を双対化して得る包含である。
いずれも二切断が線形従属なら、正次数の線束を
一切断で生成することになるため不可能である。
この分解で $H^0(\OO_Y^2)\to H^0(I)$ は同型であり、
$H^1(I)=0$ なので、$0\to3$ の積の核は元の $R_0$ となる。
$1\to3$ の積は曲面上の $\mathcal E$ の評価列から同様に $R_1$ を回収する。
従って古典型でも全合成を含めて可換基準 collection を得る。
全$\mathbb Z$-代数は
補題~\ref{lem:qa-as-serre-unique-recovery}が決めるから
有限Endの同型で止まらない。
\end{proof}

\begin{theorem}[squareとrotated squareのavatar分類]
\label{thm:qa-as-avatar-bijection}
定義~\ref{def:qa-as-square-rotated-minimal}の
最小分解型を持つAS正則$\mathbb Z$-代数について
\begin{equation}
\mathrm{AS}_{\rm rsq}/\cong\simeq
 \mathrm{Av}_{\rm rsq}/\cong,\qquad
\mathrm{AS}_{\rm sq}/\cong\simeq
 \mathrm{Av}_{\rm sq}/\cong .
\label{eq:qa-as-avatar-two-bijections}
\end{equation}
$\mathrm{Av}_{\rm rsq}$は
$D\in|2h|$（$\Sigma_2$では$D\cap S=\varnothing$）と
$P\in\Pic^0(D)\setminus\{\OO_D,\OO_D(h-2f)\}$、
$\mathrm{Av}_{\rm sq}$は順序付きrulingを持つ
$D\subset\PP^1\times\PP^1$の$(2,2)$反標準因子と
$P\in\Pic^0(D)\setminus\{\OO_D\}$からなる。
各々に一つの古典型を含める。
曲線の滑らかさ・被約性・既約性は課さない。
逆構成の有限foundationは
\[
B_{\rm rsq}=\End_{\rm dir}(P,HF^{-2},HF^{-1},H^2F^{-1}),
\qquad
B_{\rm sq}=\End_{\rm dir}(P,F,G,FG).
\]
\end{theorem}
\begin{proof}
rotated型では命題~\ref{prop:qa-as-extraction}が必要性、
命題~\ref{prop:qa-as-helix-reconstruction}が
任意avatarのAS実現、
命題~\ref{prop:qa-as-inverse-foundations}と
補題~\ref{lem:qa-as-serre-unique-recovery}が
両構成の相互逆性を与える。
代数同型は点表現・普遍線束を保ち、
avatar同型はHomと合成を保つので、
頂点を保つ代数同型類の全単射となる。

square型代数を一つ進めた
$A^{[1]}_{ji}=A_{j+1,i+1}$は
表からrotated型であり逆shiftも存在する。
一方定理~\ref{thm:qa-core-avatar-square-rotated}は
二つのavatarの同型類を相互変換する。
rotatedの一つ前の四項を曲線上の共通線束
$FP^{-1}$で正規化すると
$(L_{-1},L_0,L_1,L_2)\otimes FP^{-1}
=(P,F,G,FG)$となる。
これは曲線上の操作であって、
非可換圏上にtensor自己同値を仮定しない。
全代数のSerre復元もshiftと可換なので
square側の全単射が従う。
とくに$P=FG^{-1}$や$P=GF^{-1}$を
squareの許容条件から排除してはならない。
この共鳴はrotatedの$\Sigma_2$への移行を
記録し得るがAS正則性を破らない。
\end{proof}

\begin{proposition}[左右Noether性]
\label{prop:qa-as-avatar-noetherian}
上記二型の全代数は左右Noetherianである。
曲線型には四次の正規正則family
$g_i\in A_{i+4,i}$があり、その商は
$D$のtwisted section algebraを与える。
\end{proposition}
\begin{proof}
式~\eqref{eq:qa-as-rhom-comparison}の次数零から
$j>i$について
\[
0\to A_{j-4,i}\xrightarrow{\,g\,}A_{ji}
\to H^0(D,L_jL_i^{-1})\to0
\]
を得る。$t=S_B^{-1}[2]$ とし、三角形の最初の自然変換を
$\eta:t^{-1}\to\operatorname{id}$ と書く。
$E_i\to E_j$ の射 $a$ に対して自然性は
\[
 a\,\eta_{E_i}=\eta_{E_j}\,t^{-1}(a)
\]
を与える。$t^{-1}$ は各 Hom で同型なので、
$\eta$ を前から掛けた像と後ろから掛けた像は一致する。
上の短完全列から後合成による乗法は単射であり、
この式を使えば前合成による乗法も単射となる。
また自然性を $\eta_{E_j}$ 自身に適用すると
$\eta_{E_j}t^{-1}(\eta_{E_j})=
\eta_{E_j}\eta_{E_{j-4}}$ となる。
後合成の単射性により $t^{-1}(\eta_{E_j})=\eta_{E_{j-4}}$
である。従って $g_i=\eta_{E_{i+4}}$ は四周期の
折り畳みと整合する両側正規正則の族となる。
偏極の移動公式に加えて、正規化した初期窓では
$P=(L_0L_1^{-1})HF^{-2}$ である。
この式を次窓に適用すると
$\sigma^*P=(L_4L_5^{-1})(\sigma^*H)(\sigma^*F)^{-2}=P$
を得る。したがって周期式~\eqref{eq:qa-core-rsq-periodic-L}に
$\sigma^*F=FP^2$、$\sigma^*H=HP^4$、$\sigma^*P=P$ を代入すれば、
全整数で $L_{i+4}=\sigma^*L_i\otimes\OO_D(2h)$ となる。
$2h$ は $\Sigma_2$ でも $D\cap S=\varnothing$ により
各成分上正次数なので $D$ 上 ample である。
曲線上の数値類は各既約成分の次数で測られ、
$\sigma$ はこれらを有限置換する。
従って数値類への作用は準冪単であり、
\cite{Keeler2003,Keeler2018Corrigendum} の判定で
この偏極は $\sigma$-ample かつ $\sigma^{-1}$-ample となる。
twisted section algebra は左右 Noetherian である。
四周期をまとめた代数を $J$ とし、$g$ を四位相における
$g_i$ の和とする。$J/gJ$ の対角四隅は
$\OO_D(2h)$ のねじれ切断環である。
異なる位相の隅は、それぞれの対角隅上で
可逆層の切断加群となる。両方向の ampleness により
これらは左右で有限生成なので、$J/gJ$ は対角環の直積上で
左右有限加群となり、左右 Noetherian である。

正規正則性により、$g$-adic 随伴次数環は
$J/gJ$ 上の skew polynomial 環となる。
実際、各 $g^mJ/g^{m+1}J$ は $J/gJ$ と同定でき、
乗法は $g$ を移す正規自己同型でねじれる。
従ってこの随伴次数環は左右 Noetherian である。
内部次数で斉次な右 ideal について、その $g$-adic 初項の
ideal の有限斉次生成系を持ち上げる。
元の斉次元から初項を消すたびに $g$-adic 位数は増え、
$g$ の内部次数が正なので、その元の次数を固定すれば
有限回で停止する。よってこれらの持ち上げは斉次 ideal を生成する。
一般の右 ideal には内部次数の最高次項を使い、
次数を下げる帰納法で有限生成性を得る。左 ideal も同様である。
最後に $A$ の各射影行・列の斉次部分加群は、
折り畳み後に位相の冪等元を保つ斉次部分加群として扱える。
斉次生成元を元の添字へ戻せるので、$A$ も左右 Noetherian である。
古典型は通常の可換曲面のNoether性に帰着する。
\end{proof}

\begin{proposition}[Noether性からcoherentな導来同値へ]
\label{prop:qa-as-avatar-coherent-morita}
前定理のsquare／rotated square型の代数$A$について
$X=\Qgr A$、$\coh X=\qgr A$と置き、
$T=E_0\oplus E_1\oplus E_2\oplus E_3$、
$B=\End(T)$とする。
すると
\[
 \RHom_X(T,-):
 D^b(\coh X)\xrightarrow{\ \sim\ }D^b(\mathrm{mod}\,B)
\]
は三角圏同値である。
さらに$\coh X$の完全自己同値$\nu$があり、
$\nu E_i\simeq E_{i-4}$、
$D^b(\coh X)$のSerre関手は
$\nu[2]$と同型になる。
これは\emph{導来}Morita同値であり、
$\coh X\simeq\mathrm{mod}\,B$という
Abel圏同値を主張しない。
\end{proposition}
\begin{proof}
命題~\ref{prop:qa-as-avatar-noetherian}より
$A$は左右Noetherianである。
有限生成graded加群の有限表示のsyzygyに、
左右の最小単純分解と$\operatorname{Tor}$の検出を繰り返すと、
射影分解は長さ三以内で停止する。
実際、さらに非零のsyzygyがあれば、
左単純加群との$\operatorname{Tor}_{>3}$に非零の項を持つが、
左単純の長さ三の射影分解に矛盾する。
従って各連接対象は有限個の$E_i$のconeから生成される。
補題~\ref{lem:qa-as-serre-unique-recovery}の
四周期完全列により、これらの$E_i$は
$T$が生成する厚部分圏に入る。
$T$はstrong exceptional tiltingであり、
そのEndは有限大域次元二の$B$である。
これで導来Morita同値が$\coh X$の
有界導来圏全体へ延長される。

四周期のSerre同型
$E_{i+4}\simeq S_B^{-1}(E_i)[2]$を
全整数で自然化した$\mathbb Z$-代数のtransportは、
graded加群のexactな四頂点ずらしと可換で、
torsionを保つ。
従って商圏の完全自己同値$\nu$を与え、
$\nu E_i=E_{i-4}$となる。
有限tilting generatorとその全Endにおいて
$\nu[2]$は$B$のSerre関手と一致する。
生成性からSerre関手として同型となる。
\end{proof}

\begin{remark}[一般定理への依存と未証明の比較]
\label{rem:qa-as-classification-scope}
本節では元PDFの体上分類の
必要性・逆構成・全代数の相互逆性を移管した。
spherical helix実現の一般定理
\cite[Theorem 5]{OkawaUedaSpherical2020}と
twisted section algebraのNoether性は外部入力である。
これはAS代数の同型類の対応であり、
別に定義された内在的O-marked$\HHstack$との
pointed同値や任意の非被約基底上のAS相対分類を
意味しない。
\end{remark}

% !TeX root = ../nc-quadric-avatar.tex
% 2026-11-22 / minimal ten-arrow AS type, mutation and twelve-arrow status
\section{最小十本矢型の独立AS分類と十二本矢表示}
\label{sec:qa-as10-mutation-classification}

前節のsquare AS代数から一頂点mutationを行うと、
重み四relation potentialのreductionにより
八本矢cubic型または最小十本矢型が現れる。
ここで矢の本数だけからAS正則性を結論してはならない。
元草稿\cite[第VIII部，第X部]{AvatarVSeven2026}の
Ginzburg次数零集中・独立逆構成・評価/余評価の証明を
本稿のavatar分類に接続する。

\subsection{Mutation像によらない最小十本矢AS条件}
\label{subsec:qa-as10-minimal}

\begin{definition}[最小十本矢AS型]
\label{def:qa-as10-minimal}
$i$が奇数なら$\varepsilon_i=1$、
偶数なら$\varepsilon_i=0$とする。
positively directedで局所有限、
$A_{ii}=k$の$\mathbb Z$-代数が
$\mathrm{AS}_{10}^{\mathrm{min}}$に属するとは、
全$i$の右単純$S_i$に対して
\begin{equation}
0\to P_{i-4}\to
 P_{i-3}^{\oplus2}\oplus P_{i-2}^{\oplus\varepsilon_i}
 \to P_{i-1}^{\oplus2}\oplus P_{i-2}^{\oplus\varepsilon_i}
 \to P_i\to S_i\to0
\label{eq:qa-as10-minimal-resolution}
\end{equation}
が最小であり、さらに
$\Ext_A^q(S_i,A)=0$ ($q\ne3$)、
$\Ext_A^3(S_i,A)\simeq {}_\ell S_{i-4}$
を満たすことをいう。
同型は全整数頂点を保つ。
QP、mutation元、周期同型と分解の基底は
定義の標識ではない。
\end{definition}

\begin{lemma}[Hilbert関数と有限quiver]
\label{lem:qa-as10-hilbert-foundation}
$A\in\mathrm{AS}_{10}^{\mathrm{min}}$なら
\[
\dim A_{i+2r,i}=(r+1)^2,\qquad
\dim A_{i+2r+1,i}=(r+1)(r+2).
\]
四頂点のfinite directed algebraは大域次元二で、
最小quiverに
$0\to1\to2\to3$の各二矢と
追加の$x:1\to3$を持ち、
最小関係は$1\to3$の一二次関係と
$0\to3$の二三次関係である。
この有限foundationは全$\mathbb Z$-代数を
Serre関手で一意に回収する。
\end{lemma}
\begin{proof}
\eqref{eq:qa-as10-minimal-resolution}の
二つの$P_{i-2}$項はEuler漸化式で相殺され、
$h_n-2h_{n-1}+2h_{n-3}-h_{n-4}=0$、
$h_0=1$、$h_n=0$ ($n<0$)を得る。
生成関数は
$1/((1-t)^2(1-t^2))$であり、表示の次元が従う。
最小分解の前二段を四頂点へ制限すると、
列挙した矢と三つの関係が現れ、
関係同士の次のoverlapはないため
大域次元二である。
左右AS双対性による局所コホモロジーの証明は
補題~\ref{lem:qa-as-serre-unique-recovery}と同じで、
連続四項のSerre関手による
$E_{i+4}=S_B^{-1}(E_i)[2]$が全代数を回収する。
\end{proof}

\subsection{Ginzburg次数零集中によるAS条件}
\label{subsec:qa-as10-ginzburg}

\begin{lemma}[relation QPの次数零集中]
\label{lem:qa-as10-ginzburg-concentration}
square型、cubic型または
$\mathrm{AS}_{10}^{\mathrm{min}}$の有限foundationから
最小関係の双対戻り矢で作るrelation QPを$(Q,W)$とする。
非完備およびarrow-adic完備Ginzburg dg代数について
\[
 H^j\Gamma(Q,W)=H^j\widehat\Gamma(Q,W)=0
 \qquad(j\ne0)
\]
が成立する。
\end{lemma}
\begin{proof}
有限foundationの大域次元は二なので、
未変形の3-Calabi--Yau completion
$\Pi_3(B)=T_B\RHom_{B^e}(B,B^e)[2]$を考える。
\cite[Theorem 6.10]{KellerDeformedCY2011}の
未変形の場合はrelation Ginzburg dg代数と
$\Pi_3(B)$の準同型を与える。
その$m$番目のtensor成分は
$S_B^{-m}[2m]$の全周期Homとして計算できる。
$m=0$では有限foundationのtilting性、
$m\ge1$では
前節~\ref{lem:qa-as-serre-unique-recovery}または
前補題の前向きExt消滅によって
次数零だけに集中する。
従って直和も次数零集中を満たす。
ここで用いたのは未変形completionであり、
任意のdeformation parameterのCY性ではない。

元の矢の内部次数は$1,2$、
反対矢は$4-d$、頂点loopは$4$である。
全dg生成元の内部次数は正で、
各内部次数成分は有限次元複体。
完備化は内部次数の直積であり、
体上の直積は完全なので
各次数の集中は完備版でも保たれる。
逆向きも各内部次数成分を取れば成立する。
\end{proof}

\begin{lemma}[mutation後の最小AS分解]
\label{lem:qa-as10-keller-yang}
前補題のQPを、loopとoriented $2$-cycleを通らない
指定のmutation頂点で変換すると、
reduced QPのGinzburg dg代数も次数零に集中する。
その正次数付きJacobian algebraの頂点単純加群は
長さ三の完全な最小自己双対分解を持ち、
AS--Gorenstein条件を満たす。
\end{lemma}
\begin{proof}
\cite[Theorem 6.2]{KellerYang2011}を用いる。
必要なのはGinzburgの次数零集中、
mutation頂点にloopがないこと、および
局所的な2-cycle除外条件である。
squareの指定中間頂点にloop・2-cycleはなく、
最小十本矢型の2-cycleは別の二頂点間だけにあり、
逆mutation頂点を通らない。
前補題が第一の入力を保証する。
mutation後の完備Ginzburgの次数零集中が得られ、
正の内部次数付けで非完備版へ降りる。
Ginzburgの頂点単純の標準semifree解消を
$H^0\Gamma$に移すと、
恒等射、入る矢、最小関係、四次最上項からなる
Jacobian complexを得る。
次数零集中によりこれは\emph{実際に完全}である。
双対複体には最上次数の左単純だけが残る。
reduced QPでは微分はすべて正次数だから
最小性も成り立つ。
単なるdg導来同値から通常代数のAS正則性を
推論したのではない。
\end{proof}


% !TeX root = ../nc-quadric-avatar.tex
% 2026-11-22 / cubic reachability and derived obstruction
\section{二重壁でのcubic型への到達不能性}
\label{sec:qa-as-cubic-obstruction}

二つの乗法がともに階数三となることだけでは、
他のmutation列や別のexceptional foundationによる
cubic型への到達を排除できない。
本節ではこの強い否定方向を
\cite[第19--21節]{AvatarVSeven2026}から
必要な数値・層論・平滑化の証明を保って統合する。
分類対象は引き続き全整数頂点を保つAS代数、
比較対象は構造対象を指定しない商Abel圏・導来圏である。

\subsection{有限foundationでのcubic候補の条件}
\label{subsec:qa-as-cubic-candidate}

$A\in\mathrm{AS}_{\mathrm{sq}}$の有限foundationを
$B=\End(E_0,E_1,E_2,E_3)$、
$\mathcal D=\operatorname{per}B\simeq D^b(\qgr A)$とする。
そのEuler行列は
\[
G_{\mathrm{sq}}=
\begin{pmatrix}
1&2&2&4\\0&1&0&2\\0&0&1&2\\0&0&0&1
\end{pmatrix}.
\]
\begin{definition}[cubic候補foundation]
\label{def:qa-as-cubic-candidate}
$\mathcal D$の四対象$(Q_0,Q_1,Q_2,Q_3)$が
cubic候補foundationであるとは、
strong exceptionalであり、
その$K_0(\mathcal D)$の類が整数基底をなし、
Euler行列が
\[
G_c=
\begin{pmatrix}
1&2&4&6\\0&1&2&4\\0&0&1&2\\0&0&0&1
\end{pmatrix}
\]
で、二つの隣接合成
\begin{equation}
\Hom(Q_{i+1},Q_{i+2})\otimes\Hom(Q_i,Q_{i+1})
\xrightarrow{\ \sim\ }\Hom(Q_i,Q_{i+2})
\quad(i=0,1)
\label{eq:qa-as-cubic-two-compositions}
\end{equation}
が同型であることをいう。
fullnessは別途仮定しない。
任意のcubic AS代数の標準四対象はこの条件を満たす。
\end{definition}

\begin{lemma}[cubic候補のrankと数値線束型]
\label{lem:qa-as-cubic-candidate-numerical}
cubic候補$(Q_0,\ldots,Q_3)$が存在すれば
共通の$[1]$ shiftの後、全$Q_i$の数値rankは一、
かつ適当な$a,b\in\mathbb Z$に対して
\[
c_1(Q_i)=(a,b),\ (a+1,b),\
(a+1,b+1),\ (a+2,b+1)
\]
または二つのrulingを交換した順序を取る。
数値Chern座標を$(r,a,b,s)$と書けば
各$Q_i$は$s=ab$を満たす。
\end{lemma}
\begin{proof}
squareの標準基底に対応する数値座標を
\[
[E_0]=(1,0,0,0),\quad [E_1]=(1,1,0,0),\quad
[E_2]=(1,0,1,0),\quad [E_3]=(1,1,1,1)
\]
とする。$x^2=y^2=0,\ xy=1$の形式的Chern記号で
Serre作用は
\[
S_*(r,a,b,s)=(r,a-2r,b-2r,s-2a-2b+4r)
\]
であり、
\[
(S_*-1)^2v=8r(v)\delta,\qquad
\delta=(0,0,0,1).
\]
cubicのEuler行列$G_c$から計算した
同じ作用の二乗は、四つの候補基底ベクトルを
同一のprimitiveな交代結合
$[Q_0]-[Q_1]-[Q_2]+[Q_3]$へ送る。
従って各rankは共通整数$r$であり、
この結合が$K_0$のprimitive元なので$r=\pm1$。
共通shiftで$r=1$とする。
自己Euler値は$1+2s-2ab=1$なので$s=ab$。
二つのrank一数値類のEuler値は
$(1+a_j-a_i)(1+b_j-b_i)$。
隣接方向のEuler値は二、逆向きは零だから、
差は$(1,0)$または$(0,1)$。
二段隔てたEuler値四により差は$(1,1)$となるので、
隣接差は交互となり主張を得る。
\end{proof}

\begin{lemma}[曲線への制限とdeterminant]
\label{lem:qa-as-cubic-restriction-determinant}
$(D,F,G,P)$をsquare avatarとし、
$\rho:\operatorname{per}B\to\Perf(D)$を
四線束$(P,F,G,FG)$へのrealizationとする。
任意の$X,Y\in\mathcal D$に自然な三角形
\begin{equation}
\RHom(Y,X)^\vee[-2]\longrightarrow
\RHom(X,Y)\longrightarrow
\RHom_D(\rho X,\rho Y)
\label{eq:qa-as-cubic-spherical-restriction}
\end{equation}
がある。
数値類$[X]=(r,a,b,s)$に対し
\begin{equation}
\det(\rho X)=F^aG^bP^{\,r-a-b+s}.
\label{eq:qa-as-cubic-determinant}
\end{equation}
$D$が滑らかで$X$がrank一のexceptional対象なら
$\rho X$は次数零のshiftを除き線束である。
\end{lemma}
\begin{proof}
前節のdg双加群三角形
\eqref{eq:qa-as-dg-triangle}とSerre双対性を
$\RHom(X,-)$へ適用すると
\eqref{eq:qa-as-cubic-spherical-restriction}を得る。
$X=Y$がexceptionalなら、
右辺の自己$\RHom$は$k\oplus k[-1]$となる。
$D$が滑らかなら$\coh D$は遺伝的で、
有界複体はコホモロジー層のshiftの直和になる。
Endが$k$でrankが一だから、
制限は一つのrank-one indecomposable vector bundleの
shift、すなわち線束のshiftである。
（torsionとvector bundleの拡大は分裂する。）
rank一の符号を共通shiftで正に揃える。

数値基底の座標を解くと
\[
[X]=(r-a-b+s)[E_0]+(a-s)[E_1]
 +(b-s)[E_2]+s[E_3].
\]
perfect複体のdeterminantの加法性と
$\rho(E_0,E_1,E_2,E_3)=(P,F,G,FG)$から
\eqref{eq:qa-as-cubic-determinant}を得る。
\end{proof}

\begin{proposition}[滑らかな二重壁では任意のcubic候補が不存在]
\label{prop:qa-as-cubic-smooth-no-candidate}
$D$を滑らかな楕円曲線とし、
$P^{\otimes2}\simeq\OO_D$、
$P\not\simeq\OO_D$、$G\simeq FP$とする。
$\mathcal D=\operatorname{per}B(D,F,G,P)$には
cubic候補foundationは存在しない。
\end{proposition}
\begin{proof}
候補を仮定する。
補題~\ref{lem:qa-as-cubic-candidate-numerical}により
全rankを一とし、補題
~\ref{lem:qa-as-cubic-restriction-determinant}により
$\rho Q_i=N_i[2m_i]$と書く。
隣接する$N_i$の次数差は二で、
$\Ext$の三角形と正次数線束の$H^1=0$から
前向きHomが次数零に集中するため、
各$m_i$は共通である。
共通shiftを除き、$N_i$を四線束とする。

\eqref{eq:qa-as-cubic-determinant}で
$G=FP$、$s=ab$、$P^2=\OO_D$を用いると
$\det(\rho Q)=F^{a+b}P^{1-a+ab}$となる。
数値差が$(1,0),(0,1),(1,0)$の場合、
連続する三つの商線束は
\[
(N_1N_0^{-1},N_2N_1^{-1},N_3N_2^{-1})
=(FP^{b-1},FP^{a+1},FP^b).
\]
逆順のrulingなら
$(FP^a,FP^b,FP^{a+1})$となる。
$P$が位数二なので商は$F,FP$の二種類だけで、
最初と最後が異なる。
従って三つのうち隣接する二つは同型である。
その共通次数二線束を$L$とすると、
合成写像
\[
 H^0(D,L)\otimes H^0(D,L)\longrightarrow H^0(D,L^2)
\]
のkernelは
$\bigwedge^2H^0(D,L)$という一次元空間である。
他方\eqref{eq:qa-as-cubic-spherical-restriction}によって
候補のforward Homと合成は
この線束上のHomと合成に一致する。
これは\eqref{eq:qa-as-cubic-two-compositions}の
同型性に矛盾する。
特定のmutation列に依存せず、
全てのcubic候補を排除した。
\end{proof}

% !TeX root = ../nc-quadric-surface-avatar.tex
\section{反標準 avatar の族と二つの十本矢 chart}
\label{sec:qa-avatar-charts}

本節では幾何的な族を直接定義し、その間の開埋め込みを証明する。
代数の同型類から幾何的 avatar への対応は
\cite[定義15.2・定理15.3]{Q12Avatar2026}に記載されているが、
その同型類分類を相対 AS 代数の stack 同値として使用しない。
族上の再埋め込みは\cite[第4--5節]{Q12Weyl2026}との比較も踏まえ、
本稿の補題~\ref{lem:qa-core-curve-vanishing}、
\ref{lem:qa-core-two-pencils}と定理~\ref{thm:qa-core-avatar-square-rotated}で
必要な幾何を自足的に証明し、
開部分の判定に必要なコホモロジー、非被約因子、垂直成分および降下を以下で確認する。
統合ノート\cite[第7.11--7.12項・第9節・第13.3項]{Synthesis2026}は
これらの出典と主張の水準を整理した資料であり、独立の検証の代用とはしない。

\paragraph{AS 側の辞書の対象.}
入力する体上分類の対象は
\cite[第15.1項・定理15.3]{Q12Avatar2026} の $A_Q^{\mathrm{pres}}$ である。
すなわち、局所有限な連結 positively directed $\mathbb Z$-代数で、
大域次元三、Gorenstein parameter 四の AS 正則性を満たし、
\eqref{eq:qa-universal-potential} の十二本矢 potential
（可縮な二次 pair を許す）の展開として表示できるものを、
全整数頂点を固定する代数同型で分類する。
最小化した四型では、右射影 $\mathsf P_i=e_iA$ と右単純 $\mathsf S_i$ に対して
\[
\begin{aligned}
0\longrightarrow\mathsf P_{i-4}
 &\longrightarrow\mathsf P_{i-3}^{\oplus2}\oplus\mathsf P_{i-2}^{\oplus\epsilon_i}\\
 &\longrightarrow\mathsf P_{i-1}^{\oplus2}\oplus\mathsf P_{i-2}^{\oplus\epsilon_i}
 \longrightarrow\mathsf P_i\longrightarrow\mathsf S_i\longrightarrow0
\end{aligned}
\]
という実際に完全な最小分解がある。
$\epsilon_i=0$、奇数 $i$ のみで $\epsilon_i=1$、
偶数 $i$ のみで $\epsilon_i=1$、全 $i$ で $\epsilon_i=1$ が、
順に八本、奇数対角十本、偶数対角十本、十二本に対応する。
右側の AS--Gorenstein 条件は
$\Ext_A^q(\mathsf S_i,\mathsf P_j)=0$ for $(q,j)\ne(3,i-4)$、
$\Ext_A^3(\mathsf S_i,\mathsf P_{i-4})\simeq\kk$ である。
分解の Euler 特性だけを課す条件ではない。
Noether 性はこの定義の追加仮定ではなく、引用する分類の結論である。
potential 表示、trivial pair の係数、周期同型、分解の基底は選択データにしない
\cite[定義2.1・第15.1項]{Q12Avatar2026}\cite[定義48.1]{AvatarVSeven2026}。
この整数頂点固定の同型は、四頂点 graded 代数の同型を射とした
$\mathfrak M_r^{\mathrm{alg}}$ と同じ相対同型層であるとは仮定しない。
本節の幾何的 stack の同型との比較も、体上の同型類辞書と区別する。

\subsection{分類する幾何データ}

基礎体 $\kk$ は標数零の代数閉体とする。
射影束は一次元商を分類する規約を用いる。
次の二つの標識付き曲面を考える。
\[
\begin{array}{c|c|c}
Y&h&f\\\hline
\Sigma _0=\PP^1\times\PP^1&\OO(1,1)&\OO(1,0)\\
\Sigma _2=\PP_{\PP^1}(\OO\oplus\OO(2))&\OO_{\PP(\OO\oplus\OO(2))}(1)&\pi^*\OO_{\PP^1}(1)
\end{array}
\]
いずれも $h^2=2$, $hf=1$, $f^2=0$, $-K_Y=2h$ である。
$\Sigma _2$ の負断面は $S=h-2f$ である。
本節の $\Pic^0$ は Picard 群の単位成分を表す。

\begin{definition}[曲線型の相対 avatar]
\label{def:qa-avatar}
局所 Noetherian な $\kk$-scheme $T$ 上の
$\mathcal A_{12}^{\times}$ の対象を
\[
 (Y\xrightarrow{\pi}B\longrightarrow T;\ h,f;\ D\hookrightarrow Y;\ P)
\]
とする。ここで $B/T$ は $\PP^1$-fibration、$Y/B$ は
$T$ 上 fppf 局所で $\PP_B(\mathcal E)$ と表示される射影直線族であり、
各幾何 fiber の $(Y,h,f)$ は上表のいずれかである。
$h,f$ は相対 Picard の類で、$f$ は指定した ruling に対応する。
$D$ は $T$ 上平坦な相対有効 Cartier 因子で、相対因子類は $2h$ である。
$P$ は $\Pic^0_{D/T}$ の section で、各幾何 fiber 上で非自明とする。
射は ruling、標識、$D$ および $P$ の類を保つ $T$ 上の同型である。
線束の scalar trivialization はデータに加えない。
\end{definition}

この定義は $\Sigma _0$ と $\Sigma _2$ の間で型が跳ぶ族を許す。
実際の線束を用いる式は fppf 局所で解釈し、基底からの線束の tensor は区別しない。
以下で
\[
 p=\pi|_D:D\longrightarrow B,\qquad F=\OO_Y(f)|_D,\qquad H=\OO_Y(h)|_D
\]
と書く。特に $H$ は線束であり、内在的モジュライの記号とは無関係である。
曲面、因子および相対 Picard 類の降下により、定義\ref{def:qa-avatar}は
fppf stack を定める。本節の「開部分 stack」は任意の基底への引戻しが
開部分 scheme となる意味で用いる。

\begin{remark}[有限表示と scalar の扱い]
fppf 局所で $B=\PP^1_T$ とし、$\mathcal E=\pi_*\OO_Y(h)$ を取ると、
各幾何 fiber の分裂型は $\OO(1)^{\oplus2}$ または $\OO\oplus\OO(2)$ である。
評価列は
\[
 0\longrightarrow \mathcal R\otimes\OO_B(-1)
 \longrightarrow \mathcal V\otimes\OO_B
 \longrightarrow\mathcal E\longrightarrow0,
 \qquad\rank\mathcal R=2,\quad\rank\mathcal V=4
\]
となる。したがって線形行列、その余核の射影束上の相対反標準因子、
および相対 Picard 類によって局所的な parameter を記述できる。
この表示の基底変更をすべて新たな自己同型として保持すると、
定義\ref{def:qa-avatar}と異なる stack になる。
本節の射はあくまで標識付き曲面と因子の同型である。
\end{remark}

\subsection{非被約因子を含むコホモロジー}

\begin{lemma}[支持に対する判定]
\label{lem:qa-support}
体上で $D\in|2h|$ とする。このとき
$\omega_D\simeq\OO_D$、$h^0(\OO_D)=h^1(\OO_D)=1$ である。
真の非零部分因子 $0<Z<D$ について $Z(D-Z)\geq2$ が成り立つ。
線束 $N$ の非零切断が正確に $Z\leq D$ に支持されるなら
\[
 \deg_ZN\ \geq\ Z(D-Z).
\]
従って、線束 $L$ がすべての $0<Z\leq D$ について
\begin{equation}
 \deg_ZL+Z(D-Z)>0
 \label{eq:qa-support-vanishing}
\end{equation}
を満たせば $H^1(D,L)=0$ である。
\end{lemma}

\begin{proof}
随伴と $0\to\omega_Y\to\OO_Y\to\OO_D\to0$ により
双対化層と構造層のコホモロジーを得る。
$\Sigma _0$ では $D=2f+2g$、$Z=af+bg$ と書けて
\[
 Z(D-Z)=2(a+b-ab),\qquad 0\leq a,b\leq2.
\]
$Z=0,D$ を除いたすべての場合で右辺は二以上である。
$\Sigma _2$ では $D=2S+4f$、$Z=aS+bf$ と書くと
\[
 Z(D-Z)=2a^2+2b(1-a),\qquad0\leq a\leq2,\quad0\leq b\leq4.
\]
$a=0$ では $b\geq1$、$a=1$ では右辺が二、
$a=2$ では真部分因子なら $b\leq3$ なので、同じ結論を得る。

切断の generic annihilator によって、その正確な支持である部分因子 $Z$ を取る。
各既約成分の generic な Artin 局所環で切断の消滅次数を取り除いて $Z$ の重複度を定めると、
切断は $N|_Z(-(D-Z))$ の、各 generic point で可逆な切断になる。
$Z$ は Cohen--Macaulay なので、この切断が定める
$\OO_Z\to N|_Z(-(D-Z))$ は単射で、余核は有限長である。
Euler 特性を比較すれば不等式を得る。
これは部分因子を重複度込みで取る議論であり、$D$ の被約性を使わない。
最後は $\omega_D\simeq\OO_D$ による Serre 双対を用い、
$L^{-1}$ の切断にこの不等式を適用すればよい。
\end{proof}

\begin{lemma}[再埋め込みに必要な消滅]
\label{lem:qa-vanishing}
$R\in\Pic^0(D)$ とする。このとき
\[
 H^1(D,HR)=H^1(D,HF^{-1}R)=0,\qquad
 \chi(HR)=4,\quad\chi(HF^{-1}R)=2.
\]
また $R\not\simeq\OO_D$ なら $H^*(D,R)=0$ である。
\end{lemma}

\begin{proof}
$R$ は各部分因子上で次数零である。
$h$ は nef で全次数が四なので、補題\ref{lem:qa-support}を適用できる。
$\Sigma _0$ では $h-f$ も nef で全次数が二である。
$\Sigma _2$ では $Z=aS+bf$ 上で
\[
 \deg_Z(h-f)+Z(D-Z)=2a^2-a+3b-2ab.
\]
$a=0,1,2$ に対する右辺はそれぞれ $3b$, $1+b$, $6-b$ であり、
$Z\neq0$ の許される範囲で正である。Riemann--Roch が Euler 特性を与える。
$R$ の非零切断は真部分因子には支持されない。
全体に支持される切断の余核は次数零なので零であり、$R$ は自明になる。
$R^{-1}$ にも同じ議論を適用し、Serre 双対を使う。
\end{proof}

\subsection{Picard twist による相対的な再埋め込み}

\begin{proposition}[垂直成分がある場合も含む直像]
\label{prop:qa-pushforward}
定義\ref{def:qa-avatar}の族で、$R\in\Pic^0_{D/T}(T)$ を取る。
fppf 局所で線束によって表示すると
\[
 Rp_*(HR)\simeq\mathcal E_R:=p_*(HR)
\]
は $B$ 上の階数二の局所自由層であり、任意の基底変換と両立する。
各幾何 fiber $B_t\simeq\PP^1$ 上で次数二かつ大域生成である。
評価写像は閉埋め込み
\[
 j_R:D\hookrightarrow Y_R:=\PP_B(\mathcal E_R)
\]
を定め、像は相対反標準 Cartier 因子である。
新しい標識の制限は $F_R=F$、$H_R=HR$ である。
\end{proposition}

\begin{proof}
まず体上で $b\in B$、$G_b=\pi^{-1}(b)$ とする。
$G_b$ が $D$ の成分でなければ $D\cap G_b$ は長さ二の scheme であり、
導来 fiber の切断は二次元で次数零に集中する。
$G_b\leq D$ なら、局所 parameter $t$ と $D$ の方程式 $t^m a$、$t\nmid a$ を用いると
\[
 \operatorname{ann}_{\OO_Y/(t^ma)}(t)=(t^{m-1}a)/(t^ma).
\]
したがって $L=HR$ の $B$ 方向の導来制限は
\[
 \mathcal H^{-1}=L|_{G_b}\otimes\OO_{G_b}(-D+G_b)\simeq\OO_{\PP^1}(-1),
 \qquad
 \mathcal H^0=L|_{G_b}\simeq\OO_{\PP^1}(1)
\]
を持つ。ここで $\deg_{G_b}L=1$、$D G_b=2$、$G_b^2=0$ を使った。
次数 $-1$ の層は acyclic であるから、この場合も直像の導来 fiber は $\kk(b)^2[0]$ となる。
垂直成分を持つ $p$ は有限でも平坦でもないが、この計算にはその仮定を要しない。

$i_*L$ は Cartier immersion $i:D\hookrightarrow Y$ の線束の直像なので perfect である。
滑らかな射影的射 $\pi$ による perfect pushforward と導来基底変換を用いると、
$R\pi_*i_*L$ は perfect である。
すべての導来 fiber が同じ階数二の次数零複体なので、局所的に可逆 minor で
可縮項を消去することにより、階数二の局所自由層となる。
族の場合、$i_*L$ は $T$ 上平坦であり、まず $T$ の幾何 fiber、次に $B_t$ の一点へ
制限すれば同じ計算になる。このため非被約な基底も含めて結論が成り立つ。

補題\ref{lem:qa-vanishing}と Riemann--Roch により
\[
 \deg\mathcal E_R=\chi(HR)-2=2,\qquad
 H^1(B,\mathcal E_R(-1))=H^1(D,HF^{-1}R)=0.
\]
$\PP^1$ 上の分裂定理より、分裂型は $\OO(1)^{\oplus2}$ または $\OO\oplus\OO(2)$ である。

各 $B$-fiber での評価は、長さ二の scheme 上の完全切断の評価、
または $\OO_{\PP^1}(1)$ の評価である。従って全射で、誘導される fiber の射は閉埋め込みである。
Nakayama の補題から全評価も全射となる。$j_R$ は proper かつ quasi-finite なので有限であり、
有限射の基底変換と再度の Nakayama により
$\OO_{Y_R}\to(j_R)_*\OO_D$ は全射である。
体上では像は滑らかな曲面内の Cohen--Macaulay な余次元一部分 scheme なので Cartier である。
像の類を $a h_R+b f_R$ とすると、ruling との交点から $a=2$、
$h_R$ との交点 $\deg_D(HR)=4$ から $b=0$ を得る。
族では $D/T$ の平坦性と各 fiber の Cartier 性によって相対 Cartier 性が従う。
各 fiber で自明な差の線束は、$H^1(\OO_{Y_t})=0$ と基底変換により基底から来るので、
相対 Picard 類も $2h_R$ である。
\end{proof}

\begin{theorem}[族上の幾何的並進]
\label{thm:qa-translation}
前命題において $R=P^n$ とした構成は $\mathcal A_{12}^{\times}$ の自己同値
$\mathsf T^n$ を定める。自然な同型
\[
 \mathsf T^m\mathsf T^n\simeq\mathsf T^{m+n},\qquad
 \mathsf T^0\simeq\operatorname{id}
\]
があり、逆操作は $P^{-n}$ を用いる同じ構成である。
\end{theorem}

\begin{proof}
抽象的な $D$、ruling $p:D\to B$ および $P$ は変更せず、$H$ を $HP^n$ に変える。
二回適用した直像は $p_*(HP^nP^m)$ なので、tensor の結合則が合成則を与える。
$R=\OO_D$ の場合、制限列
\[
 0\longrightarrow\OO_Y(-h)\longrightarrow\OO_Y(h)\longrightarrow i_*H\longrightarrow0
\]
と $R\pi_*\OO_Y(-h)=0$ が元の $\mathcal E=\pi_*\OO_Y(h)$ を回収する。
相対的には左項に基底線束の tensor が入ることがあるが、同じ消滅が成り立つ。
評価の自然性により射影束だけでなく $D$ の埋め込みも回収する。
線束の代表を基底線束で tensor すると直像も同じ線束で tensor されるため、
射影束および相対標識は変わらない。
$q:D\to T$ について、各幾何 fiber で $h^0(\OO_D)=1$ であるから、
コホモロジーと基底変換により $q_*\OO_D=\OO_T$ が任意の基底変換後も成立する。
従って線束代表の自己同型は基底からの scalar であり、射影化で消える。
従って構成と合成同型は fppf 降下し、射と基底変換に整合する。
\end{proof}

この定理は\cite[定理5.4]{Q12Weyl2026}の操作である。
ここで証明したのは幾何的 avatar の stack 上の並進であり、
無限二面体群のすべての関係を stack 上へ持ち上げる主張ではない。

\subsection{二つの開 chart と四つの層}

\begin{definition}[幾何的な八本矢・十本矢 chart]
\label{def:qa-open-charts}
$\mathcal U\subset\mathcal A_{12}^{\times}$ を、全幾何 fiber が $\Sigma _0$ である部分とする。
\[
 \mathcal A_{10,\mathrm e}^{\times}:=\mathcal U,
 \qquad
 \mathcal A_{10,\mathrm o}^{\times}:=\mathsf T^{-1}(\mathcal U),
 \qquad
 \mathcal A_8^{\times}:=\mathcal U\cap\mathsf T^{-1}(\mathcal U)
\]
と置く。添字 $\mathrm e,\mathrm o$ は、後述する標準表示で残る偶数対角・奇数対角に対応する。
ここでは定義自体を幾何データだけで与えている。
\end{definition}

\begin{theorem}[avatar の開部分 stack]
\label{thm:qa-open-charts}
次の射は開埋め込みである。
\[
 \mathcal A_8^{\times}\hookrightarrow
 \mathcal A_{10,\mathrm o}^{\times}\hookrightarrow\mathcal A_{12}^{\times},
 \qquad
 \mathcal A_8^{\times}\hookrightarrow
 \mathcal A_{10,\mathrm e}^{\times}\hookrightarrow\mathcal A_{12}^{\times}.
\]
幾何点での層は次の表で与えられる。
\begin{center}
\begin{tabular}{c|c|c}
曲面&曲線型の条件&属する層\\\hline
$\Sigma _0$&$P\not\simeq F^2H^{-1}$&$\mathcal A_8^{\times}$\\
$\Sigma _0$&$P\simeq F^2H^{-1}$&$\mathcal A_{10,\mathrm e}^{\times}\setminus\mathcal A_8^{\times}$\\
$\Sigma _2$&$D\cap S=\varnothing$&$\mathcal A_{10,\mathrm o}^{\times}\setminus\mathcal A_8^{\times}$\\
$\Sigma _2$&$S\leq D$&$\mathcal A_{12}^{\times}\setminus
(\mathcal A_{10,\mathrm e}^{\times}\cup\mathcal A_{10,\mathrm o}^{\times})$
\end{tabular}
\end{center}
最後の行は幾何点集合として閉 locus である。
全 fiber がこの型となる族の充満部分関手に、閉埋め込みを自動的に与える主張ではない。
\end{theorem}

\begin{proof}
$\mathcal E=\pi_*\OO_Y(h)$ の二つの分裂型は
\[
 h^0(B_t,\mathcal E_t(-2))=
 \begin{cases}0&Y_t=\Sigma _0,\\1&Y_t=\Sigma _2\end{cases}
\]
によって区別される。コホモロジーの上半連続性から $\mathcal U$ は開である。
自己同値 $\mathsf T$ による逆像と交わりも開なので、すべての開埋め込みを得る。
残るのは幾何点での条件である。

$\Sigma _0$ で $H=FG$ と書くと
\[
 h^0(B,\mathcal E_P(-2))=h^0(D,GF^{-1}P).
\]
$Z=af+bg$ に対し、$GF^{-1}P$ の次数は $a-b$ である。
真部分因子では
\[
 Z(D-Z)-\deg_Z(GF^{-1}P)=a+3b-2ab>0.
\]
よって補題\ref{lem:qa-support}により、この線束の非零切断は全体に支持される。
全次数は零なので、非零切断があることは線束が自明であることと同値である。
従って $\mathsf T(Y)=\Sigma _2$ となる条件はちょうど $P\simeq F^2H^{-1}$ である。

$\Sigma _2$ で $D\cap S=\varnothing$ なら $H\simeq F^2$ だから
$h^0(D,HF^{-2}P)=h^0(D,P)=0$ であり、$\mathsf T(Y)=\Sigma _0$ となる。
$S\leq D$ なら $C=D-S$ と置く。$C\sim h+2f$ であり、
$\OO_Y(-C)$ は $\pi$ に沿って次数 $-1$ なので acyclic である。
従って $h^0(\OO_C)=1$, $h^1(\OO_C)=0$、標数零より $\Pic^0(C)=0$ を得る。
残余因子の完全列は
\[
 0\longrightarrow P|_C\longrightarrow\OO_D(S)P
 \longrightarrow\OO_S(-2)\longrightarrow0.
\]
$P|_C\simeq\OO_C$ なので $h^0(D,HF^{-2}P)=1$ である。
従って $\mathsf T(Y)=\Sigma _2$ のままである。
最後に $D S=0$ であるから、有効因子 $D$ は $S$ を含むか $S$ を避けるかのいずれかである。
これですべての場合を得る。閉性は二つの開部分の和の補集合であることによる。
\end{proof}

\begin{remark}[矢数の名称と体上の辞書]
幾何的な開埋め込みと表の分解は、上の証明の通り
射影束の分裂型、Picard 類および垂直成分条件から従う。
各層を $8,10,10,12$ 本の最小矢数と対応させる体上の解釈は、
$A_Q^{\mathrm{pres}}$ を対象とする
\cite[定理15.3]{Q12Avatar2026}の同型類分類を用いる。
同じ矢数の相対 QP 表示を比較する際の階数条件は、
上の幾何学的な層の名称からは
推論せず、各代数表示の積と関係について検証する。
\end{remark}


\begin{remark}[「十本矢全体」の二通りの意味]
最小十本矢の二つの層の和は $\mathcal A_8^{\times}$ を含まない。
これに対し上の $\mathcal A_{10,\mathrm o}^{\times}$ と
$\mathcal A_{10,\mathrm e}^{\times}$ は、可縮な pair を許す十本矢表示の
幾何的 chart に対応し、八本矢部分を開部分として含む。
二つの十本矢 chart の和は、最深の十二本矢層の補集合である。
この区別なしに「矢の少ないモジュライが矢の多いモジュライの開部分である」
と述べることはできない。
\end{remark}

\subsection{直積型 collection と全 square avatar}

$\mathcal A_{\mathrm{sq}}^{\times}$ を、順序付きの二つの ruling を持つ
$\Sigma _0$-族、その相対 $(2,2)$ Cartier 因子 $D$、および fiberwise 非自明な
$P\in\Pic^0_{D/T}$ の stack とする。
曲線上の ruling の制限を $F_s,G_s$ と書く。
これは $h=f+g$ と置くことによって $\mathcal U$ と同じ幾何的 stack である。

\begin{corollary}[全 square chart の相対的な対応]
\label{cor:qa-square-chart}
幾何的 stack の同値
\[
 \mathcal A_{\mathrm{sq}}^{\times}
 \xrightarrow[\ \mathsf T^{-1}\ ]{\ \sim\ }
 \mathcal A_{10,\mathrm o}^{\times}
\]
があり、制限線束は
\begin{equation}
 F=F_s,\qquad H=F_sG_sP^{-1},\qquad P=P
 \label{eq:qa-square-dictionary}
\end{equation}
となる。逆には $\mathsf T$ によって得られる $\Sigma _0$ 上で
\[
 F_s=F,\qquad G_s=HPF^{-1}
\]
を取る。square 座標で八本矢部分となる条件は
$P\not\simeq G_sF_s^{-1}$ であり、壁 $P\simeq G_sF_s^{-1}$ は
負断面を避ける $\Sigma _2$ 上の奇数対角十本矢層に移る。
\end{corollary}

\begin{proof}
定義から $\mathsf T^{-1}:\mathcal U\to\mathsf T^{-1}(\mathcal U)$ は同値であり、
その線束の式が\eqref{eq:qa-square-dictionary}である。
逆の式は $H_s=HP=F_sG_s$ による。
元の $\Sigma _0$ で $\mathsf T^{-1}$ が $\Sigma _2$ を与える条件は、
定理\ref{thm:qa-open-charts}の計算を $P^{-1}$ に対して行って
$P^{-1}=F_sG_s^{-1}$、すなわち $P=G_sF_s^{-1}$ である。
その像は $\mathcal A_{10,\mathrm o}^{\times}$ にあるので負断面を避ける。
\end{proof}

この系では反標準因子を滑らか、被約または既約と仮定していない。
従って「十本でちょうど全ての quadric の divisor と $\Pic^0$ が出る」
という記述は、\emph{奇数対角の非最小十本矢 chart と全 square avatar の対応}
として正確に定式化できる。最小十本矢だけの主張ではなく、$P=\OO_D$ も除いている。

代数側の square foundation は
\[
 B_{\mathrm{sq}}=\End_{\to}(P,F_s,G_s,F_sG_s)
\]
であり、可換な基準 collection は
$(\OO,\OO(1,0),\OO(0,1),\OO(1,1))$ である
\cite[定義1.2・定理1.3]{AvatarVSeven2026}。
最初の二項の左 mutation は曲線上で
\[
 (P^2F_s^{-1},P,G_s,F_sG_s)
\]
を与える。共通に $F_sP^{-1}$ を tensor すると
\[
 (P,F_s,F_sG_sP^{-1},F_s^2G_sP^{-1})=(P,F,H,HF)
\]
となり、上の幾何的辞書と一致する。
mutationによる指定型AS代数の対応は定理~
\ref{thm:qa-as10-global-bijection}と
\ref{thm:qa-as10-relative-families}で証明する。
元の体上および仮定付き相対対応は
\cite[定理51.4・系51.5、定理53.2]{AvatarVSeven2026}にある。
幾何的な開埋め込みは独立に証明した。

\subsection{滑らかな楕円型における三者の明示的辞書}
\label{subsec:qa-square-elliptic-pointed-dictionary}

\begin{theorem}[VdB楕円四つ組とsquare avatarの同型類の全単射]
\label{thm:qa-elliptic-vdb-square-avatar-bijection}
$k$を標数零の代数閉体とする。
滑らかな楕円曲線$E$上のVan den Bergh許容四つ組
$(E,L_0,L_1,L_2)$で、各$L_i$は次数二、
\[
 L_1L_0^{-1}\ne\OO_E,\qquad
 L_2L_0^{-1}\ne\OO_E,\qquad
 L_1L_2^{-1}\ne\OO_E
 \tag{VdB-adm}\label{eq:qa-vdb-square-admissibility}
\]
を満たすものの、曲線・三線束を保つ同型類を考える。
これは、\emph{順序付き}二rulingを持つ
滑らかなsquare avatar
$(Y=\PP^1\times\PP^1,D,P)$で
\[
 D\in|\OO_Y(2,2)|,\qquad D\text{ は滑らか},\qquad
 P\in\Pic^0(D)\setminus\{\OO_D\},\qquad
 P\not\simeq G_sF_s^{-1}
 \tag{Sq-cub}\label{eq:qa-square-elliptic-cubic-open}
\]
を満たすものの同型類と自然に全単射する。
ここで$F_s=\OO_D(1,0)$、
$G_s=\OO_D(0,1)$とし、
二rulingの順序は同型の中で保持する。

全単射の\emph{両方向の式}は
\[
 \boxed{
 \begin{aligned}
 (E;L_0,L_1,L_2)&\longmapsto
  (D\hookrightarrow\PP H^0(F_s)^\vee
           \times\PP H^0(G_s)^\vee;\ P),\\
 F_s&=L_2,\qquad
 G_s=L_1L_2L_0^{-1},\qquad
 P=L_2L_0^{-1},\\
 (D;F_s,G_s,P)&\longmapsto
  (E=D;\ L_0=F_sP^{-1},\
   L_1=G_sP^{-1},\ L_2=F_s).
 \end{aligned}}
 \tag{Sq-VdB}\label{eq:qa-square-vdb-explicit}
\]
で与えられる。
この同型類の比較は、周期同型や根を
追加標識にしないVdBのcubic$\mathbb Z$-代数
の幾何的分類を通じて、
\emph{選んだcubic normal foundation}の
相対的な$k$点の辞書になる。

差パラメータ
$\eta=L_1L_0^{-1}$、
$\delta=L_2L_0^{-1}$を用いれば、
\[
 \eta=G_sF_s^{-1},\qquad \delta=P.
 \tag{Sq-diff}\label{eq:qa-square-vdb-differences}
\]
\end{theorem}
\begin{proof}
まず$L_0,L_1,L_2$から
$P=L_2L_0^{-1}$、
$F_s=L_2$、
$G_s=L_1L_2L_0^{-1}$を作ると、
$F_s,G_s$はいずれも次数二で、
$G_sF_s^{-1}=L_1L_0^{-1}$である。
式~\eqref{eq:qa-vdb-square-admissibility}より
$F_s\not\simeq G_s$である。
滑らかな楕円曲線上の相異なる次数二線束の
完全線形系は、それぞれ二重被覆
$E\to\PP^1$を与え、その積は
$\PP^1\times\PP^1$への閉埋込みである。
像は$(2,2)$因子で、反標準因子となる。
また$P\ne\OO_E$、
$P\not\simeq G_sF_s^{-1}$である。
これは曲線同型を保ち、線形系の基底の選択を
$\PGL_2^2$で割ればavatarの同型類を定める。

逆にsquare avatarから
$L_0=F_sP^{-1}$、
$L_1=G_sP^{-1}$、
$L_2=F_s$を作ると、
式~\eqref{eq:qa-square-elliptic-cubic-open}により
$\eta=G_sF_s^{-1}\ne\OO_E$、
$\delta=P\ne\OO_E$、
$\eta\ne\delta$となる。
従ってVdBの許容条件が成立する。
二つの式を順次適用すると三線束と
二つのrulingの完全線形系をそのまま回収するので、
互いに逆である。
VdBの幾何的quintuple復元
\cite[Theorem 5.5.10]{VdB2011}を適用すれば
最後の解釈も従う。
\end{proof}

\begin{theorem}[O-markingを保つruling反射と楕円型の商全単射]
\label{thm:qa-elliptic-avatar-pointed-quotient}
前定理のsquare avatarについてさらに
\[
 G_sF_s^{-1}\not\simeq P^2
 \tag{two-cub}\label{eq:qa-square-two-cubic-open}
\]
を課した開を$\mathcal A_{\mathrm{sq,ell}}^{2c}$とする。
この開上では
\[
 \boxed{\quad
 r_1(F_s,G_s,P)
  =(G_sP^{-1},F_sP,P)
 \quad}
 \tag{Sq-R}\label{eq:qa-square-pointed-reflection}
\]
は、右辺の二つの次数二線束が定める
新しい$(2,2)$埋込みを用いて、
幾何的avatarの同型類の対合を定める。
前定理の式~\eqref{eq:qa-square-vdb-explicit}の下で、
これは論文Bの
\cite[\texttt{lem:elliptic-two-reflections}]{QuadricModuli2026}
の\emph{構造対象$\OO$を保つ}ruling交換
\[
 (L_0,L_1,L_2)\longmapsto
 (L_0L_1L_2^{-1},L_2,L_1)
\]
に一致する。
従って、$y=(\mathcal C,\OO)$が滑らかな
楕円型かつ二つのnormal foundationがともに
cubic型となる\emph{幾何的なO-marked同型類集合}は
\[
 \boxed{\quad
 \bigl|\mathcal A_{\mathrm{sq,ell}}^{2c}(k)\bigr|
    /\langle r_1\rangle
 \ \xrightarrow{\ \sim\ }\
 \bigl|\mathfrak H_{2,\mathrm{ell}}^{2c}(k)\bigr|
 \quad}
 \tag{Sq-H2}\label{eq:qa-square-pointed-bijection}
\]
と全単射する。
左辺はavatar同型類の集合をとってから
$r_1$の軌道で割ることを意味する。
ここでは$S_2$の商の全相対stack同値を
この幾何的全単射だけから推論しない。

一方、helixの一段の移動に対応する
\[
 r_0(F_s,G_s,P)=(G_s,F_sP,P)
 \tag{Sq-T}\label{eq:qa-square-unpointed-shift}
\]
はVdB四つ組の
$(L_0,L_1,L_2)\mapsto(L_1,L_2,L_3)$に
対応するが、一般に$\OO$を変更する。
従って\emph{O-marked}同型類を得るために
$r_0$まで同一視してはならない。

この商の記述は、一般の二つのavatarの
$r_1$-軌道が二元であることを要求しない。
特殊な二捩れ慣性点では二つが同型になる。
\end{theorem}
\begin{proof}
式~\eqref{eq:qa-square-vdb-explicit}を
$r_1$の右辺へ適用すると
\[
 \begin{aligned}
 L'_0&=G_sP^{-2}
       =L_0L_1L_2^{-1},\\
 L'_1&=F_s=L_2,\qquad
 L'_2=G_sP^{-1}=L_1
 \end{aligned}
\]
を得る。新しい二ruling線束が相異なる条件は
$G_sP^{-1}\not\simeq F_sP$、
すなわち式~\eqref{eq:qa-square-two-cubic-open}
に他ならない。
もう一度適用すると$(F_s,G_s,P)$を回収するため
対合である。

論文Bの二つの内在的foundationの系
\cite[\texttt{cor:two-intrinsic-foundations}]{QuadricModuli2026}と、
pointed groupoid比較の系
\cite[\texttt{cor:h2-point-groupoid-foundation}]{QuadricModuli2026}は、指定$\OO$から始まるnormal foundationは
exact同値の下でちょうど二択であり、
二つともcubic型なら他方は上記$R$で得られること、
またfoundation同型はpointed exact同型を
完全に検出することを示す。
Van den Berghの四つ組復元を合わせれば、
二cubic楕円型のpointed同型類に関する
核と像はちょうど$r_1$の同型類軌道となる。
これが式~\eqref{eq:qa-square-pointed-bijection}である。

同様に$r_0$の右辺から得る
$(L'_0,L'_1,L'_2)$は
\[
 (G_sP^{-1},F_s,G_s)=(L_1,L_2,L_3),
 \quad L_3=L_0^{-1}L_1L_2
\]
である。論文Bの同じ補題では、この添字移動は
underlying Abel圏の同値類は保つが、
一般に構造対象$\OO$を保たない。
従って$r_0$を割るのは別の
unpointed moduli問題である。
\end{proof}

\begin{corollary}[OU potentialからsquare avatar、O-marked対象への一般次数]
\label{cor:qa-ou-potential-square-avatar-pointed-degrees}
論文C
~\ref{thm:qa-ou-rigid-potential-quintuple-generic-etale}
の正則半単純な滑らかな楕円開を、
前定理の二cubic開~\eqref{eq:qa-square-two-cubic-open}
および一般慣性同型の稠密開へ制限する。
この開ではpotentialの幾何的四重表示が
VdB四つ組$(E,L_0,L_1,L_2)$を与え、
式~\eqref{eq:qa-square-vdb-explicit}により
\emph{square avatarへの相対的な対応}がある。
全ての線束・曲線が相対的に構成される
étale chart上では、これは基底変更と可換な
代数スタックの射となり、
相対的な曲線同型を通じて降下する。
一方、対応の幾何的fiberの一般次数は
\[
 \boxed{
 \begin{array}{ccc}
 \text{有効potential表示}
   &\xrightarrow{\;4\;}&
  \text{順序付きsquare avatar}\\
 &&\big\downarrow\vcenter{\rlap{$\scriptstyle 2$}}\\
 &&\text{O-marked }\mathfrak H_2
 \end{array}}
 \tag{Pot-Sq-H2}\label{eq:qa-potential-square-pointed-degrees}
\]
であり、potentialからO-marked対象への
一般次数は$4\cdot2=8$である。
ここで上の$4$はpotential/quintupleの
\emph{表示間}の一般四重射、
下の$2$は二つのnormal foundationの忘却であり、
$\Pic^0$の四つの平方根の選択を
追加標識として数えた次数ではない。
\end{corollary}
\begin{proof}
Cの
~\ref{thm:qa-ou-rigid-potential-quintuple-generic-etale}
による一般有限étale四重射は、
指定された幾何quintupleの相対四重テンソルを与える。
滑らかな正則半単純開では、C
~\ref{cor:qa-ou-relative-regular-semisimple-potential}
の点曲線と二rulingの相対構成が、
基底変更と可換な三つの次数二線束を与える。
それらのテンソル積と二つの完全線形系で
式~\eqref{eq:qa-square-vdb-explicit}の
square avatarを相対的に作る。
曲線・線束の同型を保つので基底変更と降下に
関して自然である。
VdBの四つ組の幾何的同型類とsquare avatarの
全単射は定理
~\ref{thm:qa-elliptic-vdb-square-avatar-bijection}で証明した。
これとCの
~\ref{cor:qa-quintuple-h2-generic-degree-two}の
二つのfoundationの一般次数を合成すれば
表示された幾何一般次数が従う。
\end{proof}

\begin{remark}[一cubic境界と古典点での辞書の射程]
\label{rem:qa-square-avatar-boundary-cautions}
$G_sF_s^{-1}=P^2$では、$r_1$後の
$F'_s=G_sP^{-1}$と$G'_s=F_sP$が一致し、
square$(2,2)$埋込みを構成できない。
これは論文C
~\ref{cor:qa-ou-one-cubic-fiber-and-nonfiniteness}
の一cubic境界であり、
二つ目のnormal foundationはnoncubicとして
内在的な選択被覆
~\ref{thm:qa-global-two-ruling-etale-cover}
に残る。

また$P=\OO_D$を任意の$(2,2)$因子に
そのまま許すと、これらの曲線から得る
可換な四項foundation代数は$D$に依存せず、
同じ可換quadricを重複して数える。
従って全同型類の全単射には古典点を
別途一点として追加し、元のavatarの
$P=0$座を単に追加するだけでは足りない。
$\Sigma_2$の古典型は別のweighted
stackを生み、無条件には$\HHstack$に入らない。
$\Sigma_2$曲線型でのO-marked適合性、
非被約反標準因子を含む全対応、
およびpotentialの安定化・自己同型群の比較は
式~\eqref{eq:qa-square-pointed-bijection}
の証明には含まれない。
\end{remark}

\begin{proposition}[二つのwallは異なり、noncubic avatar壁には追加認識が必要]
\label{prop:qa-avatar-walls-and-noncubic-dimension}
滑らかな楕円型square avatarについて、
$\eta=G_sF_s^{-1}$、$\delta=P$と書く。
次の二つのwallを区別する。
\begin{enumerate}
\item $\eta=2\delta$、$\delta\notin E[2]$：
式~\eqref{eq:qa-square-vdb-explicit}の四つ組は
geometric cubic型であるが、
$r_1$後の第二foundationはnoncubicになる。
したがって元のO-marked点は
\emph{cubic locusに属する}。
\item $\eta=\delta\ne0$：
式~\eqref{eq:qa-square-vdb-explicit}の
$L_1=L_2$なので元のfoundationが
geometric cubic型ではない。
このwallはsquareの八本矢開の補集合であり、
曲線型avatarの十本矢・$\Sigma_2$ chartへ
移行する。
\end{enumerate}
後者の全smooth楕円avatar同型類の
パラメータ空間は、楕円曲線の$j$と
非零点$\delta\in E$を局所座標として
\emph{二次元}である。
一方、論文Bの内在noncubic境界
$\mathfrak N\subset\HHstack$は
\emph{純一次元}のArtinスタックである。
従ってこのwall\emph{全体}を
内在noncubic境界と表現可能quasi-finiteな
モジュライ同型で結び付けることはできない。
O-marked適合性はこのwall上で
さらに制限されるか、あるいは
同じpointed対象へ正の相対次元を持って
潰れるかのいずれかである。

特に非零二捩れ$P$の条件$P^2=\OO_E$は
楕円曲線のmoduli上で一次元を定め、
次に検討すべき\emph{候補}ではあるが、
それが内在noncubic境界をちょうど尽くすことは
ここでは証明していない。
\end{proposition}
\begin{proof}
最初のwallでは$L_0=F_sP^{-1}$、
$L_1=G_sP^{-1}$、$L_2=F_s$が
互いに相異なり、VdB許容条件を満たす。
$r_1$による新しい二線束は
$G_sP^{-1}=F_sP$と$F_sP$で一致するので、
第二はnoncubicとなる。
これはC
~\ref{cor:qa-ou-one-cubic-fiber-and-nonfiniteness}
の一cubic境界に一致する。
次のwallでは$G_s=F_sP$だから
$L_1=G_sP^{-1}=F_s=L_2$となる。
$\mathcal A_8^\times$の条件
~\eqref{eq:qa-square-elliptic-cubic-open}
に反することも直接に分かる。
$\Sigma_2$との対応は
系~\ref{cor:qa-square-chart}に従う。

最後の次元を計算する。
楕円曲線$E$のmoduliは一次元で、
$F_s\in\Pic^2(E)$の選択は曲線の
translationによって動かせる。
標数零では$[2]:E\to E$が全射なので、
このtranslationの軌道は$\Pic^2(E)$全体である。
$\eta=\delta$の下では$G_s=F_sP$だから
残るパラメータは
$E$と$0\ne P\in\Pic^0(E)$で、
そのmoduliの一般次元は$1+1=2$。
一方、論文B
\cite[\texttt{thm:noncubic-boundary-smooth-codim-two}]{QuadricModuli2026}
は全$\HHstack$が純次元三で
$\mathfrak N$が滑らかな余次元二、
すなわち純次元一であることを証明している。
有限型quasi-finite射の像は源の各成分と同じ
次元を持つので、この二次元wall全体を
一次元の境界にquasi-finiteに写すことはできない。
最後の二捩れ条件の次元も
有限étaleな$E[2]\setminus0$から従う。
\end{proof}

\begin{theorem}[二捩れsquare壁とnoncubic二根対の幾何同型類の全単射]
\label{thm:qa-avatar-two-torsion-root-pair-bijection}
$k$を標数零の代数閉体とする。
次の二集合を考える。
\begin{enumerate}
\item 滑らかな楕円曲線$E$、次数二線束$F_s$、
非自明な$P\in\Pic^0(E)[2]$の三つ組
$(E,F_s,P)$の曲線・線束・Picard類を
保つ同型類集合。$G_s=F_sP$とすると
$F_s\not\simeq G_s$だから、二完全線形系は
$(2,2)$埋込みを定め、対応するsquare avatarは
壁$P=G_sF_s^{-1}$上にある。
\item 射影直線$B\simeq\PP^1$上の、
互いに交わらない\emph{被約}二点因子の
非順序対$\{B_1,B_2\}$で、
$\deg B_i=2$かつ$B_1+B_2$は
四つの相異なる点からなるものの、
$\PGL_2$同型類集合。
\end{enumerate}
この二集合の間に自然な全単射がある。
対応は、$F_s$の二重被覆
\[
 \pi_F:E\longrightarrow\PP H^0(E,F_s)^\vee
\]
の四つの分岐点に対し、
$P$に対応する非零二捩れ平行移動が
誘導する二組の分割を取ることで与えられる。
逆方向には、四分岐点$B_1+B_2$の
二重被覆$\pi_F:E\to B$を取り、
その分割に対応する唯一の
非零$E[2]$の元を$P$とする。
このとき$F_s=\pi_F^*\OO_B(1)$である。

論文B
\cite[\texttt{thm:noncubic-intrinsic-presentation-equivalence}]{QuadricModuli2026}
の非cubic根対表示によれば、
従ってこの二捩れsquare壁の幾何的avatar同型類は、
\emph{四根が相異なるnoncubic O-marked H2境界}
の同型類とも全単射する。

最後の全単射は\emph{論文Bの根対分類を介した}
同型類の対応である。
この幾何辞書だけから
avatar側のAS$\mathbb Z$-代数の
$\qgr$とそのnoncubic O-marked heartの
同値、または対応するモジュライstackの射は
まだ結論しない。
\end{theorem}
\begin{proof}
次数二の線束$F_s$は楕円曲線から$\PP^1$への
二重被覆を与える。
その被覆対合を$\iota_F$と書くと、
その固定点集合$R_F$は四点からなり、
像$B_F$が四つの相異なる分岐点となる。
$E$は標識なしでも$\Pic^1(E)$のtorsorで、
$P\in\Pic^0(E)[2]\setminus0$は
$E$上の固定点のない位数二平行移動$t_P$を
定める。
次数二線束に対する平行移動の作用は
$t_P^*F_s\simeq F_sP^{\pm2}\simeq F_s$なので、
$t_P$は$\iota_F$と可換であり、
$R_F$を保つ。
その四点への作用は自由であり、
ちょうど二つの長さ二の軌道に分ける。
二つの軌道の$\pi_F$による像が
$B_F=B_1+B_2$の二分割となる。
これは線束の基底・楕円原点の選択に依存せず、
曲線と線束の同型に自然である。

逆に四つの相異なる分岐点$B_1+B_2$から、
$B$上の二重被覆は、
$\mathcal O_B(2)$とその四点を零点に持つ
切断を用いて、被覆同型を除き一意に構成される。
$F_s=\pi_F^*\OO_B(1)$は次数二である。
四つの分岐点は$E$上の被覆対合の固定点と
一対一に対応する。任意の一つを楕円原点に
選べば、その四点が$E[2]$となる。
$E[2]\simeq(\mathbb Z/2)^2$の各非零元の
平行移動は四点を二つの組に分割し、
三つの非零元と三つの二分割との間に
自然な全単射を与える。
従って指定の$\{B_1,B_2\}$から
一意の非零二捩れ$P$を回収する。
原点を変えても$\Pic^0(E)$における$P$の
同型類は変わらない。
二方向の構成が互いに逆であることも
この四点の分割から直ちに従う。

具体的に、適当な射影座標で
$B_1=\{0,\infty\}$、$B_2=\{1,\lambda\}$、
$\lambda\notin\{0,1,\infty\}$とすれば、
$E$はLegendre曲線
\[
 y^2=x(x-1)(x-\lambda)
\]
で与えられる。$B_1,B_2$をそれぞれ
集合として保ち、その二点を交換する
$\PP^1$の対合は$x\mapsto\lambda/x$である。
その固定点を持たないliftは
\[
 (x,y)\longmapsto
 \left(\frac{\lambda}{x},
             -\frac{\lambda y}{x^2}\right)
 \tag{Root-2}\label{eq:qa-avatar-root-pair-translation}
\]
であり、$E$上では非零二捩れ点
$(0,0)$によるtranslationである。
この式からも二分割と二捩れ点の
相互復元を直接確認できる。

最後に論文Bの非cubic被約境界は、
二つの二根集合の非順序対を
$\PGL_2$で割った群oidとして記述されており、
四根相異なる開ではまさに第二集合である。
従ってBの分類との合成が最後の
同型類全単射を与える。
これはBのO-marked対象への
逆構成による対応であり、
元のavatar AS代数の直接認識とは区別する。
\end{proof}

\subsection{二根対のquartic resolventと退化avatarの非可逆Picardパラメータ}
\label{subsec:qa-rootpair-resolvent-singular-avatar}

\begin{theorem}[noncubic根対の三重quartic resolventと三つの退化型]
\label{thm:qa-noncubic-rootpair-quartic-resolvent}
$k$を標数零の代数閉体、$W$を二次元ベクトル空間とし、
$\mathcal D_2=\PP(\operatorname{Sym}^2W^\vee)$、
$\mathcal D_4=\PP(\operatorname{Sym}^4W^\vee)$と置く。
二つの二次形式をその順序を忘れて掛ける射
\[
 \mu:\operatorname{Sym}^2(\mathcal D_2)\longrightarrow
             \mathcal D_4,\qquad
 \{[f_1],[f_2]\}\longmapsto[f_1f_2]
 \tag{Res-4}\label{eq:qa-rootpair-resolvent-map}
\]
は全射有限射であり、一般次数は$3$である。
根が相異なる四次形式$q$の上の三点は
四根の非順序二根対への三分割を表す。
アフィンmonic chartで
\[
 q(x)=x^4+a x^3+b x^2+c x+d
\]
とすると、四根$\alpha_1,\ldots,\alpha_4$の
二組分割から得る
$\alpha_i\alpha_j+\alpha_k\alpha_\ell$の
三つの値は
\[
 R_q(z)=z^3-bz^2+(ac-4d)z
             -(a^2d+c^2-4bd)
 \tag{Res-cub}\label{eq:qa-rootpair-resolvent-cubic}
\]
の三根に一致する。
この三重射は古典的なquartic resolventの
幾何的な表示である
\cite[Section 4.3]{Cremona1999}。

$\mathcal D_2^{\mathrm{sf}}$を二つの相異なる
根を持つ二次形式の開部分とする。
論文Bのnoncubic根対表示の
\[
 \mathcal R_{\mathrm{sf}}=
 \operatorname{Sym}^2(\mathcal D_2^{\mathrm{sf}})
\]
上では$q=f_1f_2$の幾何的根配置は
次の三型に限られ、それぞれ
\[
 C_q:\quad y^2=f_1f_2
 \quad\text{in }\operatorname{Tot}\OO_{\PP^1}(2)
 \tag{Res-curve}\label{eq:qa-rootpair-spectral-curve}
\]
の射影的二重被覆が次の形となる。
\[
 \begin{array}{c|c|c}
 \text{根配置}&q\text{の重複度}&C_q\\\hline
 B_1\cap B_2=\varnothing&(1,1,1,1)&\text{smooth genus }1\\
 |B_1\cap B_2|=1&(2,1,1)&\text{既約nodal genus }1\\
 B_1=B_2&(2,2)&\text{二つの}\PP^1\text{が二節点で交わる}
 \end{array}
 \tag{Res-strata}\label{eq:qa-rootpair-three-strata}
\]
三型とも$C_q$は連結なGorenstein曲線で
算術種数$1$である。
この構成は$\operatorname{PGL}(W)$と
根対交換の作用に自然であり、論文Bの
根対スタック
$\mathfrak R=[((\mathcal D_2^{\mathrm{sf}})^2)/
(\operatorname{PGL}(W)\times S_2)]$
からquartic二重被覆のスタックへの射を与える。
ただし全$\operatorname{Sym}^2(\mathcal D_2)$の
有限射と、平方因子を除いた
$\mathcal R_{\mathrm{sf}}$上の制限射は区別する。
\end{theorem}
\begin{proof}
四次形式の任意の因数分解は
その根の重複度付き四点集合から
二点ずつを選ぶ有限個の操作なので、
$\mu$の幾何fiberは有限である。
$\operatorname{Sym}^2(\mathcal D_2)$は射影的、
従って$\mu$はproperかつquasi-finiteで
有限である。代数閉体上ではすべての四次形式が
二次形式の積となるので全射。
四根が相異なれば分割は三つあり、
一般次数は三となる。
式~\eqref{eq:qa-rootpair-resolvent-cubic}は
四根の基本対称式を展開すれば得られる。
例えば三値を
$u_{12|34}=\alpha_1\alpha_2+\alpha_3\alpha_4$
などと置くと、
$\sum u=b$、
$\sum_{r<s}u_ru_s=ac-4d$、
$\prod u=a^2d+c^2-4bd$である。

$f_1,f_2$の根が各々相異なるので、
積の各根の重複度は高々二である。
共通根の個数が$0,1,2$の三場合が
表の各行に一致する。
四単根なら二重被覆はsmoothでRiemann--Hurwitzにより
genus一。一つの二重根の近くでは
$y^2=t^2u(t)$、$u(0)\ne0$なので、
étale局所で$y'^2=t^2$という普通節点になる。
残り二単根から二重被覆は既約で、
正規化はgenus零。
二つの二重根なら$f_1=f_2$（射影類として）で、
$y^2=f_1^2$は二つのsection
$y=\pm f_1$の和である。
各sectionは$\PP^1$で、$f_1$の二単根で横断する。
二重被覆の代数
$\OO_{\PP^1}\oplus\OO_{\PP^1}(-2)y$
から$\chi(\OO_{C_q})=1+(-1)=0$であり、
全型で$p_a(C_q)=1$。
局所的に平面曲線の節点またはsmoothなので
Gorensteinでもある。
群の作用については積と二重被覆の
構成が自然であることから従う。
\end{proof}

\begin{theorem}[根対からの相対torsion-free spectral avatar]
\label{thm:qa-noncubic-rootpair-torsionfree-spectral-module}
論文Bの二形式表示
\[
 (f_1,f_2)\in
 \bigl(\operatorname{Sym}^2W^\vee\bigr)_{\mathrm{nd}}^2
\]
上で、$B=\PP(W)$と置き、
$\widetilde p:\widetilde C\to B$、
$p:C\to B$を以下の有限射、
$r:\widetilde C\to C$をその間の自然な射とする。
\[
 \widetilde C=
 \operatorname{Spec}_{\PP(W)}
  \frac{\OO_{\PP(W)}[c,d]}
       {(c^2-f_1,\ d^2-f_2)},\qquad
 C=\operatorname{Spec}_{\PP(W)}
  \frac{\OO_{\PP(W)}[y]}{(y^2-f_1f_2)}
\]
を取る。ただし$c,d$は局所的に
$\OO_{\PP(W)}(-1)$の生成元、$y=cd$は
$\OO_{\PP(W)}(-2)$の生成元と解釈する。
$\widetilde C$から$C$への有限射は
$(c,d)\mapsto(-c,-d)$の商であり、
$\OO_C$-加群としての反不変部分と
その$B$への直像はそれぞれ
\[
 \mathcal M=(r_*\OO_{\widetilde C})^-,
 \qquad
 p_*\mathcal M
 \cong\OO_B(-1)c\oplus\OO_B(-1)d
 \tag{Res-M}\label{eq:qa-rootpair-spectral-module}
\]
であり、$\mathcal M$は
\[
 y c=f_1 d,\qquad y d=f_2 c
 \tag{Res-MF}\label{eq:qa-rootpair-spectral-matrix-factorization}
\]
を満たす。全二形式族と任意の基底変更にわたり、
$C$は有限平坦二重被覆、$\mathcal M$は
\emph{相対平坦なrank-one torsion-free sheaf}
（各幾何fiber上で次数零）となる。
この構成は二形式の基底変更・スカラー変更・
二形式交換に関して自然なので、
論文Bのnoncubic表示スタック
$\mathfrak N^{\mathrm{pres}}=[P/G]$から
「種数一の有限二重被覆とrank-one
torsion-free sheaf」の相対群oidへの射を与える。

二根対が交わらない滑らかなfiberでは、
$\widetilde C\to C$はétale $\mu_2$-torsor、
$\mathcal M$は非自明な
$\Pic^0(C)[2]$の線束である。
これは定理
~\ref{thm:qa-avatar-two-torsion-root-pair-bijection}
で定めた分割の2-torsionと一致する。
他方、根対が一根を共有するとき
$\mathcal M$はその節点で局所自由でなく、
二根対が一致するときは二つの節点で
局所自由でない。
従って\emph{全noncubic根対配置にわたって}
自然に存在するのは通常の$\Pic^0$線束ではなく、
この相対torsion-free spectral moduleである。
\end{theorem}
\begin{proof}
二つの二次平方根の代数は、
$\OO_B$上の基底$1,c,d,cd$を持つ
階数四の局所自由代数である。
$\widetilde p_*\OO_{\widetilde C}$の
同時符号変換の不変部分は$1,cd$、
反不変部分は$c,d$なので
式~\eqref{eq:qa-rootpair-spectral-module}を得る。
$y=cd$による積は
式~\eqref{eq:qa-rootpair-spectral-matrix-factorization}
で与えられ、形式的には
\[
 \begin{pmatrix}0&f_2\\f_1&0\end{pmatrix}^2
      =f_1f_2 I_2
\]
というmatrix factorizationである。
$p_*\mathcal M$は$B$上で
$\OO(-1)^{\oplus2}$だから、
二形式パラメータとその任意基底変更に対して平坦。
各幾何fiberでは二次拡大のgeneric点でrank一で、
局所的に二次平方根代数の直和因子であるから
torsion-freeである。
Euler特性は
\[
 \chi(C,\mathcal M)
 =\chi(\PP^1,\OO(-1)^{\oplus2})=0
 =\chi(C,\OO_C)
\]
なので次数零である。

$f_1,f_2$に共通零点がなければ、
$c,d$が同時に零となる点がないため
同時符号対合は自由で、
$\widetilde C\to C$は有限étaleの二重torsor。
このとき反不変部分$\mathcal M$は
invertibleかつ$\mathcal M^2\simeq\OO_C$。
$H^0(C,\mathcal M)=H^0(B,p_*\mathcal M)
 =H^0(\PP^1,\OO(-1)^{\oplus2})=0$
なので非自明である。
$\widetilde C$の三つの二次部分被覆の
分岐分割から、これは前定理の2-torsionに一致する。

共有根の局所では単元の平方根をétale局所に取り、
\[
 \OO_{C,o}^{\wedge}
  \simeq k[[x,y]]/(y^2-x^2),\qquad
 \mathcal M_o\cong
   \langle c,d\mid yc=xd,\ yd=xc\rangle .
\]
$\mathcal M_o/(x,y)\mathcal M_o$の次元は二なので、
rank一局所自由ではない。
一致根対では$C$は二つのsectionの和で、
その二つの交点の局所方程式も同じになり、
$\mathcal M$は各節点で非局所自由である。
二形式の交換は$c,d$の交換、
スカラー変更は平方根のスカラー変更として
作用し、構成は元の表示群の作用に整合する。
\end{proof}

\begin{corollary}[noncubic中央二重平方環の可換点曲線からの幾何射]
\label{cor:qa-noncubic-central-square-point-curve-spectral-avatar}
論文A・Bのnoncubic中央二重平方環
\[
 B_f=k[u,v]\langle c,d\rangle/
          (c^2-f_1(u,v),\,d^2-f_2(u,v))
\]
に対し、通常の可換化を施した環と
その射影曲線は
\[
 B_f^{\mathrm{ab}}=
 k[u,v,c,d]/(c^2-f_1,\,d^2-f_2),\qquad
 E_f=\operatorname{Proj}(B_f^{\mathrm{ab}})
 \subset\PP^3
\]
である。これは定理
~\ref{thm:qa-noncubic-rootpair-torsionfree-spectral-module}
の四重被覆$\widetilde C$と自然に同型であり、
\[
 \boxed{\quad
 E_f\xrightarrow{/\,\mu_2}C_{f_1f_2}
 \xrightarrow{\;2:1\;}\PP^1,
 \qquad
 \mathcal M=(r_*\OO_{E_f})^- .
 \quad}
 \tag{NC-point-avatar}\label{eq:qa-noncubic-pointcurve-spectral-factorization}
\]
$\mu_2$は$(c,d)\mapsto(-c,-d)$として作用する。
この対合の固定点は$f_1,f_2$の
共通根の上に限られるので、二根対が
非交差なら第一射はétale、
一根共有または一致根対では
固定点に対応する節点でétaleでなくなる。
この場合、有限平坦な二重torsorと
みなしてはならない。

この射はnoncubic線形代数$B_f$から
rank-one torsion-free spectral avatarへの
\emph{具体的な幾何構成}である。
非可換な$B_f$の$\qgr$と
$\OO$-marked avatar側のheartが同値である
と主張するものではない。
\end{corollary}
\begin{proof}
$u,v$は中央にあるので、可換化において
追加される関係は$cd=dc$だけである。
そのProJの$\PP^1=\operatorname{Proj}k[u,v]$上の
相対座標環は$1,c,d,cd$を基底とする
二重平方根代数となる。
これは前定理の$\widetilde C$そのものである。
同時符号の不変部分は$1,cd$、
反不変部分は$c,d$なので、
商と$\mathcal M$の表示も一致する。
非自明な同時符号の固定点は、
射影座標$(u,v)$が非零である限り
$c=d=0$を満たす点に限られ、
これは$f_1=f_2=0$と同値である。
$(u,v)=(0,0)$では二つの平方関係から
$c=d=0$となり射影点は存在しない。
従って固定点の条件が従う。
\end{proof}

\begin{proposition}[共有根退化での可逆Picard延長の局所障害]
\label{prop:qa-rootpair-nodal-picard-nonextension}
$R=k[[s]]$上の二根対の退化
\[
 f_1=x(x-1),\qquad f_2=(x-s)(x-2)
\]
を取り、$C_R:y^2=f_1f_2$とする。
その一般fiberは滑らかな楕円曲線であり、
特殊fiberは既約な一節点曲線となる。
一般fiberの自然な非零2-torsion線束
$\mathcal M_K$は、
\emph{$C_R$全体の可逆線束には延長できない}。
従って二根対から
$C:y^2=f_1f_2$と
\emph{可逆な}$P\in\Pic^0(C)$を
割り当てる滑らかなavatar辞書は、
この族に沿ってそのまま相対射としては延長できない。

この障害は、$C_R$の節点における
表面$A_1$特異点と、
そこでの非自明なrank-one reflexive module
によって検出される。
$C_R$以外の曲線へのavatar再埋込み、
$\Sigma_2$へのelementary transform、
またはcompactified Picard型データ
\cite{AltmanKleiman1980}による
延長の可能性は排除しない。
ただし同じ一径数基底の上で全空間の$A_1$特異点を
最小解消するだけでは不十分である。
命題~\ref{prop:qa-nodal-resolution-odd-exceptional-degree}の
反不変moduleの可逆化は例外曲線上の次数が奇数であり、
任意の垂直Cartier因子でtwistしても
相対的な多重次数零線束にはならない。
第~\ref{sec:qa-spectral-boundary-models}節の二traitでは，
任意の有限DVR基底変更後の最小正則模型でも二乗自明な線束への延長は存在しない
（定理~\ref{thm:qa-boundary-input-two-torsion-line-obstruction}）．
線束自体への延長と，非可逆TF層・二次積を保持する標的は区別する．
\end{proposition}
\begin{proof}
$s=0$では$f_1f_2=x^2(x-1)(x-2)$であり、
特殊曲線は既約nodal genus一である。
共通根$x=s=0$の完成局所環では、
単元の平方根で座標を正規化すると
\[
 A=k[[x,s,y]]/(y^2-x(x-s)).
\]
二形式を別々に平方根化する被覆の完成環は
$B=k[[c,d]]$で、
\[
 x=c^2,\quad x-s=d^2,\quad y=cd,\quad
 A\simeq B^{\{\pm1\}}
\]
となる。$B$上の同時符号の反不変部分
$M=Ac+Ad$は、商$A$上の
rank-one maximal Cohen--Macaulay、従って
rank-one reflexive moduleである。
その最小生成元は$c,d$の二つなので非自由であり、
$A$は表面$A_1$特異点を持つ。
これは前定理の$\mathcal M$の完成stalkである。

相対曲線の全空間$C_R$はnormalであり、
その唯一の特異点で上記$A_1$局所型を持つ。
$\mathcal M$はnormal曲面上の
rank-one reflexive sheafで、
一般fiberでは$\mathcal M_K\in\Pic^0(C_K)[2]$。
もし可逆線束$L$が全空間上に存在して
$L_K\simeq\mathcal M_K$となれば、
normal曲面上のrank-one reflexive sheafの
Weil因子記述から、$L$と$\mathcal M$の差は
特殊fiberに支持されるWeil因子となる。
この特殊fiberは\emph{既約}かつ
パラメータ$s$で切られたCartier因子なので、
その整数倍は主因子である。
従って$L\simeq\mathcal M$となり、
節点での非局所自由性に矛盾する。
最後の留保は$C_R$を固定したときの
可逆延長障害と、別のavatarモデルの存在を
混同しないためである。
\end{proof}

\begin{theorem}[全noncubic根対と二捩れsquare avatarの体上同型類全単射]
\label{thm:qa-noncubic-rootpair-all-geometric-square-avatar-bijection}
$k$を標数零の代数閉体とする。
論文Bのnoncubic根対の全幾何点、すなわち
$\PP^1$上の\emph{各々は被約}な二点因子の
非順序対$\{B_1,B_2\}$の
$\PGL_2$-同型類集合を考える。
これは、次の条件を持つsquare avatarの
\emph{同型類集合}と自然に全単射する。
\[
 \begin{gathered}
 (Y=\PP^1\times\PP^1;\
  D\in|\OO_Y(2,2)|;\ F_s,G_s,P),\\
 P=G_sF_s^{-1}\ne\OO_D,\qquad
 P^{\otimes2}\simeq\OO_D,\\
 D\text{ はsmooth genus }1,\quad
 \text{既約一節点、または二}\PP^1
 \text{が二節点で交わる曲線}.
 \end{gathered}
 \tag{Sq-all-nc}\label{eq:qa-all-noncubic-square-avatar-bijection}
\]
ここで$F_s=\OO_D(1,0)$、
$G_s=\OO_D(0,1)$であり、
二rulingの順序は保持する。
この全単射は、論文Bの
\cite[\texttt{thm:noncubic-intrinsic-presentation-equivalence}]{QuadricModuli2026}
を介して、\emph{全noncubic O-marked H2幾何同型類}の
集合論的なavatar分類を与える。

構成は次の通りである。
まず$q=f_1f_2$から
式~\eqref{eq:qa-rootpair-spectral-curve}の
二重被覆$p:C_q\to\PP^1$と
$F_s=p^*\OO_{\PP^1}(1)$を取る。
\begin{enumerate}
\item 四根相異なる場合は、
定理~\ref{thm:qa-avatar-two-torsion-root-pair-bijection}
の分割に対応する非零
$P\in\Pic^0(C_q)[2]$を取る。
\item 一根共有または二根対一致の場合は
$\Pic^0(C_q)\simeq\mathbb G_m$なので、
その\emph{唯一の}非自明二捩れ点$P=-1$
を取る。
\end{enumerate}
いずれの場合も$G_s=F_sP$は次数二で
大域生成され、二つの完全線形系の積
\[
 C_q\longrightarrow
 \PP H^0(F_s)^\vee\times\PP H^0(G_s)^\vee
\]
は$(2,2)$反標準Cartier因子への閉埋込みとなる。
この逆操作は、square avatarの
第一rulingの二重被覆$D\to\PP^1$の
分岐四次因子と$P$から、
四根相異なる場合には三分割の一つを、
一根共有・一致根対では唯一の
被約二根対分割を復元することである。

\emph{重要：} これは\emph{体上の同型類の全単射}である。
上の二場合で特別に選んだ
$P=-1\in\mathbb G_m$は、
必ずしも滑らかな族の2-torsion線束の極限でない。
実際、命題
~\ref{prop:qa-rootpair-nodal-picard-nonextension}
の一根共有族では延長できない。
従ってこの全単射を、そのまま
通常のPicard avatarスタックへの
相対モジュライ射、または元のavatar
AS代数の$\qgr$のpointed同値と
読み替えてはならない。
\end{theorem}
\begin{proof}
smoothな場合は定理
~\ref{thm:qa-avatar-two-torsion-root-pair-bijection}
による。以下、特異な二場合を調べる。

\emph{既約一節点の場合。}
$C_q$の正規化は$\nu:\PP^1\to C_q$で、
節点の二逆像を$a,b$とする。
$\Pic^0(C_q)$の線束は、
正規化上の自明線束の
$a,b$での貼合せスカラー
$\lambda\in k^\times$で分類される。
従って$\Pic^0(C_q)\simeq\mathbb G_m$で、
非自明二捩れは貼合せ比$-1$だけである。

$F_s=p^*\OO(1)$は正規化上次数二で、
節点での貼合せ比を$1$に規格化できる。
$G_s=F_sP$は正規化上の$\OO(2)$で、
貼合せ比は$-1$となる。
$H^0(C_q,G_s)$は$H^0(\PP^1,\OO(2))$の
$a,b$における値が符号反対になる
二次元部分空間である。
これはbasepoint-freeである。
実際、$a,b$の二評価汎関数の
非自明な線形結合は、
二次Veroneseの双対円錐上で
単一点の評価汎関数には一致しない。
（異なる二点を結ぶ割線が二次曲線と
交わるのはその両端のみである。）
従って$|G_s|$も有限な次数二射を与える。

$F_s\not\simeq G_s$なので、
この二射の積は総次数二ではあり得ず、
$C_q$上の一般点で単射である。
像の二rulingへの次数は各二だから
$(2,2)$因子で、算術種数は一。
元の$C_q$も算術種数一なので、
有限双有理射の構造層の余核はEuler特性零、
従って零であり、積射は閉埋込みとなる。

\emph{二根対一致の場合。}
$C_q=C_+\cup C_-$は二つの$\PP^1$が
相異なる二節点で交わる曲線である。
正規化上の各成分の次数零線束は自明であり、
二つの節点での貼合せ比
$(\lambda_1,\lambda_2)$を
成分ごとの定数倍で同一視すると、
$\lambda_1/\lambda_2\in\mathbb G_m$
だけが残る。
よって$\Pic^0(C_q)\simeq\mathbb G_m$、
非自明二捩れはこの比が$-1$のものに限る。
$F_s$は各成分で$\OO_{\PP^1}(1)$であり、
$G_s=F_sP$も各成分で次数一、
二節点での値の二条件から
$H^0(C_q,G_s)$は二次元で大域生成される。
二つの成分の像は、
第一$\PP^1$から第二$\PP^1$への
二つの$(1,1)$グラフとなる。
貼合せ比$-1$ゆえ二グラフは異なり、
ちょうど二つの節点で横断する。
従って像は元の$C_q$と同型な
$(2,2)$Cartier因子である。

逆向きは平方freeな二つの二次因子の
分割を数えることで得られる。
四単根なら三通りで非零$E[2]$の三点に対応する。
$(2,1,1)$なら
$(a,b)\mid(a,c)$の一通り、
$(2,2)$なら
$(a,b)\mid(a,b)$の一通りしか
二因子ともsquarefreeにする分割はない。
同じ各曲線の$\Pic^0[2]\setminus0$の
点数が$3,1,1$なので、逆は一意である。
元の第二rulingは
$G_s=F_sP$で復元されるから、
曲線・線束の同型はavatarの同型となる。
最後にBのnoncubic根対分類と合成すれば
O-marked同型類への全単射が得られる。
相対射に関する留保は直前命題の
非延長例により必要である。
\end{proof}

\begin{theorem}[全noncubic O-marked同型類の$\Sigma_2$反標準avatar分類]
\label{thm:qa-all-noncubic-sigma2-anticanonical-avatar-bijection}
$k$を標数零の代数閉体とする。
$\Sigma_2=\PP_{\PP^1}(\OO\oplus\OO(2))$の
負断面を$S=h-2f$とし、
$D\in|2h|=|-K_{\Sigma_2}|$を
$D\cap S=\varnothing$を満たす
相対反標準Cartier因子とする。
このとき、次の二集合は自然に全単射する。
\[
 \boxed{
 \begin{aligned}
 &\{\text{全noncubic O-marked quadricの
             幾何exact同型類}\}\\
 &\quad\longleftrightarrow
 \left\{
 \begin{array}{l}
 (D\subset\Sigma_2,P),\quad
 D\cap S=\varnothing,\\
 0\ne P\in\Pic^0(D)[2]
 \end{array}
 \right\}
/\operatorname{Aut}(\Sigma_2;h,f).
 \end{aligned}}
 \tag{NC-Sigma2}\label{eq:qa-all-noncubic-sigma2-avatar}
\]
右辺で$D$はあらかじめ滑らか、既約、被約と
仮定しない。しかし$P$の非自明二捩れ性から、
実際には$D$はsmooth genus一、
既約一節点、または二つの$\PP^1$が二節点で
交わる型のいずれかになる。
この三型は論文Bの二根対の
非交差／一根共有／一致に正確に対応する。

具体的な同型類の対応は
定理~\ref{thm:qa-noncubic-rootpair-all-geometric-square-avatar-bijection}
のsquare avatar
$(D_s;F_s,G_s,P)$に、
系~\ref{cor:qa-square-chart}の
幾何的変換$\mathsf T^{-1}$を適用したものである。
$\Sigma_2$側の制限線束を
$F=\OO_{\Sigma_2}(f)|_D$、
$H=\OO_{\Sigma_2}(h)|_D$とすると、
\[
 H\simeq F^2,\qquad
 F_s=F,\quad G_s=FP,\quad
 H_s=F_sG_s=F^2P
 \tag{NC-Sigma2-dict}\label{eq:qa-noncubic-sigma2-square-dictionary}
\]
が全対応の明示公式となる。

これは\emph{体上の集合論的全単射}であり、
$\Sigma_2$ avatarから元のAS代数の
pointed heartを直接回復する定理でも、
Picard line bundleを使う相対スタック同値でもない。
smooth域でこのsquare換算を指定AS10 mutationと一番頂点の
pointingへ接続した対象は、
定理~\ref{thm:qa-spectral-avatar-pointed-involution}により
元の対象の二同種双対に対応する。
元とpointed同値となるのは$\chi=0,2$だけであり、
この対応先は二つのnoncubic境界でH2族にも延長できない
（系~\ref{cor:qa-spectral-involution-no-h2-extension}）。
命題~\ref{prop:qa-rootpair-nodal-picard-nonextension}
の共有根退化では、固定されたspectral曲線族上の
可逆$P$への相対延長が不可能なことに注意する。
\end{theorem}
\begin{proof}
$D\cap S=\varnothing$のとき、
$S=h-2f$の定義から
$\OO_D(S)\simeq\OO_D$となり、
$H\simeq F^2$が従う。
また$D$はどの$\PP^1$-fiberも成分として含まない。
実際、fiberは$S$と交わるので、
そのような成分は$D\cap S=\varnothing$に反する。
従って$p:D\to\PP^1$は有限平坦次数二である。
そのtrace-zero直和分解を取ると
\[
 p_*\OO_D\simeq\OO_{\PP^1}\oplus
                \OO_{\PP^1}(-2)
\]
となる。後項の次数は
$\chi(\OO_D)=0$と$p_a(D)=1$から決まる。
よって$D$は四次式$q$を用いた
$y^2=q$型の二重被覆になる。
$q$が四単根、一つの二重根、二つの二重根なら
前定理の三型になる。
それ以外の四次根配置は
$(3,1)$のcuspidal型、
$(4)$の二枝接触型、
または$q=0$の非被約double sectionである。
これらの$\Pic^0$は標数零で
加法型$\mathbb G_a$となり、
非自明二捩れ点を持たない。
この主張は正規化と特異点での単元の貼合せ、
非被約double sectionでは
$1+\OO_{\PP^1}(-2)$の加法的拡大から得られる。
従って$P\ne\OO_D,\ P^2=\OO_D$は
三つのsmooth/nodal/I2型を強制する。

幾何的変換については
系~\ref{cor:qa-square-chart}が、
$\Sigma_2$の負断面を避けるavatarを
square chartへ送る操作$\mathsf T$を
任意基底変更と可換な同値として構成している。
そこで$H=F^2$を代入すると
square側の線束は
\[
 F_s=F,\qquad G_s=HPF^{-1}=FP
\]
となり、$P=G_sF_s^{-1}$かつ$P^2=\OO_D$である。
逆にこの二捩れsquare wallのデータに
$\mathsf T^{-1}$を適用すれば
$H=F_sG_sP^{-1}=F_s^2$となり、
系~\ref{cor:qa-square-chart}と
定理~\ref{thm:qa-open-charts}により
負断面を避ける$\Sigma_2$に入る。
両操作は互いに逆であり、
曲面・ruling・因子・Picard類の同型を保つ。
これと前定理の全noncubic二根対の
幾何同型類全単射、論文Bの
\cite[\texttt{thm:noncubic-intrinsic-presentation-equivalence}]{QuadricModuli2026}
を合成して結論を得る。
相対射に関する留保は直前の
$A_1$局所非延長の結果による。
\end{proof}

\begin{corollary}[noncubic良moduliの同じ点を区別する二つの$\Sigma_2$ avatar]
\label{cor:qa-sigma2-avatar-separates-special-noncubic-classes}
論文Bのnoncubic境界の良moduli
$\chi:\mathfrak N\to\mathbb A^1$において、
一根共有と二根対一致はともに$\chi=1$へ写るが、
これらは異なるpointed exact同型類である。
定理~\ref{thm:qa-all-noncubic-sigma2-anticanonical-avatar-bijection}
の$\Sigma_2$ avatarによれば、
前者は\emph{既約nodal}反標準因子$D$と
唯一の非零二捩れ$P$、
後者は\emph{二$\PP^1$が二節点で交わる}
反標準因子$D$と唯一の非零二捩れ$P$に対応する。
従って曲線$D$の既約成分構成が、
良moduliでは同一視される二つの
O-marked同型類を区別する。
\end{corollary}
\begin{proof}
良moduliの特殊fiberの二幾何同型類と
$\chi=1$への写像は論文B
\cite[\texttt{prop:noncubic-good-moduli-fibers}]{QuadricModuli2026}
による。
根対の共通根の個数に応じた
二重被覆$C_q$の特異型の分類は
定理~\ref{thm:qa-noncubic-rootpair-quartic-resolvent}
にあり、$\Sigma_2$での再埋込みは
前定理で証明した。
既約性・成分数が異なるので
二avatarは同型になり得ない。
\end{proof}

\begin{remark}[節点型avatarの二つの明示$(2,2)$模型]
\label{rem:qa-singular-square-avatar-explicit-models}
定理~\ref{thm:qa-noncubic-rootpair-all-geometric-square-avatar-bijection}
の特異な二型は、いずれも
$\PP^1_X\times\PP^1_Z$上で明示できる。
既約一節点型では、正規化$\PP^1_t$の
$t=1,-1$を貼り合わせ、非自明二捩れの
貼合せ比を$-1$とする。
$F_s$の二切断は$1,t^2$、
$G_s$の二切断は$t,1-t^2$であり、
\[
 X=t^2,\qquad Z=\frac{1-t^2}{t}
 \quad\Longrightarrow\quad
 XZ^2=(1-X)^2 .
 \tag{Sq-node}\label{eq:qa-nodal-square-avatar-model}
\]
これは$(X,Z)=(1,0)$に普通節点を持つ
既約$(2,2)$曲線である。
式は各$\PP^1$で斉次化するものとする。

一致根対型では、$C_+\cup C_-$を二つの
$\PP^1$の和として、その二交点を
$X=0,\infty$に取る。$F_s$は両成分で
$X$座標、非自明二捩れ$P$による
$G_s$は第一成分で$Z=X$、
第二成分で$Z=-X$となる。
従って
\[
 (Z-X)(Z+X)=0
 \tag{Sq-I2}\label{eq:qa-i2-square-avatar-model}
\]
は二つの$(1,1)$曲線が
$(0,0),(\infty,\infty)$で横断する
反標準$(2,2)$因子を与える。
この二模型は、特異な場合の
$P=-1\in\Pic^0(D)[2]$が本当に
square avatarの二rulingを作ることを
直接に示している。
\end{remark}

\begin{proposition}[Picard零切断の全体を古典pointとして追加できない理由]
\label{prop:qa-avatar-zero-picard-collapse}
$Y=\PP^1\times\PP^1$、
$D\in|\OO_Y(2,2)|$を任意の反標準Cartier因子とし、
$P=\OO_D$と置く。
このとき曲線上の四線束
\[
 (\OO_D,\ F_s,\ G_s,\ F_sG_s)
\]
の有向End代数（各積を含む）は、
$D$の選択に依存せず、
$Y$上の可換quadricの標準四項collection
\[
 (\OO_Y,\OO_Y(1,0),\OO_Y(0,1),\OO_Y(1,1))
\]
の有向End代数と同型である。
従って$\Pic^0$の零切断をそのまま
曲線型avatarに追加すると、
異なる$D$が同一の可換foundation代数を与え、
O-marked分類への単射性が失われる。
古典的な可換quadricは、曲線型とは別に
一点として追加する必要がある。
\end{proposition}
\begin{proof}
相異なる四項の間のHomに現れる
$\OO_Y(a,b)$は、
$a,b\in\{-1,0,1\}$の範囲にある。
各々についてCartier制限列
\[
 0\to\OO_Y(a-2,b-2)\to
 \OO_Y(a,b)\to\OO_D(a,b)\to0
\]
を考えると、$a-2\le-1$、$b-2\le-1$であり、
Künnethの公式から
$H^0(Y,\OO(a-2,b-2))=
H^1(Y,\OO(a-2,b-2))=0$である。
従って四項間の全Homの制限は同型となる。
構成は線束の合成に関して自然なので、
有向End代数の積も保存する。
右辺は$D$に依存しない。
\end{proof}

\begin{remark}[二つの $\Sigma _2$ への移行]
\label{rem:qa-rotated-square}
square helix の foundation を一つ進めた rotated square の偏極は
\[
 F_r=F_s,\qquad H_r=F_sG_sP
\]
である（本稿の幾何的証明は定理~\ref{thm:qa-core-avatar-square-rotated}；
原草稿\cite[命題11.1・定理15.1]{AvatarVSeven2026}とも整合する）。
一方、\eqref{eq:qa-square-dictionary}の偏極は $H=F_sG_sP^{-1}$ である。
前者の $\Sigma _2$ 条件は $P=F_sG_s^{-1}$、後者の条件は $P=G_sF_s^{-1}$ となる。
これらは一般には別の壁である。
foundation の移動や mutation とともに選ばれる構造対象も変わるため、
固定された構造対象の内在的一点 Hilbert 関手を同じものとして扱わない。
\end{remark}

\subsection{古典点と最深境界の扱い}

本節の stack は一貫して $P$ が各 fiber で非自明な曲線型を扱った。
体上の分類では、$P=\OO_D$ とした有限 directed algebra は $D$ の選択に依存せず、
\[
 \mathrm{Cl}_0:\operatorname{coh}(\PP^1\times\PP^1),\qquad
 \mathrm{Cl}_2:\operatorname{coh}(\PP(1,1,2))
\]
という二つの古典型に対応する
\cite[注意11.3・定理15.3]{Q12Avatar2026}。
後者の $\PP(1,1,2)$ は重み付き射影 stack である。
これは同型類の説明であり、相対 Picard の零切断を二点へ自動的に縮約する
stack の構成ではない。零切断では $D$ とその自己同型が残るため、
相対的に古典点を加えるには別途の収縮または同一視を定義する必要がある。

最深の十二本矢層についても、三頂点の有限安定表現のモジュライが $\Sigma _2$ になることと、
標準 heart における内在的一点 Hilbert 曲面が $\Sigma _2$ になることは異なる。
この区別のための内在的な検査として，
例~\ref{ex:qa-weighted-not-h2}では重み付き古典模型の全標準頂点で，
例~\ref{ex:qa-hirzebruch-not-h2}では $\coh\Sigma_2$ 自身の全可逆層で，
反Serre軌道の生成条件が破れることを示した。
古典型の負断面上には、一つの foundation では安定であっても、
次の foundation では最後の頂点から生成されなくなる表現が存在する
\cite[命題14.1]{Q12Avatar2026}。
したがって本節の開部分 stack だけから、全 $\mathcal A_{12}^{\times}$ の
内在的 quadric モジュライへの射、あるいはその射の定義域が開であることを結論しない。
そこでは構造対象、標準 heart、豊富性、Serre 平方根および相対 AS 持上げを
別に検証しなければならない。

% !TeX root = ../nc-quadric-avatar.tex
% 2026-11-22 / two genuinely distinct genus-one curves in noncubic comparison
\section{Noncubic の点曲線と spectral 曲線の二重 isogeny}
\label{sec:qa-point-spectral-isogeny}

noncubic の中央二重平方環から得る可換点曲線と，
二根対の四次 resolvent から得る spectral 曲線は，
滑らかな場合でさえ一般に同型ではない．
両者の関係を明示することは，avatar の幾何的分類を
内在的 O-marked heart の同値と取り違えないために必要である．

\begin{theorem}[点曲線とspectral曲線の二重 isogeny]
\label{thm:qa-point-spectral-dual-two-isogeny}
$k$ を標数 $0$ の代数閉体，
$f_1,f_2\in k[u,v]_2$ を squarefree かつ互いに
根を共有しない二次形式とする．
\[
 \begin{split}
 E_f&=\{c^2=f_1(u,v),\ d^2=f_2(u,v)\}
                   \ \subset\PP^3_{[u:v:c:d]},\\
 C_q&=\{y^2=f_1(u,v)f_2(u,v)\}
             \ \longrightarrow\PP^1_{[u:v]}
 \end{split}
\]
と置く．$E_f$ は
定理~\cite[\texttt{thm:ore-anticommuting-central}]{QuadricIntrinsic2026}
に現れる中央二重平方環の可換点曲線であり，
$C_q$ は
式~\eqref{eq:qa-rootpair-spectral-curve}の
spectral 曲線である．
次が成り立つ．
\begin{enumerate}
\item $E_f$ と $C_q$ は滑らかな連結種数一曲線である．
 同時符号対合
 $\tau(u,v,c,d)=(u,v,-c,-d)$ は $E_f$ に自由に作用し，
 \[
   r:E_f\longrightarrow C_q,\qquad
   [u:v:c:d]\longmapsto([u:v],y=cd)
   \tag{Point--spectral}\label{eq:qa-point-spectral-cover}
 \]
 はその有限 \'{e}tale 二重商である．
 原点を選べば $r$ は次数二の isogeny となる．
\item $\alpha(c,d)=(-c,d)$ と $\beta(c,d)=(c,-d)$ とする．
 これらの商は滑らかな conic
 \[
    Q_d=\{d^2=f_2(u,v)\},\qquad
    Q_c=\{c^2=f_1(u,v)\}
 \]
 であり，商射を $p_d:E_f\to Q_d$，
 $p_c:E_f\to Q_c$ と書く．
 $Q_d,Q_c\simeq\PP^1$ 上で各々唯一の次数一線束
 $N_d,N_c$ を取り，
 $F_d=p_d^*N_d,\ F_c=p_c^*N_c$ とすると
 \[
   \deg F_d=\deg F_c=2,\quad
   F_d^{\otimes2}\simeq F_c^{\otimes2}
       \simeq\OO_{E_f}(1),\quad
   0\ne P_E:=F_cF_d^{-1}\in\Pic^0(E_f)[2].
   \tag{Point-avatar}\label{eq:qa-point-curve-double-wall}
 \]
 従って $(E_f,F_d,F_c,P_E)$ 自身が
 二重壁上の square avatar を与える．
 また $\tau$ は $P_E$ による非自明二捩れ平行移動である．
\item
 $M_C=(r_*\OO_{E_f})^-$ は
 定理~\ref{thm:qa-noncubic-rootpair-torsionfree-spectral-module}
 の非自明二捩れ線束である．
 その $\Pic^0(C_q)[2]$ の点は，
 $r$ の\emph{双対 isogeny}の核を生成する：
 \[
     0\longrightarrow
       \{\OO_{C_q},M_C\}
       \longrightarrow\Pic^0(C_q)
       \xrightarrow{\ r^*\ }\Pic^0(E_f)
       \longrightarrow0.
       \tag{Dual-2}\label{eq:qa-point-spectral-dual-isogeny}
 \]
 特に原点の選択の下で
 $E_f\simeq C_q/\langle M_C\rangle$ である．
\end{enumerate}
\end{theorem}
\begin{proof}
$E_f\to\PP^1$ は $c,d$ の符号を独立に変える
biquadratic な次数四被覆である．
$f_1,f_2$ の四分岐点が互いに異なるため，
各分岐点では一方の平方根だけが分岐し，
$E_f$ は滑らかで連結である．
Riemann--Hurwitz により
\[
 2g(E_f)-2=4(-2)+4\cdot2=0.
\]
同時符号対合の固定点では$c=d=0$を要するが，
これは$f_1=f_2=0$を意味して仮定に反する．
この対合の不変関数は
$u,v$ と $cd=y$ で，$y^2=f_1f_2$であるから，
商は$C_q$である．
商は非分岐二重被覆で，$C_q$も種数一となる．

$\alpha$ と $\beta$ の不変式から，それぞれの
商が $Q_d,Q_c$ であることが直ちに従う．
squarefree 性によって conic は滑らかであり，
$k$ は代数閉だから $\PP^1$ に同型である．
Conic の超平面線束は $\OO_{\PP^1}(2)$なので，
その平方根となる次数一線束は一意に存在する．
$E_f$ の埋込の線形座標
$(u,v,d)$，$(u,v,c)$ はそれぞれ
$Q_d$，$Q_c$ の超平面切断の引戻しだから
$F_d^2=F_c^2=\OO_{E_f}(1)$となる．
それぞれ次数二で大域生成されており，
相異なる二つの pencil は
補題~\ref{lem:qa-core-two-pencils}により
$E_f$を$\PP^1\times\PP^1$の$(2,2)$因子として
閉埋め込みする．
もし $F_d\simeq F_c$ なら二つの次数二線形系は
基底の射影変換を除いて同じ射を定め，
その deck involution も一致するはずである．
しかし $\alpha\ne\beta$ なので非自明性を得る．

滑らかな種数一曲線の次数二線束 $L$ は，
原点を選んで$x\mapsto \ell_L-x$という
involution を定める．
従って二線束$F_d,F_c$の involution
$\alpha,\beta$の積は，その差
$F_cF_d^{-1}=P_E$による平行移動である
（符号の違いは二捩れゆえ消える）．
これで第二項が従う．

最後に$E_f\to C_q$は \'{e}tale な
$\mu_2$-torsorであり，
$r_*\OO_{E_f}=\OO_{C_q}\oplus M_C$，
$M_C^{\otimes2}\simeq\OO_{C_q}$である．
$H^0(C_q,M_C)=0$は
式~\eqref{eq:qa-rootpair-spectral-module}から従う．
従って $M_C$ は非自明である．
有限 \'{e}tale 同種写像の双対性，
または線束の Galois descent から，
$\ker r^*=\{\OO_{C_q},M_C\}$となる．
次数二の isogeny の双対も次数二なので
$r^*$ は全射となり，最後の完全列を得る．
\end{proof}



\begin{proposition}[非自由MCM二型の行列核は点曲線の二pencilを回収する]
\label{prop:qa-point-mcm-matrix-pencils}
前定理の仮定の下で，$f_1=-\ell_1\ell_2$，
$f_2=-m_1m_2$と因数分解し，
$\ell_1,\ell_2$および$m_1,m_2$を
それぞれ独立な一次形式とする．
中央二重平方環 $B_f$ で
\[
 M_c=\begin{pmatrix}c&\ell_1\\\ell_2&-c\end{pmatrix},
 \qquad
 M_d=\begin{pmatrix}d&m_1\\m_2&-d\end{pmatrix}
 \tag{MCM-matrices}\label{eq:qa-point-mcm-matrices}
\]
と置くと，$M_c^2=M_d^2=0$であり，
それぞれは元のnoncubic normal helixに使う
二つの線形matrix factorizationの表示となる．
$E_f$上では両行列が全点で階数一となり，
その右核の直線から定まる射
\[
 \kappa_c:E_f\longrightarrow\PP^1,\qquad
 \kappa_d:E_f\longrightarrow\PP^1
\]
は，それぞれ前定理の
$p_c:E_f\to Q_c$，$p_d:E_f\to Q_d$
と conic の同型 $Q_c,Q_d\simeq\PP^1$
との合成に一致する（終域の射影変換を除く）．
従って
\[
 \kappa_c^*\OO_{\PP^1}(1)\simeq F_c,\qquad
 \kappa_d^*\OO_{\PP^1}(1)\simeq F_d .
 \tag{MCM-pencils}\label{eq:qa-point-mcm-pencil-classes}
\]
これは\emph{matrix factorizationの点上のkernel}と
幾何的な二pencilの一致であり，
ASのfinite directed End代数全体の同型や
$\pi B_f$を保つpointed心の同型を
まだ結論しない．
\end{proposition}
\begin{proof}
$u,v$ は $B_f$ の中心にあり，
$c^2=f_1=-\ell_1\ell_2$，
$d^2=f_2=-m_1m_2$だから
行列の平方はそれぞれ零である．
これは
\cite[\texttt{thm:ore-anticommuting-central}]{QuadricIntrinsic2026}
の二つの非自由MCM型に用いる
線形行列分解の表示である．
$E_f$の点において$M_c$が零行列となるなら
$c=\ell_1=\ell_2=0$が必要だが，
$\ell_1,\ell_2$は$\PP^1_{[u:v]}$上で
共通零点を持たない．
従って行列の階数は一である．
同じ議論が$M_d$にも適用できる．

$M_c$の成分は$u,v,c$だけに依存するので，
$\kappa_c$は conic $Q_c$ を経由する．
$\ell_1,\ell_2$を $s,t$ とする座標で
$Q_c$の式は $c^2+st=0$である．
その標準パラメータ表示
\[
 [s:t:c]=[a^2:-b^2:ab]
\]
の下で
\[
 M_c=
 \begin{pmatrix}a\\-b\end{pmatrix}
 \begin{pmatrix}b&a\end{pmatrix}
\]
となるから，右核は直線$[a:-b]$である．
これは conicのパラメータ$[a:b]$への
射影線形自己同型であり，
$\kappa_c$は$Q_c\simeq\PP^1$を通る次数二射となる．
よって$\kappa_c^*\OO(1)=F_c$を得る．
$m_1,m_2,d$を用いた同じ計算で
$F_d$の主張も従う．
\end{proof}

\begin{proposition}[MCMの隣接積と二pencilのVeronese積の一致]
\label{prop:qa-point-mcm-composition-veronese}
前命題の仮定の下で，中央二重平方環の
一次部分を
\[
 V=(B_f)_1=\langle u,v,c,d\rangle_k
     \xrightarrow{\ \sim\ }
     H^0(E_f,\OO_{E_f}(1))
 \tag{Degree-one}\label{eq:qa-point-degree-one-restriction}
\]
と同定する．
二つの非自由MCM型 $X_c,X_d$ に関する
noncubic normal helix の隣接合成
\[
 \begin{aligned}
 \mu_c:\;&
 \Hom(\pi X_c,\pi B_f(1))
 \otimes\Hom(\pi B_f,\pi X_c)
 \longrightarrow\Hom(\pi B_f,\pi B_f(1)),\\
 \mu_d:\;&
 \Hom(\pi X_d,\pi B_f(1))
 \otimes\Hom(\pi B_f,\pi X_d)
 \longrightarrow\Hom(\pi B_f,\pi B_f(1))
 \end{aligned}
\]
の像をそれぞれ $U_c,U_d\subset V$ とする．
一方，$E_f$ 上の二つの次数二pencilに対する
通常の対称積
\[
 m_c:\Sym^2H^0(E_f,F_c)\to H^0(E_f,F_c^2),
 \qquad
 m_d:\Sym^2H^0(E_f,F_d)\to H^0(E_f,F_d^2)
\]
の像を前式の$V$へ写す．すると
\[
 \boxed{
 \im\mu_c=\im m_c=U_c=\langle u,v,c\rangle,\qquad
 \im\mu_d=\im m_d=U_d=\langle u,v,d\rangle.}
 \tag{Two-rank3}\label{eq:qa-point-two-rank-three-subspaces}
\]
特に
\[
 \dim U_c=\dim U_d=3,\qquad
 U_c\cap U_d=\langle u,v\rangle,\qquad U_c+U_d=V.
\]
よってnoncubic normal helixの二つのMCM型に伴う
\emph{最小十本矢の隣接合成像}は，点曲線$E_f$の
二つのconic商pencilのVeronese像と一致する．
ただし，この命題は二つの像部分空間と
MCMのkernel写像を同定するもので，
四頂点のfinite directed End代数の
\emph{すべての積}の同型を意味しない．
\end{proposition}
\begin{proof}
$E_f\subset\PP^3$は二quadricの完全交叉で，
一次の斉次関係はない．
よって$V=H^0(E_f,\OO_{E_f}(1))$となる．
定理
~\cite[\texttt{thm:normal-helix} および
 \texttt{cor:normal-cubic-defect}]{QuadricIntrinsic2026}
における非自由MCMの隣接合成は，
線形matrix factorizationの
$2\times2$成分を一次切断とみなしたものだから，
式~\eqref{eq:qa-point-mcm-matrices}より
\[
 \im\mu_c=\langle c,\ell_1,\ell_2\rangle
          =\langle u,v,c\rangle,\qquad
 \im\mu_d=\langle d,m_1,m_2\rangle
          =\langle u,v,d\rangle
\]
である．ここで各一次形式対は独立である．

一方，前命題から
$F_c=p_c^*N_c$，$N_c\simeq\OO_{\PP^1}(1)$であり，
$H^0(E_f,F_c)=H^0(Q_c,N_c)$である．
実際，$p_{c*}\OO_{E_f}$の不変部分は$\OO_{Q_c}$で，
反不変部分の負次数の線束には切断がない．
従って$m_c$は
\[
 \Sym^2H^0(Q_c,N_c)\longrightarrow
 H^0(Q_c,N_c^2)=H^0(Q_c,\OO_{Q_c}(1))
\]
の通常のVeronese積の引戻しである．
この積は三次元切断空間への同型であり，
$Q_c\subset\PP^2_{[u:v:c]}$の
超平面切断の像は
$\langle u,v,c\rangle$である．
同様に$m_d$の像は$\langle u,v,d\rangle$となる．
二つの部分空間の共通部分と和は明らかである．
\end{proof}

\begin{lemma}[noncubic 点曲線への制限と四頂点Endの同型]
\label{lem:qa-pointed-abelianization-four-end}
標数零の代数閉体上、$f_1,f_2$がそれぞれ相異なる二根を持ち、
互いに根を共有しないとする。
$B=B_f$、$R=B/(cd-dc)$、$E=\operatorname{Proj}R=E_f$とする。
$f_1=-\ell_1\ell_2$, $f_2=-m_1m_2$に対する
二つの右$B$-加群$X_c=\operatorname{coker}M_c$,
$X_d=\operatorname{coker}M_d$を前命題のように取る。
$F=F_c$, $G=F_d$と書けば、次の自然な制限は
頂点と全合成を保つ有限次元代数の同型である：
\begin{equation}
\End_{\rm dir}^{\qgr B}
 (\pi B,\pi X_c,\pi B(1),\pi X_d(1))
 \ \xrightarrow{\ \sim\ }\
\End_{\rm dir}^{E}
 (\OO_E,F,F^2,F^2G).
\label{eq:qa-pointed-four-end-abelianization}
\end{equation}
ここで$F^2=G^2=\OO_E(1)$であり、
左辺では$B$の通常シフト$(1)$を記すが、
これ自体を適合Serre平方根とは主張しない。
\end{lemma}
\begin{proof}
$R_c=k[u,v,c]/(c^2-f_1)$とおく。
$R=R_c\oplus dR_c(-1)$は$R_c$上自由である。
滑らかなconic $Q_c=\operatorname{Proj}R_c\simeq\PP^1$
上で$M_c$は$s=\ell_1,t=\ell_2$により
\[
c^2+st=0,\qquad
[s:t:c]=[a^2:-b^2:ab],\qquad
M_c=\binom{a}{-b}(b\ \ a)
\]
となる。従って$\operatorname{coker}M_c$を
$R_c$へ制限した次数付き加群は
$\bigoplus_{n\ge0}H^0(Q_c,\OO_{\PP^1}(2n+1))$である。
$R_c\to R$への自由な拡大と
$p_{c*}\OO_E\simeq\OO_{\PP^1}\oplus\OO_{\PP^1}(-2)$から
\[
X_c\otimes_B R
 \simeq\bigoplus_{n\ge0}H^0(E,F\OO_E(n)),
\qquad
X_d\otimes_B R
 \simeq\bigoplus_{n\ge0}H^0(E,G\OO_E(n)).
\tag{Spinor--sections}\label{eq:qa-pointed-spinor-sections}
\]
両式は単にsheafificationの同型ではなく、
非負次数のsection moduleの同型である。
特に右辺は有限長torsionを持たない。
$B\to R$の核は次数二の$cd-dc$で生成されるので、
$B_0,B_1$ならびに$(X_c)_0,(X_c)_1,
(X_d)_0,(X_d)_1$は、それぞれ$R$への制限で不変である。

Aの適合表示をここで明確にする。
分類論文Aの
\cite[\texttt{thm:ore-anticommuting-central}]{QuadricIntrinsic2026}
で構成するZhang自己同型$\sigma$は、
線形座標変換を伴って$c,d$を交換し、
$B^\sigma$の通常シフトを適合根とする。
右Zhang同値でそのnormal helixの四連続対象を$B$側へ
戻すと、自由加群のshiftの作用を$\sigma$で輸送し、
二つのspinor型も交換するので、
添字を選んで
$(\pi B,\pi X_c,\pi B(1),\pi X_d(1))$
を得る。この輸送は次数付き加群の単なる同一視ではなく、
右作用の$\sigma$-twistを含む。
従って
\cite[\texttt{thm:normal-helix},
\texttt{cor:normal-cubic-defect}]{QuadricIntrinsic2026}
のHom計算と、深さ三によるgraded/tails比較が適用できる。
四頂点の上三角Hom次元は
\[
\begin{pmatrix}
1&2&4&6\\0&1&2&4\\0&0&1&2\\0&0&0&1
\end{pmatrix}.
\tag{Four-End-dims}\label{eq:qa-pointed-four-end-dimensions}
\]
右辺の曲線上の線束も次数$0,2,4,6$であり、
$E$が滑らかな種数一曲線なので同じ次元を持つ。

左辺の各Homは生成元の像で決まり、
$B_0,B_1,(X_d)_0,(X_d)_1$のいずれかに値を持つ。
特に$\Hom(X_c,X_d(1))$の射は
$X_c$の次数零の二生成元の$(X_d)_1$内の像で決まる。
上記の低次数の制限単射性と
\eqref{eq:qa-pointed-spinor-sections}により、
どのHomも右辺のsheaf Homへ単射となる。
両辺の次元が一致するため全て同型である。
$-\otimes_B R$は右完全で、個々の加群射の合成を保つため、
上の六つの同型は同時に全積を保つ。

最初の二段階合成はさらに直接確認できる。
$X_c$の次数零生成元を$e_1,e_2$とし、
$\Hom_B(X_c,B(1))$の基底を
$(c,\ell_1),(-\ell_2,c)$と取ると、
合成の$2\times2$行列は
\[
\begin{pmatrix}
c&\ell_1\\-\ell_2&c
\end{pmatrix}.
\]
従って唯一の二次関係は二つの対角道の差であり、
像は$\langle c,\ell_1,\ell_2\rangle$である。
上のconic表示で行列は
$\binom{a}{b}(b\ \ a)$となり、
$H^0(F)\otimes H^0(F)\to H^0(F^2)$の
反対称核と一致する。この計算は像の一致だけでなく、
二つの二次元Homの\emph{双線形積そのもの}を同定する。
\end{proof}

\begin{theorem}[滑らかなnoncubicの位相を合わせたpointed認識]
\label{thm:qa-point-smooth-noncubic-pointed}
\label{prop:qa-point-mcm-pointed-finite-end-reduction}
前補題の仮定の下で、
square avatar $(E,F,G,P)$を
$F=F_c$, $G=F_d$, $P=GF^{-1}\ne\OO_E$で定める。
$P^2=\OO_E$かつ$G=FP$である。
このsquare型AS$\mathbb Z$-代数を論文Cの指定mutationで
最小十本矢型$A_{\rm av}$に移し、
mutation後のhelixを$(H_i)_{i\in\mathbb Z}$とする。
このとき
\begin{equation}
 \End_{\rm dir}^{\qgr B_f}(E_0,E_1,E_2,E_3)
 \ \simeq\
 \End_{\rm dir}^{\qgr A_{\rm av}}(H_1,H_2,H_3,H_4),
\qquad E_0=\pi B_f,
\label{eq:qa-pointed-shifted-four-end}
\end{equation}
という頂点と全合成を保つ同型が存在する。
論文Aのnormal-helixと本稿のAS10 Serre完成・Noether性・AZ再構成を
適用すると、
\begin{equation}
(\qgr B_f,\pi B_f)\simeq
(\qgr A_{\rm av},\,\pi e_1A_{\rm av})
\tag{Smooth-pointed}\label{eq:qa-pointed-smooth-equivalence}
\end{equation}
というpointed exact同値を得る。
ここで構造対象の像はmutation\emph{後の一番頂点}であり、
ゼロ番頂点ではない。
\end{theorem}
\begin{proof}
square側の曲線上の初期四項は$(P,F,G,FG)$である。
補題~\ref{lem:qa-as10-all-index-mutual-inverse}で扱うevaluation mutationでは
\[
 H_0=K_0,\quad H_1=P,\quad H_2=G,\quad H_3=FG,
\quad
 K_0=\ker\bigl(H^0(FP^{-1})\otimes P\to F\bigr)
      =P^2F^{-1}.
\]
次のsquare項は
系~\ref{cor:qa-core-square-complexes}の式から
$L_4=F^2G^2P,\ L_5=F^3G^2P^2$である。
ゆえに一周期後のevaluation kernelは
\[
 H_4=K_1=L_4^2L_5^{-1}=FG^2.
\]
$P^2=\OO_E$, $G=FP$を代入し、
四項$(H_1,H_2,H_3,H_4)=(P,FP,F^2P,F^3)$を
共通線束$P^{-1}$でtensorすると
\[
(\OO_E,F,F^2,F^3P^{-1})
   =(\OO_E,F,F^2,F^2G).
\]
共通線束のtensorは曲線上の全Homとその積を保つ。
前補題と本稿の全整数mutationのEnd同定を合わせれば
\eqref{eq:qa-pointed-shifted-four-end}を得る。
これは線束上の計算であり、任意のnc heartに
線束tensor自己同値を仮定するものではない。

Aのnormal helixはfullでample、
Cのmutation後のhelixもfull very strongで、
Noether性とAZ amplenessを持つ。
有限Endの同型はtilting Morita同値を与え、
Serre関手と自然に可換する。
従って四項をSerre完成して得られる全整数の
前向きHomと積は一致する。
補題~\ref{lem:qa-as-serre-unique-recovery}と
補題~\ref{lem:qa-as10-all-index-mutual-inverse}および
\cite[\texttt{thm:normal-helix}]{QuadricIntrinsic2026}
の再構成を適用して
\eqref{eq:qa-pointed-smooth-equivalence}を得る。
\end{proof}

\begin{remark}[位相の取り違えと相対境界]
\label{rem:qa-pointed-phase-obstruction}
同じnoncubic normal四項をmutation後の
$(H_0,H_1,H_2,H_3)$に\emph{頂点固定で}
同定することはできない。
前者の$0\to1\to2$合成の核は一次元だが、
後者は最小十本矢型の定義により
$0\to2$に最小関係を持たず、核は零である。
従って一添字shiftは単なる記法でなく必要条件である。
なお定理は四根相異なるsmooth noncubicの
\emph{体上のpointed同値}を述べる。
一根共有・二根対一致における同じ式の相対化、
非被約底上のgroupoid、$C_q$側の通常Picardへの延長、
およびpotential全体の相対忘却は主張しない。
\end{remark}

\begin{corollary}[非零二捩れ付き楕円曲線上の双対二重商]
\label{cor:qa-point-spectral-atkin-lehner}
滑らかな楕円曲線と非零二捩れの組
$(T,L)$ の同型類に対し
\[
 \mathcal W_2(T,L)=
 \bigl(T/\langle L\rangle,\ L^\vee\bigr)
 \tag{Dual-wall}\label{eq:qa-point-spectral-dual-wall}
\]
と置く．ここで $L$ を原点選択によって
位数二の点とみなし，$L^\vee$は
商 isogeny の双対核を指定する非零二捩れである．
この対応は同型類上の involution：
$\mathcal W_2^2=\mathrm{id}$ を定める．
前定理の根対から生じる二つの組は
\[
 \mathcal W_2(C_q,M_C)=(E_f,P_E),\qquad
 \mathcal W_2(E_f,P_E)=(C_q,M_C)
\]
を満たす．
これは楕円曲線と二捩れ\emph{だけ}の対応であり，
$F_d,F_c$と$p^*\OO_{\PP^1}(1)$という
異なる degree-two pencil まで
自動的に同一視するものではない．
\end{corollary}
\begin{proof}
有限 \'{e}tale 二重 isogeny
$T\to T/\langle L\rangle$に対し，
双対 isogeny の核は位数二である．
さらに双対を取れば，楕円曲線の主偏極による
二重双対の自然同型から元の isogeny に戻る．
従って商と双対核の操作は二回で元に戻る．
前定理で$r:E_f\to C_q$の核は$P_E$であり，
$r^\vee:C_q\to E_f$の核は$M_C$と示したので，
二つの表示が得られる．
\end{proof}

\begin{example}[二曲線は一般に同型でない]
\label{ex:qa-point-spectral-j-distinction}
アフィン座標$x=u/v$で
$f_1=x(x-1)$，$f_2=(x-2)(x-3)$とする．
$C_q$の射影模型は
\[
 y^2=x(x-1)(x-2)(x-3).
\]
一方 $E_f$の関数体で
$t=d/(x-3)$と置けば
\[
 x=\frac{3t^2-2}{t^2-1},\qquad
 z=(t^2-1)c,\qquad
 z^2=(3t^2-2)(2t^2-1)=6t^4-7t^2+2.
\]
四次式 $az^4+bz^3+cz^2+dz+e$ の
不変式を $I=12ae-3bd+c^2$，
$J=72ace+9bcd-27ad^2-27b^2e-2c^3$
とすると
$j=1728\cdot4I^3/(4I^3-J^2)$である．
$C_q$の四次式には $(I,J)=(13,-70)$，
$E_f$の上の四次式には $(I,J)=(193,-5362)$を得るので，
\[
 \boxed{j(C_q)=\frac{35152}{9},\qquad
        j(E_f)=\frac{28756228}{3}.}
 \tag{Distinct-j}\label{eq:qa-point-spectral-distinct-j}
\]
従って両曲線は同型でない．
この有理数の計算は
\texttt{computations/verify\_rootpair\_resolvent\_spectral.py}
にも厳密算術の検算を置く．
\end{example}

\begin{remark}[pointed 認識への帰結と退化境界]
\label{rem:qa-point-spectral-pointed-warning}
定理~\ref{thm:qa-noncubic-rootpair-all-geometric-square-avatar-bijection}
で使った square avatar の曲線は \emph{$C_q$} である．
一方，noncubic 中央二重平方環そのものの
可換点曲線は \emph{$E_f$} であり，
第二項は $E_f$ 上にも自然な
二捩れ double-wall avatar があることを示す．
式~\eqref{eq:qa-point-spectral-distinct-j}により，
二つの構成を同一曲線とみなして
avatar由来AS代数と$\qgr B_f$の
pointed同値を結論することはできない．
\emph{どちらの曲線が特定のnormal helixの
AS avatarとして抽出されるか}，
およびそのexceptional foundationの対応は
まだ直接同定していない．これは次の認識問題である．

$f_1,f_2$ が一根を共有すると，
$\tau$には固定点が生じ，
$r$は有限 \'{e}tale 二重isogenyではなくなる．
その際に正しい相対対象は
式~\eqref{eq:qa-rootpair-spectral-module}の
torsion-free $M_C$とその二次積であり，
双対isogenyの短完全列を特異fiberへ
そのまま延長してはならない．
\end{remark}

% !TeX root = ../nc-quadric-avatar.tex
% 2026-11-22: 共有根 A1 の最小解消に残る相対 Pic^0 障害
\section{共有根の最小解消と相対 Picard の奇偶障害}
\label{sec:qa-nodal-resolution-parity}

定理~\ref{thm:qa-noncubic-rootpair-torsionfree-spectral-module}
の相対 spectral module は共有根で非局所自由となる．
この非自由性は，その点を最小解消して線束にすれば
自動的に通常の相対 $\Pic^0$ に戻る，という意味ではない．
以下は元の一径数基底を変えない最小解消での具体的な障害である．

\begin{proposition}[最小解消後にも残る奇数次数]
\label{prop:qa-nodal-resolution-odd-exceptional-degree}
$k$ は標数零の代数閉体とし，
$A=k[[a,b,y]]/(ab-y^2)$，
$s=a-b$ とおく．
符号商
\[
 A=k[[c,d]]^{\{\pm1\}},\qquad
 a=c^2,\quad b=d^2,\quad y=cd
\]
の反不変部分
$M=k[[c,d]]^- =Ac+Ad$
を考える．$\pi:\widetilde X\to\Spec A$ を
$A_1$ 特異点の最小解消，$E\simeq\PP^1$ を
例外曲線とする．
\[
 L=(\pi^*M)/(\text{torsion})
\]
は $\widetilde X$ 上可逆で，その例外曲線への制限は
\[
                    L|_E\simeq\OO_{\PP^1}(1).
\]
特に任意の $n\in\mathbb Z$ に対し
$\deg_E L(nE)=1-2n$ は奇数であり，
相対 $\Pic^0$ に属する多重次数零の線束にはならない．

さらに，$s$ を基底パラメータとする適切な大域曲線族で
特殊曲線が既約一節点であれば，最小解消後の特殊
fiber は二つの成分 $D'$ と $E$ を持ち，
$D'\cdot E=2$ である．
同じ generic fiber 上で $M$ と一致する
任意の可逆延長は，$L$ から垂直 Cartier 因子で
tensor したものであるから，
$E$ 上の次数の奇偶は変わらない．
従って\emph{元の一径数基底上で最小解消するだけでは}，
非自明二捩れの generic fiber を
多重次数零の可逆線束の相対族として延長できない．
\end{proposition}
\begin{proof}
最小解消の標準二枚のアフィン chart は
\[
\begin{array}{c|ccc}
 &a&y&b\\ \hline
 U_c&u&ut&ut^2\\
 U_d&vw^2&vw&v
\end{array},\qquad
v=ut^2,\quad w=t^{-1}.
\]
例外曲線は $u=0$ と $v=0$ を貼り合わせた
$\PP^1_{[c:d]}$ で，自己交点は $E^2=-2$ である．
$M$ の生成元を形式的に $e_c=c,\ e_d=d$ と書く．
$A$ 上の関係 $yc=ad$ と $yd=bc$ を引き戻して
torsion を除くと，
$U_c$ では $e_d=t e_c$，
$U_d$ では $e_c=w e_d$ となる．
従って各 chart では一つの自由生成元を持ち，
重なりでの遷移関数は
$e_d=t e_c$ である．これは
$\OO_{\PP^1}(1)$ の標準遷移である．
$\OO(E)|_E=\OO_E(-2)$ だから
$\deg_E L(nE)=1-2n$ が従う．

$s=a-b$ を引き戻すと
$U_c$ では $s=u(1-t^2)$，
$U_d$ では $s=v(w^2-1)$ となる．
従って特殊 fiber には例外成分 $E$ が重複度一で現れ，
元の節点の二枝の strict transform が
$E$ の $t=1,-1$ の二点に横断的に交わる．
大域的に元の特殊曲線が既約一節点なら，
この二枝は同じ既約 strict transform $D'$ に属する．
ゆえに $D'\cdot E=2$ である．
二つの可逆延長が generic fiber 上同型なら，
正則曲面上のその比は垂直 Cartier 因子で表される．
その成分は $E$ と $D'$ の整数係数の和なので，
$E$ 上の次数の変化は
$nE^2+mD'\cdot E=-2n+2m$ で偶数である．
よって多重次数零への延長は存在しない．
\end{proof}

\begin{corollary}[相対 avatar の修正に必要な追加操作]
\label{cor:qa-nodal-resolution-requires-more}
共有根の noncubic spectral avatar を
通常の相対 $\Pic^0$ へ移すには，
元の基底を変えずに全空間の $A_1$ を最小解消するだけでは
不十分である．
基底変更，非可逆 rank-one torsion-free 層を許す
compactified Picard，または安定モデルと偏極を
同時に変更する構成のいずれかを別途検討する必要がある．
この命題は，それらの追加操作を施した場合にも
相対射が存在しないと主張するものではない．
\end{corollary}
\begin{proof}
命題~\ref{prop:qa-nodal-resolution-odd-exceptional-degree}
で可逆延長の例外成分上の次数が奇数になる．
相対 $\Pic^0$ の多重次数零の線束にはなり得ない．
\end{proof}


\subsection{Spectral 二捩れ乗法の欠損と正規化上の極限}
\label{subsec:qa-spectral-compactified-limit}

前の奇偶障害は，線束の極限が存在しないという意味ではなく，
\emph{可逆線束に限った極限が存在しない}ことを意味する．
相対 torsion-free sheaf 自身には自然な二次積があり，
その余核によって境界を scheme-theoretic に測定できる．
以下では定理~\ref{thm:qa-noncubic-rootpair-torsionfree-spectral-module}
の記号を使い，二形式は各々二つの相異なる根を持つとする．

\begin{proposition}[乗法欠損の長さと自己双対性]
\label{prop:qa-spectral-product-defect}
$B=\PP^1_k$，$f_1,f_2\in H^0(B,\OO_B(2))$ は
非零で各々 squarefree，$k$ は標数零の代数閉体とする．
$C=\{y^2=f_1f_2\}$ と
$\mathcal M=(r_*\OO_{\widetilde C})^-$ を取る．
反不変元の積から生じる
\[
 \mu:\mathcal M\otimes_{\OO_C}\mathcal M\longrightarrow\OO_C
\]
について，次が成り立つ．
\begin{enumerate}
\item $\operatorname{im}\mu=\mathcal I=(f_1,f_2,y)$ であり，
  $T=V_B(f_1,f_2)$ を $y=0$ によって $C$ に埋め込むと
  \[
     \coker\mu\simeq\OO_T,\qquad
     \operatorname{length}(\coker\mu)
       =\deg\gcd(f_1,f_2)\in\{0,1,2\}.
     \tag{Sq-defect}\label{eq:qa-spectral-square-defect}
  \]
  この余核は二根対の共通零点の
  \emph{scheme-theoretic} な構造層である．
\item 各幾何 fiber において，積の随伴射
  \[
   \mu^\flat:\mathcal M
       \longrightarrow
       \mathcal H om_C(\mathcal M,\OO_C)
  \]
  は同型である．つまり $\mathcal M$ は自己双対だが，
  $\mu$ が全射となるのは二根対が交わらない場合に限る．
\end{enumerate}
特に $\mu^\flat$ の同型と
可逆線束としての $\mathcal M^{\otimes2}\simeq\OO_C$
とは，節点では異なる主張である．
\end{proposition}
\begin{proof}
相対射 $r:\widetilde C\to C$ に関する
同時符号の反不変元は局所的に $c,d$ で生成され，
$c^2=f_1,\ d^2=f_2,\ cd=y$ である．
従って積の像は記載のイデアルで，
\[
 \OO_C/(f_1,f_2,y)
 \simeq\OO_B/(f_1,f_2)
\]
が従う．$f_i$ は各々 squarefree なので，
その共通零点は $0,1,2$ 個の被約点である．

自己双対性は，二形式が共通根を持たない開集合では
通常の非分岐二重 torsor の二次積から従う．
共通根では完備局所座標と単元の平方根を選ぶと，
\[
 A_0=k[[x,y]]/(y^2-x^2),\quad
 M_0=\langle c,d\mid yc=xd,\ yd=xc\rangle,
 \qquad c^2=d^2=x,\quad cd=y .
\]
ここで $u=c+d,\ v=c-d$ とすれば
\[
  M_0\simeq A_0/(y-x)\,u
          \ \oplus\ A_0/(y+x)\,v,
  \qquad u^2=2(x+y),\quad
  v^2=2(x-y),\quad uv=0.                 \tag{Node-M}
\]
環 $A_0$ において
\[
 \begin{split}
  \Hom_{A_0}(A_0/(y-x),A_0)&=(x+y),\\
  \Hom_{A_0}(A_0/(y+x),A_0)&=(x-y)
 \end{split}
\]
であり，積からの $u\mapsto 2(x+y)$ と
$v\mapsto2(x-y)$ は各直和因子で同型となる．
異なる枝間の Hom は零なので，$\mu^\flat$ は同型である．
各共通根は同じ局所型であり，主張が従う．
\end{proof}

\begin{theorem}[noncubic 三型の compactified spectral limit]
\label{thm:qa-spectral-three-fiber-limits}
前命題の仮定を置き，$\nu:\widehat C\to C$ を
節点での正規化とする．
\begin{enumerate}
\item 二根対が交わらない smooth 型では
  $\mathcal M$ は非自明な $\Pic^0(C)[2]$ の線束であり，
  $\End_C(\mathcal M)=k$ である．
\item ちょうど一根共有する既約 nodal 型では
  $\widehat C\simeq\PP^1$ であり，
  \[
       \mathcal M\simeq\nu_*\OO_{\PP^1}(-1),
       \qquad \End_C(\mathcal M)=k .
       \tag{Node-limit}\label{eq:qa-spectral-nodal-limit}
  \]
  これは次数零の rank-one torsion-free \emph{単純}層だが，
  $\Pic^0(C)\simeq\Gm$ の非自明二捩れ線束
  （節点の貼合せ比 $-1$）とは同型でない．
\item 二根対が一致する $I_2$ 型では
  $\widehat C=C_+\amalg C_-$，$C_\pm\simeq\PP^1$ であり，
  \[
       \mathcal M\simeq
       \nu_*(\OO_{C_+}(-1)\oplus\OO_{C_-}(-1)),
       \qquad \End_C(\mathcal M)\simeq k\oplus k .
       \tag{I2-limit}\label{eq:qa-spectral-i2-limit}
  \]
  従って $\Aut_C(\mathcal M)\simeq\Gm^2$ で，
  この層は\emph{単純でない}．
\end{enumerate}
三型とも $\chi(\mathcal M)=0$ である．
従って，前定理の全相対族を受けるモジュライを
\emph{単純な} rank-one torsion-free sheaf に限定すると，
$I_2$ 型の点が脱落する．
ただし適切な偏極に関する半安定層の粗モジュライを
排除するものではない．
\end{theorem}
\begin{proof}
(1) は定理
~\ref{thm:qa-noncubic-rootpair-torsionfree-spectral-module}
により従う．
各節点での式 (Node-M) は
正規化の二枝の直像の局所模型そのものである．
従って非局所自由な全節点で正規化すると，
$\mathcal M=\nu_*Q$ となる線束 $Q$ が存在する．

(2) の正規化は種数零の滑らかな連結曲線であり，
$Q\simeq\OO_{\PP^1}(d)$ と書ける．
既知の
$p_*\mathcal M\simeq\OO_{\PP^1}(-1)^{\oplus2}$
によって $\chi(\mathcal M)=0$ なので，
$\chi(Q)=d+1=0$，従って $d=-1$ である．
また
$\mathcal E nd_C(\nu_*Q)\simeq\nu_*\OO_{\widehat C}$
であるから，大域 End は $k$ となる．
$\Pic^0(C)$ の元は節点でも局所自由なので，
非局所自由な $\mathcal M$ とは同型にならない．

(3) では比例する二形式をスカラーで規格化して
$f_1=f_2=f$ とできる．曲線は $y=\pm f$ の二成分で，
正規化後の各成分は $B\simeq\PP^1$ に同型である．
$y$ の $c,d$ 上の作用は
$\left(\begin{smallmatrix}0&f\\f&0\end{smallmatrix}\right)$
なので，$c+d$ と $c-d$ がそれぞれ
$y=f$ と $y=-f$ に支持される固有生成元となる．
$p_*\mathcal M=\OO_B(-1)c\oplus\OO_B(-1)d$ から
各固有成分は $\OO_B(-1)$ である．
異なる既約成分間には非零な sheaf Hom はなく，
各成分の自己準同型環は $H^0(\OO_{\PP^1})=k$ だから
End は $k\oplus k$ となる．
最後の留保は，simple/fine と
semistable/coarse の compactification を
混同しないためである．
\end{proof}

\begin{remark}[相対 avatar に必要なモジュライの選択]
\label{rem:qa-spectral-compactification-choice}
既約 nodal 型では $\nu_*\OO_{\PP^1}(-1)$ が
単純な compactified Picard の境界点となる一方，
$I_2$ 型では automorphism が $\Gm^2$ になる．
rank-one torsion-free sheaf の compactification と
安定性条件の一般論は
\cite{AltmanKleiman1980,ChristPayneShen2023} を参照する．
本稿の固有な入力は式~\eqref{eq:qa-spectral-square-defect}と
二根対の三型に沿った explicit な極限表示である．
今後の正確な比較対象は，
相対曲線 $C/B$ と $\mathcal M$ のみならず，
\emph{自己双対な二次積 $\mu^\flat$ と
長さ $0,1,2$ の積欠損}も保持する群oidである．
この群oidと avatar AS 心の同値は未証明である．
\end{remark}

% !TEX root = ../nc-quadric-surface-avatar.tex
% 2026-11-23（canonical square完成の相対tails認識を含む）
\section{内在的なモジュライへの忘却}
\label{sec:qa-intrinsic-forgetting}

有限 foundation、周期的 AS 代数、avatar、および構造対象だけを指定した
Abel 圏は、保持する構造も同型の意味も異なる。
本節では、前節までの開部分と内在的モジュライとの関係を、
既に構成した射と、相対的な認識を要する射とに分けて述べる。

\subsection{標的となる強い内在的クラス}

本稿の \(\HH\)、\(\HHstack\) は
\cite[\texttt{def:h2}]{QuadricIntrinsic2026} と
\cite[\texttt{def:relative-h2}]{QuadricModuli2026} の意味で用いる。
特に、任意の三次元 AS 正則 \(\mathbb Z\)-代数の tails 圏を
無条件に \(\HH\) と呼ぶわけではない。

\begin{definition}[体上の強い内在的クラス]
\label{def:qa-h2}
\(\HH\) の対象は \((\mathcal C,\OO)\) であり、次を満たす。
\begin{enumerate}
\item \(\mathcal C\) は \(\kk\)-線形 noetherian Abel 圏で、
\(D^b(\mathcal C)\) は Hom 有限で Serre 関手 \(S\) を持つ。
\item \(\OO\) は rank-one torsion-free exceptional 対象である。
ここで $\mathcal C_0$ は有限長対象の部分圏、$\mathcal C_1$ は
$\mathcal C/\mathcal C_0$ で有限長となる対象の逆像とし、
$q:\mathcal C\to\mathcal C/\mathcal C_1$ とする。
条件は $q\OO$ が非零単純かつ $\Hom(\mathcal C_1,\OO)=0$ である。
また exceptional とは $\End(\OO)=\kk$、$\Ext^{i}(\OO,\OO)=0$ ($i>0$) をいう。
\item 完全自己同値 \(s\) が存在し、
\(s^2\simeq S^{-1}[2]\) かつ \((\OO,s)\) は Artin--Zhang ample である。
\item 同じ \(s\) に対して、全軌道環
\[
 B(\mathcal C,\OO,s)
   =\bigoplus_{n\ge0}\Hom(\OO,s^n\OO)
\]
が \(A/(g)\) と同型である。
ここで \(A\) は左右 noetherian、二次 Koszul AS 正則、
型 \((1,4,6,4,1)\) の環、\(g\in A_2\) は左右正則正規元である。
\end{enumerate}
\(s\)、根同型、\(A\)、\(g\)、商表示は指定構造ではない。
射は \(\OO\) の同型を伴う完全同値であり、
その同型を保つ自然同型によって一射を同一視する。
\end{definition}

ここで型 \((1,4,6,4,1)\) は、自明加群の最小解消の階数列が
この列で、global dimension と Gorenstein parameter がともに四であることを意味する。
\(g\) の中心性は要求しない。
また、ある商圏との同値だけでなく、指定した \(\OO\) と同じ \(s\) の
\emph{全軌道環} に対する表示を要求する。

相対版 \(\HHstack\) は、標数零の代数閉体上の局所有限型スキームの
fppf site 上で定義する。
圏の基底変更は Ind 圏における加群対象の圏によるものであり、
Hom を tensor するだけの基底変更ではない。
対象と完全列の fppf descent、有限 Hom と flat base change、
平坦な \(\OO\) を要求する。
局所 canonical compact 導来圏は相対 smooth proper で、
相対 Serre 双対性と導来基底変更に適合する。
全幾何ファイバーは定義~\ref{def:qa-h2} を満たす。
さらに、\(s\)、AZ 条件、全軌道環の AS4 被覆は
\emph{非被約基底も含めて fppf 局所に存在する}ことを要求する。
根や被覆を重なり上で同一視することは要求しない。
これらは標識ではなく存在条件だからである。

\begin{remark}[相対存在は幾何点条件より強い]
\label{rem:qa-relative-root}
全幾何ファイバーが \(\HH\) に入っても、
相対平方根、AZ ampleness、相対 AS4 被覆が存在することは別問題である。
例外対象の無障害性や有限代数の \(\operatorname{HH}^3\) の消滅だけでは、
これらの条件は得られない。
\cite[\texttt{prop:chosen-root-does-not-lift}]{QuadricModuli2026} は、可換 quadric 上の
ある適合根が、特定の一次 Abel 変形へ持ち上がらない例を与える。
これは別の根の存在まで否定する例ではなく、
\(\HHstack\) の非代数性や許容部分の非開性を証明するものでもない。
\end{remark}

\subsection{canonical square 完成で確立した定義域}

定理~\ref{thm:qa-relative-square-completion-tails} と
命題~\ref{prop:qa-relative-square-unfolding} により，実際の曲線 avatar
から構成した有限 square 代数の canonical 完成と全整数展開には，
標準対象を保つ相対 tails 認識がある．これは幾何点の同値だけではなく，
本稿の site 上の module-object 基底変更と coherent に整合する完全同値である．
左右 Noether 性と係数束付きの実際の AS 解消も証明済みである．
補題~\ref{lem:qa-classified-fiber-finite-square-recognition} は，
古典型を含む分類済み square fiber を持つ有限局所自由族まで扱う．
従ってこの canonical 部分の相対 heart 回復を未解決のままにしない．
一方，任意の別途与えた AS/QP 族との全整数同一視には
注意~\ref{rem:qa-square-unfolding-vertex-lines} の頂点線束・周期比較が必要であり，
以下の一般 presentation stack 全体の認識と区別する．

さらに，定理~\ref{thm:qa-square-normal-foundation-stack-equivalence} は
有限 square と normal foundation の全スタックを同定し，
系~\ref{cor:qa-h2-square-foundation-open} は既存の内在的開をこの表示へ移す．
全 $\HHstack$ については
定理~\ref{thm:qa-universal-h2-rank81-completion} が，ruling の選択を要しない
階数八十一の通常代数と pointed tails を与える．
ここまでの有限表示・canonical 完成の認識は確立済みである．
定理~\ref{thm:qa-square-explicit-h2-open}と
系~\ref{cor:qa-square-h2-twenty-polynomials}は，この全$H_2$開を
三十二次元線形表示上の四組・計二十式以内で記述する．
任意の別途与えた AS/QP presentation 群oidとの比較は後続課題である．

\subsection{忘却射の正確な定義域}

\begin{definition}[内在的条件に適合する族]
\label{def:qa-compatible-chart}
\(\mathfrak U\) を、周期的代数または avatar から
基底変更と同型に関して自然に \((\mathcal C,\OO)\) を構成する
fppf パラメータスタックとする。
\(\mathfrak U_{\HH}\) は、その構成した族が
上記の \(\HHstack\) の相対条件を満たす対象の全部分スタックとする。
この定義において、\(\mathfrak U_{\HH}\) が代数的または開であることは
仮定も結論もしない。
\end{definition}

\begin{proposition}[適合族からの忘却]
\label{prop:qa-compatible-forgetting}
定義~\ref{def:qa-compatible-chart} の構成は射
\[
 \pi_{\mathfrak U}:\mathfrak U_{\HH}\longrightarrow\HHstack,
 \qquad u\longmapsto(\mathcal C_u,\OO_u)
\]
を与える。
二つの chart 間の変換が、基底変更と両立する pointed exact 同値を
与えるなら、二つの忘却射はその変換に沿って自然に二同型となる。
\end{proposition}
\begin{proof}
\(\mathfrak U_{\HH}\) の条件は fppf 局所的で、
pointed exact 同値と基底変更に安定である。
従って同型で飽和した全部分スタックを定める。
構成した圏の系、構造対象、その比較同型を保持して
残る標識を忘れれば、定義された標的への functor を得る。
二つの構成を結ぶ pointed exact 同値が任意の基底変更と整合する場合には、
それらの同値がそのまま二つの functor の二同型を与える。
\end{proof}

この命題は忘却の意味を明示するものであり、
未知の条件を有限個の方程式で認識したという定理ではない。
特に \(\mathfrak F_{\le12}\) は有限代数だけを分類するので、
そこから直接 \(\mathcal C\) を取ることはできない。
\(\operatorname{Perf}(A)\) の構成だけでは、元の曲面の heart は選ばれていない。
有限代数の標準的な加群 heart と、tilting によって移された曲面の heart を
同一視してはならない。

「十二本矢の適当な開部分から \(\HHstack\) への自然な射」という主張を
有限な表示で実現するためには、少なくとも次を示す必要がある。
\begin{enumerate}
\item 使用する periodic algebra、grading、foundation と heart を
相対的に構成し、全次数の平坦性、noetherian 性、基底変更を確認する。
\item 指定した \(\OO\) の相対例外性・rank 条件、相対 Serre 公式、
適合する完全平方根と AZ ampleness を確認する。
\item 同じ根の全軌道環が、相対 AS4 環の正則正規二次切断となることを確認する。
\item これらの条件を持つ部分関手が、考えている表示スタックの
不変開部分で表されることを証明し、表示変更の全ての射を降下させる。
\end{enumerate}
根と全軌道環が既に与えられた場合、固定軌道環への AS4 被覆を選ぶ空間は
\cite[\texttt{thm:relative-finite-as-cover}]{QuadricModuli2026} により、
\(\operatorname{Gr}(6,K)\) 内の有限表示局所閉 incidence として記述できる。
しかし、その被覆を忘れた fppf 像の開性・代数性や、根そのものの存在は
この有限判定の結論ではない。

\subsection{十二本矢の全体を送ることへの具体的な障害}

\begin{lemma}[反 Serre 軌道の生成は必要である]
\label{lem:qa-antiserre-necessary}
\((\mathcal C,\OO)\in\HH\) ならば、
\(t=S^{-1}[2]\) に対する \((\OO,t)\) も AZ ample である。
特に、非零対象 \(M\) に対して
\(\Hom(t^{-n}\OO,M)=0\) が全ての \(n\ge0\) で成立することはない。
\end{lemma}
\begin{proof}
定義の根 \(s\) を選び、\(E=\OO\) と置く。
\(B=A/(g)\) は次数一で生成される。
AZ の生成条件を \(sE\) に適用し、得られた全射に
\(s^{-1}\) を作用させると、\(E\) は
\(s^{-m}E\)（\(m>0\)）の有限直和の商になる。
次数一生成により各成分 \(s^{-m}E\to E\) は、
最後の段が \(s^{-1}E\to E\) である合成の線形結合である。
従って評価 \(\Hom(s^{-1}E,E)\otimes s^{-1}E\to E\) は全射となる。
Hom 有限性により、この評価の源は有限直和である。
これは \cite[\texttt{lem:ample-square-from-generation}]{QuadricIntrinsic2026}
の評価全射である。

任意の対象への AZ 生成全射で奇数指数 \(m\) の項があれば、
この評価全射を \(s^{-m}\) で移して \(s^{-(m+1)}E\) の
有限直和に置き換える。従って偶数指数の項だけで生成できる。
指数に任意の下限を課す場合は、先に対象を正に shift して
生成全射を作り、負の shift で戻す。
全射に対する Hom の最終的な全射性は
元の AZ 条件を偶数部分列に制限すれば従う。
\(s^2\simeq t\) なので \((E,t)\) は AZ ample であり、
最後の主張はその生成条件から従う。
\end{proof}

\begin{example}[重み付き十二本矢と residual gerbe]
\label{ex:qa-weighted-not-h2}
\(S=\kk[x,y,z]\)、\(\deg x=\deg y=1\)、\(\deg z=2\) とし、
\(\mathcal X=\PP(1,1,2)\) を重み付き射影 \emph{stack} とする。
この環は global dimension 三、Gorenstein parameter 四の
noetherian AS 正則環である。
しかし
\[
 (\coh\mathcal X,L)\notin\HH\qquad(L\in\Pic(\mathcal X)).
\]
\end{example}
\begin{proof}
通常の三変数 Koszul 解消に各変数の重みを入れると
\[
 0\longrightarrow S(-4)
 \longrightarrow S(-3)^2\oplus S(-2)
 \longrightarrow S(-1)^2\oplus S(-2)
 \longrightarrow S\longrightarrow\kk\longrightarrow0
\]
を得る。これは AS 条件と Gorenstein parameter を与える。

\((\OO,\OO(1),\OO(2),\OO(3))\) の有限 End 代数を考える。
隣接する \(x_i,y_i\) が六本、距離二の \(z_0,z_1\) が二本であり、
最小関係は
\[
 \begin{aligned}
 x_1y_0-y_1x_0&=0,&x_2y_1-y_2x_1&=0,\\
 z_1x_0-x_2z_0&=0,&z_1y_0-y_2z_0&=0.
 \end{aligned}
\]
従って relation completion は十二本矢である。
前節の四写像表示では \(a,b\) はともに階数三で、核は交代テンソルの直線、
\(c,d\) は全射である。よってこの点は実際に
\(\mathfrak F_{\le12}\) の最小十二本矢 locus に属する。
この四対象が full strong であることも直接確認できる。
重み付き多項式環の局所コホモロジーから
\(H^1(\mathcal X,\OO(n))=0\) であり、
Serre 双対により \(n\ge-3\) では \(H^2(\mathcal X,\OO(n))=0\) である。
これで strong 性が従う。
三変数の sheaf Koszul 完全列は四つの連続する twist から
全ての twist を生成し、有限生成 graded \(S\)-加群の有限自由解消により
これらは \(D^b(\coh\mathcal X)\) 全体を生成する。

一方、重み付き Euler 列から
\(\omega_{\mathcal X}=\OO(-4)\) であり、\(t=-\otimes\OO(4)\) である。
点 \([0:0:1]\) の residual gerbe を
\(i:B\mu_2\hookrightarrow\mathcal X\) とし、
\(\chi\) を \(\mu_2\) の非自明な一次元表現とする。
任意の可逆層 \(L\) に対し、非零 coherent 対象
\(M_L=i_*(i^*L\otimes\chi)\) を取る。
\(\OO(4n)\) の gerbe 上の指標は自明なので、引戻しと直像の随伴により
\[
 \Hom(L(-4n),M_L)=
 H^0(B\mu_2,i^*\OO(4n)\otimes\chi)=0\qquad(n\ge0).
\]
特に標準 helix の四位相すべてで生成条件が破れる。
補題~\ref{lem:qa-antiserre-necessary} に反するので、主張を得る。
\end{proof}

\begin{example}[\(\Sigma_2\) 自身の構造対象]
\label{ex:qa-hirzebruch-not-h2}
\(Y=\Sigma_2\) の負断面を \(C\)、ファイバーを \(f\) とする。
\(C^2=-2\)、\(-K_Y=2C+4f\) なので、
\(\OO_Y(-K_Y)|_C\simeq\OO_C\) である。
任意の可逆層 \(L\) に対し、
\(M_L=i_*(i^*L\otimes\OO_C(-1))\ne0\) とおくと
\[
 \Hom\bigl(L\otimes\OO_Y(nK_Y),M_L\bigr)
   =H^0(C,\OO_C(-1))=0\qquad(n\ge0).
\]
同じ必要条件から \((\coh\Sigma_2,L)\notin\HH\) となる。
これは \emph{avatar の一点 Hilbert 曲面が \(\Sigma_2\) である}ことと、
元の Abel 圏が \(\coh\Sigma_2\) であることを混同してよいという意味ではない。
Hilbert 曲面の型だけから、対応する非可換圏の \(\HH\) 適合性は決まらない。
\end{example}

この二例は、十二本矢の全 locus、AS 正則性、あるいは
full strong foundation の存在だけで忘却の標的が \(\HHstack\) になる、
という主張を否定する。
適切な開部分が存在する可能性を否定するものではなく、
その開部分には構造対象の反 Serre 軌道に対する生成条件も必要であることを示す。

\subsection{標識、射、既存の表示スタック}

\(\HHstack\) への忘却では \(\OO\) を保持する。
従って、foundation の mutation が構造対象を別の対象へ移す場合には、
そのまま同じ pointed 対象の表示変更と扱えない。
QP の mutation が dg 導来同値を与えても、
指定した heart の完全同値と \(\OO\) の同型は別途必要である。
\cite[\texttt{cor:pointed-derived-exact}]{QuadricIntrinsic2026} は体上の \(\HH\) 対象の間の判定であり、
任意の相対 periodic family について同じ結論を自動的に与えるものではない。

同様に、消去可能な二本の矢 \(x,y\) と項 \(yx\) を付加した生の QP には、
\((x,y)\mapsto(\lambda x,\lambda^{-1}y)\) という自己同型がある。
消去後の表示ではこの自己同型は見えなくなる。
従って、生の安定化 QP groupoid の消去を、
射を補正しないままスタック同値とすることはできない。
有限 foundation の頂点スカラーの剛体化、
周期的環の次数スカラーの除去、
構造対象を保つ自己同値全体での同一視は、それぞれ異なる操作である。

現在の内在的研究との関係をまとめると、
\cite[\texttt{prop:foundation-stack}]{QuadricModuli2026} の有限二写像表示 \(\mathfrak F_0\) は、
一方の長さ二の積を可逆にした foundation chart である。

\begin{proposition}[幾何点における十本矢 chart の被覆]
\label{prop:qa-h2-normal-ten-chart}
体上の \(\HH\) の任意の対象は、
\(\mathfrak F_{\le10}^{(a)}\) に属する normal foundation を持つ。
これは幾何点での主張であり、任意の \(\HHstack\) の族における
foundation の相対的な存在・降下を結論するものではない。
\end{proposition}
\begin{proof}
\cite[\texttt{thm:normal-helix}]{QuadricIntrinsic2026} と
\cite[\texttt{prop:ruling-involution-mutation}]{QuadricModuli2026} により、
\((E_0,E_1,E_2,E_3)=(\OO,X,s\OO,Z)\) という foundation があり、
本節の \(b\) は同型、\(c\) は全射である。
\(E_0,E_1,E_2\) は rank-one torsion-free 対象である。
このような対象間の非零射は単射である。
実際、像が低い Gabriel--Krull 段階に入れば標的の torsion-free 性に反し、
従って rank-one 商では非零射、すなわち同型を誘導する。
その核は低い段階に入るが、始域の torsion-free 性によって零である。
従って二つの隣接する非零射の合成は零にならず、
\(\ker a\) は非零の分解可能テンソルを含まない。
\(b\) は同型なので他方の核は零であり、
\(c\) の全射性から定義~\ref{def:qa-foundation-open} の全条件を満たす。
\end{proof}

相対的な ruling 交換を \(\tau\) とすると、
\[
 \mathfrak Q_{\mathrm{fnd}}=[\mathfrak F_0/S_2]
\]
は有限型 Artin スタックであり、体上では
\(\HH(\kk)\) をその full subgroupoid として実現する
（\cite[\texttt{cor:h2-point-groupoid-foundation}]{QuadricModuli2026}）。
この幾何点の結果だけから、\(\mathfrak Q_{\mathrm{fnd}}\) 全体または
その未特定の開部分が \(\HHstack\) と同値であるとはいえない。

相対的には、同論文の反 Serre 認識によって全忠実な射
\[
 \HHstack\hookrightarrow
 \mathfrak P=[X_3/\PGL_9]
\]
がある。
ここで \(X_3\subset\operatorname{Gr}(56,81)\) は、
九生成元・五十六二次関係の環の三次成分が階数四十九となる
scheme-theoretic な局所閉 locus である。
\[
 X_{\mathrm{as}}=X_3\times_{\mathfrak P}\HHstack
\]
は不変 fppf 部分層として定義され、
\(\HHstack\simeq[X_{\mathrm{as}}/\PGL_9]\) が成立する。
\(X_{\mathrm{as}}\) の有限型・滑らか・既約な純次元83のスキームとしての表現可能性は
相対モジュライ論文の定理~\cite[\texttt{thm:global-h2-smooth-finite-type}]{QuadricModuli2026} と
系~\cite[\texttt{cor:global-xas-smooth-scheme}]{QuadricModuli2026} で証明されている。
ただし $X_{\mathrm{as}}\to X_3$ は一致根対で開ではなく、
局所閉性と明示的な偏極認識方程式は未解決である。
従って本稿の八・十・十二本矢の明示的な開部分は、
この相対認識問題への有限表示の候補を与えるが、
それ自体で内在的モジュライの代数性を解決するものではない。

% !TeX root = ../quiver-avatar-moduli.tex
% 2026-11-22 / v0.5（OU z2座標とPicard根の明示対応、Kummer障害を追加）
\section{既存GITと内在的モジュライの比較}
\label{sec:qa-git-comparison}

本節は Okawa--Ueda \cite[Theorem 1.1]{OU2014} の
quintuple と conifold potential の GIT 商を、
本稿の有限 foundation と内在的 $\HHstack$ に対する比較問題へ接続する。
二つの既存 GIT 商自体を本稿の新しい成果とはしない。
また、表示を忘れた $\HHstack$ への射は
定義~\ref{def:qa-compatible-chart} の適合する領域でのみ構成し、
GIT compactification 全体で自動的に存在するとは主張しない。

\subsection{既存の次数四射と三つの忘却}

$V_0,V_1,V_2,V_3$ を二次元、$W$ を一次元の $k$-ベクトル空間とし、
\[
 T=(V_0\otimes V_1\otimes V_2\otimes V_3)\otimes W^\vee,\qquad
 \widetilde G=\prod_{i=0}^3\GL(V_i)\times\GL(W)
\]
と置く。Okawa--Uedaの生のquintuple表示スタックを
$\mathfrak M_{\mathrm{qui}}^{\mathrm{raw}}=[T/\widetilde G]$ と書く。
$T$ の零元を含む表示スタックの記法であり、
$\HHstack$ を定義するものではない。
同論文のpotential表示を $\mathfrak M_{\mathrm{pot}}$ とし、
その compact GIT 商を $\overline M_{\mathrm{pot}}$、
quintuple側の商を $\overline M_{\mathrm{qui}}$ とする。
既知結果は
\[
 \overline M_{\mathrm{pot}}\simeq\PP(1,2,3,4),\qquad
 \overline M_{\mathrm{qui}}\simeq\PP(2,4,4,6),\qquad
 \overline M_{\mathrm{pot}}\longrightarrow\overline M_{\mathrm{qui}}
 \text{ が有限次数 }4
\]
である。右端の射は潜在的なSerre平方根の忘却射ではなく、
\emph{potentialからquintupleへの表示の射}である。

本研究では区別すべき射が少なくとも三つある。
potentialからquintupleへの射、quintupleから四項foundationへの
表示上の射、そして適合foundationから pointed exact 対象
$(\mathcal C,\OO)$ への忘却射である。
最後の射では、quintupleの標識、ruling、Serre平方根とその同型、
ambientを保持しない。これらの射のどの組合せが
相対スタック射として存在するかは、対象と自然同型を個別に調べる必要がある。

\subsection{スカラー慣性の核}

\begin{proposition}[生のquintuple表示の四次元スカラー核]
\label{prop:qa-ou-scalar-gerbe}
次の中央部分群スキームは $\widetilde G$ に含まれる。
\[
 K=\left\{
 \bigl(\lambda_0\mathrm{id}_{V_0},\ldots,
       \lambda_3\mathrm{id}_{V_3},
       (\lambda_0\lambda_1\lambda_2\lambda_3)\mathrm{id}_{W}\bigr)
 :\lambda_i\in\Gm
 \right\}\simeq\Gm^4.
 \tag{GIT-K}\label{eq:qa-ou-scalar-kernel}
\]
$K$ は $T$ 全体に自明に作用する。従って
\[
 [T/\widetilde G]\longrightarrow[T/(\widetilde G/K)]
 \tag{GIT-rig}\label{eq:qa-ou-rigidification}
\]
は $K$ を相対慣性とする fppf gerbe である。

さらに、quintupleから対応する $\mathbb Z$-代数のtails圏と
添字零の構造対象を作る通常の構成が実際の相対族を与え、
その対象が $\HHstack$ の公理を満たす領域では、
$K$ の全元は pointed exact 自己同値の\emph{恒等射}へ写る。
従って、この領域の生の表示スタックから $\HHstack$ への
標準的な忘却射は表現可能ではない。
少なくともこのスカラー慣性を消さずに
全忠実な相対比較射を得ることはできない。
\end{proposition}
\begin{proof}
$\widetilde G$ のスカラー部分の作用は
$T$ 上で
\[
 (\lambda_0,\ldots,\lambda_3,\mu)\cdot t
   =(\lambda_0\lambda_1\lambda_2\lambda_3)\mu^{-1}t
\]
である。従って式~\eqref{eq:qa-ou-scalar-kernel} は
$T$ への作用の核に含まれる中央な四次元トーラスである。
$\widetilde G/K$ は群スキームとして存在し、
$\widetilde G\to\widetilde G/K$ は忠実平坦なので、
$K$ の作用が自明である場合の商スタックの一般的な構成から
式~\eqref{eq:qa-ou-rigidification} はgerbeになる。

後半は表現を変えるスカラーが圏に与える作用を直接調べる。
quintupleを使って構成した $\mathbb Z$-代数を $D$ とし、
その隣接Hom空間を $V_i$（または採用する双対規約）で表す。
スカラー変更は、各頂点の単元を変えない
\[
 D_{ij}\longrightarrow D_{ij},\qquad
 a\longmapsto\gamma_i^{-1}\gamma_j\,a
\]
という次数ごとの共役に帰着する。
ここで隣接空間での所定のスカラーに合うよう
$\gamma_i\in k^\times$ を漸化的に選ぶ。
四つの隣接スカラーの積が一次関係線 $W$ のスカラーと一致する
ことが、quintupleの関係を保つ条件である。
右graded $D$-加群の次数 $i$ の成分に
$\gamma_i^{-1}$ を掛ける自然変換は、この共役による
twist functorと恒等 functorの自然同型を与える。
したがってtailsへの降下でも恒等と自然同型であり、
添字零の構造対象の同型も同じスカラーで規格化できる。
相対族でも単元を用いて同じ計算ができる。

一方、元の表示スタックでは任意の $K$ の元は
各点の自己同型として残る。
忘却による相対慣性の核が $K$ を含むため、
射が表現可能であるための必要条件（相対慣性の自明性）に反する。
\end{proof}

\begin{remark}[GITの次数と慣性を混同しない]
\label{rem:qa-ou-degree-four-not-root}
Okawa--Ueda \cite[Theorem 1.1(4)]{OU2014} の次数四は、
compact GIT 商の有限射の一般ファイバーの次数である。
式~\eqref{eq:qa-ou-scalar-kernel} の群は
quintuple表示の自己同型に既に存在し、
pointed忘却では消える。
したがってその次数四を、標識を忘れた
$\HHstack$ の一般ファイバーの位数、
あるいは可換点でのSerre平方根の個数と同一視できない。
適合する族の範囲、表示側の剛体化、全pointed同型の像と核、
および閉軌道の選択が別途必要である。
\end{remark}

\subsection{相対ファイバーの群oidとしての公式}

\begin{proposition}[忘却射の幾何ファイバーと慣性核]
\label{prop:qa-forgetful-homotopy-fiber}
$F:\mathfrak X\to\mathfrak Y$ を $k$-上の groupoid stack の射とし、
$y\in\mathfrak Y(k)$ を取る。
幾何ファイバーの groupoid
$\mathfrak X_y=\mathfrak X\times_{\mathfrak Y}\operatorname{Spec}k$
の対象は、$x\in\mathfrak X(k)$ と同型
$\alpha:F(x)\xrightarrow{\sim}y$ の組である。
その点 $(x,\alpha)$ の自己同型群は
\[
 \operatorname{Aut}_{\mathfrak X_y}(x,\alpha)
 =\ker\bigl(\operatorname{Aut}_{\mathfrak X}(x)
       \longrightarrow\operatorname{Aut}_{\mathfrak Y}(F(x))\bigr).
 \tag{Fib-I}\label{eq:qa-fiber-inertia}
\]
さらに、$F(x)\simeq y$ となる $\mathfrak X(k)$ の同型類を
一つずつ選ぶと、$\mathfrak X_y(k)$ の同型類の集合は
\[
 \coprod_{[x]:\,F(x)\simeq y}
 \operatorname{Aut}_{\mathfrak Y}(y)\Big/
  \alpha\,F(\operatorname{Aut}_{\mathfrak X}(x))\,\alpha^{-1}
 \tag{Fib-C}\label{eq:qa-fiber-cosets}
\]
と同一視される。右辺は群の商ではなく、一般には剰余\emph{集合}である。
従って同型類のファイバーの濃度、スタックの相対慣性、
GITの有限射の次数はそれぞれ異なる量である。
\end{proposition}
\begin{proof}
二つの対象 $(x,\alpha),(x',\alpha')$ の間の射は
$f:x\to x'$ で $\alpha'F(f)=\alpha$ を満たすものである。
$x=x'$、$\alpha=\alpha'$ とすれば
式~\eqref{eq:qa-fiber-inertia} が直ちに従う。
固定した $x$ について同型 $\alpha:F(x)\to y$ を一つ選ぶと、
他の全ての同型は $g\alpha$、
$g\in\operatorname{Aut}_{\mathfrak Y}(y)$ の形である。
二つの選択は $\operatorname{Aut}_{\mathfrak X}(x)$ の像による作用で
同型になるため、式~\eqref{eq:qa-fiber-cosets} が得られる。
\end{proof}


\subsection{次数四射の $S_4$-monodromy と Picard 平方根}

\begin{proposition}[OU次数四射の一般Galois閉包]
\label{prop:qa-ou-fourfold-s4-monodromy}
標数零の代数閉体上で、Okawa--Ueda \cite[\S5, (5.8)]{OU2014}
のGIT有限射を、対角化による稠密な一般開で考える。
同論文の表示は
\[
 \PP^3/S_3\longrightarrow
 \PP^3/\bigl(S_3\ltimes(\mathbb Z/2)^3\bigr)
 \tag{OU-4}\label{eq:qa-ou-fourfold-diagonal}
\]
であり、$\PP^3$ では同時符号反転が自明に作用する。
従って有効な符号変更群と全群はそれぞれ
\[
 \Gamma=(\mathbb F_2)^3/
   \langle(1,1,1)\rangle\simeq(\mathbb F_2)^2,\qquad
 G=\Gamma\rtimes S_3\simeq\operatorname{AGL}_2(\mathbb F_2)
 \simeq S_4.
\]
$S_3$ は $\Gamma$ の三つの非零元を置換する標準的な
$\operatorname{GL}_2(\mathbb F_2)$-作用を持つ。
一般の射~\eqref{eq:qa-ou-fourfold-diagonal} は次数四で、
\emph{Galois射ではなく}、その関数体拡大のGalois閉包の群は $S_4$ である。
$\Gamma$ を自明化する表示の方へ引き戻した場合にのみ、
$\Gamma$-torsorとしての四点の符号変更を直接見ることができる。
\end{proposition}
\begin{proof}
$\mathbb F_2^3$ の標準基底の剰余類を $e_1,e_2,e_3$ とする。
商 $\Gamma$ では $e_1+e_2+e_3=0$ なので、
三つはまさに非零元である。$S_3$ がこれらを置換する作用は忠実で、
$\operatorname{Aut}_{\mathbb F_2}(\Gamma)
=\operatorname{GL}_2(\mathbb F_2)$ は位数六なので同型となる。
$\Gamma$ の四元上への平行移動とこの線型作用を合わせた
$G$ は四元上の忠実で推移的な位数 $4\cdot6=24$ の置換群、
従って $S_4$ である。
$H=S_3$ は原点の安定化群で、$[G:H]=4$。
標数零の一般対角座標はこの有限群作用に対し自明な
幾何安定化群を持つので、
関数体拡大 $k(\PP^3)^H/k(\PP^3)^G$ の次数は四である。
$H$ の $G$ における core（全共役の共通部分）は、
四元上の推移作用の核と一致し、忠実性から自明である。
従ってこの非正規拡大のGalois閉包は $k(\PP^3)$ であり、
Galois群は $G\simeq S_4$ となる。
最終のtorsorに関する文は、
$H$ の選択を消した一般の基底では
$\Gamma$ が一定群として作用する四重Galois被覆を
無条件に得られない、という意味である。
置換群の全24元とcoreは
\path{computations/verify_ou_fourfold.py} で再現できる。
\end{proof}

\begin{proposition}[適合Picard平方根の幾何fiberの濃度]
\label{prop:qa-picard-roots-degeneration-count}
論文Bの非線型geometric cubic quintupleの相対曲線
$C\to S$ と $N\in\Pic^0_{C/S}(S)$ を考え、
同論文 \cite[\texttt{lem:picard-halves}]{QuadricModuli2026}
の平方根空間
\[
 T=S\times_{\Pic^0_{C/S},[-2]}\Pic^0_{C/S}
 \tag{PR-4}\label{eq:qa-picard-root-torsor}
\]
を取る。これは $S$ 上 étale 全射である。
滑らかな楕円曲線の locus $S_{\mathrm{sm}}$ 上では
$T_{S_{\mathrm{sm}}}\to S_{\mathrm{sm}}$ は
$\Pic^0_{C/S}[2]$-torsorであり、有限 étale 次数四である。
各代数閉幾何点での平方根の個数は、曲線が
滑らかな楕円曲線、既約一節点の種数一曲線、
既約一尖点の種数一曲線である場合、それぞれ
\[
 \begin{array}{c|ccc}
 C_{\bar s}&\text{smooth elliptic}&\text{nodal}&\text{cuspidal}\\ \hline
 \Pic^0(C_{\bar s})&E&\mathbb G_m&\mathbb G_a\\
 \#T_{\bar s}&4&2&1
 \end{array}
 \tag{PR-count}\label{eq:qa-picard-root-count}
\]
となる。従って滑らかなfiberと節点fiberが同じ連結基底族に
現れるとき、$T\to S$ はその族上で有限 étale
\emph{次数四のまま}延長することはできない。
\end{proposition}
\begin{proof}
$\Pic^0_{C/S}$ は滑らかな相対一次元可換群代数空間で、
$[-2]$ は étale 全射であることを論文Bの
\cite[\texttt{lem:picard-halves}]{QuadricModuli2026} が示す。
滑らかな種数一のfiberのPicard零成分は楕円曲線の
相対Jacobianであり、$[-2]$ は有限 étale
次数 $2^{2}=4$ のisogenyである。
そのpullbackは四点の$E[2]$-torsorとなる。

一節点曲線の正規化は $\PP^1$ で、節点の二つの原像の
fiberの貼合せスカラーにより
$\Pic^0\simeq\mathbb G_m$ となる。
$[-2]$ は $z\mapsto z^{-2}$ なので、
代数閉体の各非零値の逆像は二点である。
一尖点曲線では正規化との貼合せの一次jetの加法群から
$\Pic^0\simeq\mathbb G_a$ となる。
$[-2]$ は $z\mapsto-2z$、従って一意の平方根を持つ。
いずれも標数零では étale である。
最後の主張は有限 étale 射の次数が連結基底上で
局所定数となることから従う。
\end{proof}

\begin{remark}[四点の抽象的な一致と比較すべきtorsor class]
\label{rem:qa-ou-picard-fourfold-not-identical}
OU側の $\Gamma\simeq\mathbb F_2^2$ と
楕円曲線上の $E[2](\bar k)\simeq\mathbb F_2^2$ は、
$S_3=\operatorname{GL}_2(\mathbb F_2)$ の
三非零元に対する作用も含めて同じ抽象モデルを持つ。
これは次数四の一致を説明する\emph{候補}であって、
OUの有限射とPicard平方根空間の同型を証明するものではない。
特にOU側は一般に$S_4$-monodromyの非Galois四重射であり、
Picard側は変動する群スキーム$E[2]$のtorsorである。
二つが同じ基底上に引き戻され、
$E[2]$の局所系とOUの符号変更局所系が同定できた後に、
両torsorの相対同型を確認する必要がある。
標識、自己同型群、さらにpointed同値による商は
その後に初めて比較できる。
節点・尖点における式~\eqref{eq:qa-picard-root-count} は
Picard根空間そのもののfiberの個数であり、
OU射の分岐fiberや$\HHstack$への忘却fiberを
数えたものではない。
\end{remark}

\begin{proposition}[OU楕円四つ組の $z_2$ と適合Picard根]
\label{prop:qa-ou-elliptic-half-coordinate}
滑らかな楕円曲線 $E$ に原点 $o$ を取り、
Okawa--Ueda \cite[\S4, (4.12)--(4.13)]{OU2014}
の楕円四つ組を
\[
 (E,\mathcal L_0,\mathcal L_1,\mathcal L_2)
 =
 (E,\OO_E(2o),\OO_E(2z_1),\OO_E(2z_2))
\]
と書く。論文Bの
$N=\mathcal L_2\mathcal L_0^{-1}\in\Pic^0(E)$ に対する
$a^{-2}=N$ の解集合は
\[
 \left\{
 \OO_E(o-z_2)\otimes\xi:\xi\in E[2]
 \right\}
 \tag{OU-PR}\label{eq:qa-ou-elliptic-half-coordinate}
\]
である。従ってOUの楕円四つ組の birational 表示において、
$\mathcal L_0,\mathcal L_1,\mathcal L_2$ を固定したまま
$z_2$ に2-torsionを加える操作は、
\emph{対応するPicard平方根の四つの選択}に一致する。
これは式~(4.13)の $(\mathbb Z/2)^4$ のうち
$z_2$ に作用する一つの $(\mathbb Z/2)^2$ の因子である。

ただし、式~\eqref{eq:qa-ou-fourfold-diagonal} の
potential GITからquintuple GITへの次数四の
\emph{四つのsheet}がこの $z_2$ 因子のtorsorと同型であることは、
この命題の結論ではない。
\end{proposition}
\begin{proof}
$N=\OO_E(2z_2-2o)$ であるから
$a_0=\OO_E(o-z_2)$ は
$a_0^{\otimes(-2)}=N$ を満たす。
二つの根$a,a_0$の比は二乗が自明な
$\Pic^0(E)[2]=E[2]$ の元であり、逆にその全元は解を与える。
$z_2$ の取り換え $z_2\mapsto z_2+t$（$2t=0$）は
$\OO_E(2z_2)$ を変えず、$a_0$ を
$\OO_E(-t)$ 倍する。従って四つの選択と自由かつ推移的に対応する。
OUの式~(4.11)--(4.13)は $z_1,z_2$ の独立した
2-torsion平行移動を含むので、最後の因子の同定も従う。
\end{proof}

\begin{lemma}[Picard平方根のKummer class]
\label{lem:qa-picard-kummer-torsor}
滑らかな楕円曲線の相対族 $E\to S$ と
$N\in E(S)$ に対し、平方根空間
$T_N=S\times_{E,[-2]}E$ の
$E[2]$-torsor class は、étale siteのKummer列
\[
 0\longrightarrow E[2]\longrightarrow E
 \xrightarrow{[-2]}E\longrightarrow0
\]
から得られる連結写像による
$\delta_{[-2]}(N)\in H^1_{\mathrm{\acute et}}(S,E[2])$
である。
従って、同じ基底上でOUの四重射に
$E[2]$-torsor構造を与える局所系の同型を
別途構成できた場合に限り、そのtorsorとPicard根空間が同型となる
必要十分条件は二つの $H^1$ class の一致である。
\end{lemma}
\begin{proof}
$[-2]$ は滑らかな楕円群スキーム上の有限étale isogenyで
核は$E[2]$なので、示した短完全列はétale層として完全。
$N$ の局所持上げの差から得るČech 1-cocycleこそ
$T_N$ のtorsor classである。
比較先にも同じ構造群スキームを同定した後は、
torsorの同型類の分類が$H^1$で与えられることから従う。
\end{proof}

\subsection{一般非線型 cubic の実際の相対比較射}

\begin{theorem}[geometric quintuple から pointed 相対スタックへの射]
\label{thm:qa-geometric-cubic-relative-forgetting}
$\mathscr G^{\mathrm{nl}}$ を、式~\eqref{eq:qa-ou-scalar-kernel}
のテンソル表示空間 $T\setminus\{0\}$ のうち
Van den Bergh の\emph{非線型 geometric quintuple}
を与える $\widetilde G$-不変開部分とする。
このとき局所有限型 $k$-スキーム上の fppf site において
\[
 \mathcal F_{\mathrm{geo}}:
 [\mathscr G^{\mathrm{nl}}/\widetilde G]
 \longrightarrow\HHstack,\qquad
 (V_0,V_1,V_2,V_3,W\hookrightarrow
       V_0\otimes\cdots\otimes V_3)
 \longmapsto
 (\operatorname{qgr}C,\pi e_0C)
 \tag{GIT-H2}\label{eq:qa-geometric-cubic-comparison-map}
\]
という\emph{実際の相対群oid射}が定義される。
ここで $D$ はquintupleのcubic $\mathbb Z$-代数、
右加群・向きの規約を論文Bに合わせて
$C_{ij}=D_{-j,-i}$ とする。
$\operatorname{qgr}C$ の相対heartは
module-object 基底変更・fppf descent を伴って構成される。

この射は $K\simeq\Gm^4$ の慣性を自明にする自然な
二同型を持ち、その剛体化を通じて
\[
 \overline{\mathcal F}_{\mathrm{geo}}:
 [\mathscr G^{\mathrm{nl}}/(\widetilde G/K)]
 \longrightarrow\HHstack
 \tag{GIT-H2-rig}\label{eq:qa-geometric-cubic-rigid-map}
\]
と因数分解する。

さらに、論文Bの相対Picard二等分空間
$T_{\mathrm{root}}\to\mathscr G^{\mathrm{nl}}$ で選んだ
二周期 $\beta$ と Serre 平方同型を使って得る
$\HHstack$-対象は、選択にかかわらず
式~\eqref{eq:qa-geometric-cubic-comparison-map}
の像の引戻しである。
特に滑らかな楕円曲線上の四つのPicard根は、
\emph{固定quintupleの四つの異なるpointed対象ではない}。

この定理は $\overline{\mathcal F}_{\mathrm{geo}}$ の
表現可能性・全忠実性・幾何ファイバーの有限性、
およびGIT compactification全体への延長を主張しない。
\end{theorem}
\begin{proof}
$\mathscr G^{\mathrm{nl}}$ 上の普遍quintupleは
相対的なcubic $\mathbb Z$-代数 $D$ を定める。
双対・添字反転により $C_{ij}=D_{-j,-i}$ とするのは
論文B \cite[\texttt{thm:geometric-cubic-relative-effectivity}]{QuadricModuli2026}
の右tails規約である。
同論文 \cite[\texttt{prop:general-cubic-canonical-tails}]{QuadricModuli2026}
により、二周期を選ぶ\emph{前}から
$\operatorname{qgr}C$ と $\pi e_0C$ は
全係数基底変更に整合する相対Abel族であり、
compact Morita理論・Serre双加群も基底上に構成される。

一方、同論文
\cite[\texttt{lem:picard-halves},
\texttt{thm:relative-compatible-root}]{QuadricModuli2026}
は普遍基底を被覆するétale局所で
実際の二周期 $\beta$ とそのSerre平方比較を構成し、
\cite[\texttt{thm:geometric-cubic-relative-effectivity}]{QuadricModuli2026}
が同じ相対圏と構造対象に対して
AZ性・全軌道AS4被覆を含む $\HHstack$ の条件を証明する。
これらは\emph{存在条件}で、追加標識ではなく
fppf局所で判定される。
従って基底上にもともとある
$(\operatorname{qgr}C,\pi e_0C)$ が
$\HHstack$ の実際の対象を定める。
ここでは選んだ $\beta$ を降下させたのではなく、
既に存在する圏と構造対象に対して、
必要な局所witnessの存在を確認している。
同じ議論は任意の局所有限型の非被約基底変更で成立する。

quintupleの同型は $\mathbb Z$-代数の同型を誘導し、
右tailsの相対heartおよび $\pi e_0C$ の同型を保つ。
この操作は恒等、合成、任意の基底変更と両立する。
従って $\widetilde G$-equivariantな族として
式~\eqref{eq:qa-geometric-cubic-comparison-map} が
fppf descentにより構成される。

命題~\ref{prop:qa-ou-scalar-gerbe} の頂点ごとの
共役による自然同型は、添字零でスカラーを1に規格化すれば
pointedである。各 $\gamma_i$ は四つのスカラーの
積による指標として選べるので、
$K$ の積に関する二同型の整合性も成立する。
従って $K$ の作用を剛体化した表示スタックへ
この射は降下する。
最後に根空間上で$\beta$を替えても
$\operatorname{qgr}C$と$\pi e_0C$は元の基底から
引き戻した\emph{同じ}対象であり、結論が従う。
\end{proof}

\begin{corollary}[quintuple の剛体化後の相対次元]
\label{cor:qa-quintuple-rigid-dimension}
前定理の剛体化された表示スタックは
$k$ 上smoothで、純次元三である。
その $\HHstack$ への射~\eqref{eq:qa-geometric-cubic-rigid-map}
は同じ局所次元三のsmooth Artinスタック間の射である。
この次元の一致は相対射の étale 性や
pointed同値群の全忠実性を単独で意味しない。
\end{corollary}
\begin{proof}
$\mathscr G^{\mathrm{nl}}$ は16次元テンソル空間の
smoothな開部分である。
$\dim\widetilde G=4\cdot\dim\GL_2+\dim\GL_1=17$、
$\dim K=4$ なので
$\dim(\widetilde G/K)=13$ である。
従ってそのquotient stackの次元は $16-13=3$、
smooth性は群と表示開のsmooth性から従う。
一方、論文B
\cite[\texttt{thm:global-h2-smooth-finite-type}]{QuadricModuli2026}
が $\HHstack$ のsmooth純次元三を与える。
両者の間の相対射の微分・慣性核の計算は
この数値だけからは従わない。
\end{proof}

\begin{remark}[Cで次に必要な具体的計算]
\label{rem:qa-git-next-fiber}
実際の相対射は定理~\ref{thm:qa-geometric-cubic-relative-forgetting} で構成した。
次の段階では、式~\eqref{eq:qa-ou-fourfold-diagonal} の
$S_4$-monodromyと式~\eqref{eq:qa-picard-root-torsor} の
相対$E[2]$-torsorを同じsmooth elliptic quintuple基底へ引き戻し、
二つの四重射のモノドロミーとtorsor classを比較する。
そのうえで一般geometric cubic quintupleの
適合する有限型開において
$\operatorname{Aut}(\mathcal C,\OO)$ と
標識を保存する自己同型群の像を計算する。
次に上の二式を用いて、rulingと四窓を忘れる
pointed相対ファイバーを求める。
その後に、potentialからquintupleへの既存次数四射と照合する。
楕円型の構造対象の格子・特殊な自己同型、
noncubic境界への特殊化は別の段階とする。
$\HHstack$ 全体のgood moduliの存在、
標識適合領域の開性・全表示との相対同値は
現時点では未証明として保持する。
\end{remark}


\section{比較を完成させるための認識問題}
\label{sec:qa-remaining}

本稿の有限代数の開部分と、幾何的 avatar の開部分は、それぞれ独立に定義される。
両者を相対 AS 分類として同定するには、avatar から得る有限代数が
全周期にわたる相対 AS 代数を再構成し、逆向きの幾何構成と任意底変更に可換すること、
および同型の層と自己同型群が一致することを証明する必要がある。
体上の全単射はこの最後の同型層の比較を代替しない。

さらに $\HHstack$ へ進むには、選んだ構造対象に対する標準 heart、
相対 Serre 双対、豊富な exact 平方根、その同じ平方根の全軌道環の
正則正規 AS4 被覆を認識する必要がある。
これらを満たす部分には忘却射があるが、
その部分が十二本矢表示モジュライの開部分であることは別の命題である。
平方根や被覆を選んだ有限表示の incidence が存在することから、
それらの局所存在だけを要求する部分の開性や代数性を推論してはならない。

従って次の研究課題は、矢数だけを減らす追加の消去ではなく、
相対 AS 再構成と内在的認識の二つの比較を閉じることである。
一般 cubic quintuple の表示群のスカラー核を除いた後で、
pointed exact 自己同型との対応と相対ファイバーを決定することが先決となる。
既存 GIT の次数四を内在的な忘却次数と読み替えない。
どちらの場合も、marked な表示の自己同型から消える核、
構造対象の mutation による移動、非被約底での基底変更を同時に追跡する。
この順序により、有限型の表示空間を先に構成したことと、
内在的モジュライを表現したことの論理的な境界が明確になる。

\bibliographystyle{alpha}
\bibliography{references}
\end{document}
